% Genesis — Generated by generate_bible.py
% NET Bible text, scripture package formatting
\gdef\bbookchaptable{\vspace{4pt}%
\parbox[t]{\linewidth}{\centering\small\hyperlink{ch-genesis-1}{\textit{\small 1}}\hspace{0.4em}\hyperlink{ch-genesis-2}{\textit{\small 2}}\hspace{0.4em}\hyperlink{ch-genesis-3}{\textit{\small 3}}\hspace{0.4em}\hyperlink{ch-genesis-4}{\textit{\small 4}}\hspace{0.4em}\hyperlink{ch-genesis-5}{\textit{\small 5}}\hspace{0.4em}\hyperlink{ch-genesis-6}{\textit{\small 6}}\hspace{0.4em}\hyperlink{ch-genesis-7}{\textit{\small 7}}\hspace{0.4em}\hyperlink{ch-genesis-8}{\textit{\small 8}}\hspace{0.4em}\hyperlink{ch-genesis-9}{\textit{\small 9}}\hspace{0.4em}\hyperlink{ch-genesis-10}{\textit{\small 10}}\hspace{0.4em}\hyperlink{ch-genesis-11}{\textit{\small 11}}\hspace{0.4em}\hyperlink{ch-genesis-12}{\textit{\small 12}}\hspace{0.4em}\hyperlink{ch-genesis-13}{\textit{\small 13}}\hspace{0.4em}\hyperlink{ch-genesis-14}{\textit{\small 14}}\hspace{0.4em}\hyperlink{ch-genesis-15}{\textit{\small 15}}\hspace{0.4em}\hyperlink{ch-genesis-16}{\textit{\small 16}}\hspace{0.4em}\hyperlink{ch-genesis-17}{\textit{\small 17}}\hspace{0.4em}\hyperlink{ch-genesis-18}{\textit{\small 18}}\hspace{0.4em}\hyperlink{ch-genesis-19}{\textit{\small 19}}\hspace{0.4em}\hyperlink{ch-genesis-20}{\textit{\small 20}}\hspace{0.4em}\hyperlink{ch-genesis-21}{\textit{\small 21}}\hspace{0.4em}\hyperlink{ch-genesis-22}{\textit{\small 22}}\hspace{0.4em}\hyperlink{ch-genesis-23}{\textit{\small 23}}\hspace{0.4em}\hyperlink{ch-genesis-24}{\textit{\small 24}}\hspace{0.4em}\hyperlink{ch-genesis-25}{\textit{\small 25}}\hspace{0.4em}\hyperlink{ch-genesis-26}{\textit{\small 26}}\hspace{0.4em}\hyperlink{ch-genesis-27}{\textit{\small 27}}\hspace{0.4em}\hyperlink{ch-genesis-28}{\textit{\small 28}}\hspace{0.4em}\hyperlink{ch-genesis-29}{\textit{\small 29}}\hspace{0.4em}\hyperlink{ch-genesis-30}{\textit{\small 30}}\hspace{0.4em}\hyperlink{ch-genesis-31}{\textit{\small 31}}\hspace{0.4em}\hyperlink{ch-genesis-32}{\textit{\small 32}}\hspace{0.4em}\hyperlink{ch-genesis-33}{\textit{\small 33}}\hspace{0.4em}\hyperlink{ch-genesis-34}{\textit{\small 34}}\hspace{0.4em}\hyperlink{ch-genesis-35}{\textit{\small 35}}\hspace{0.4em}\hyperlink{ch-genesis-36}{\textit{\small 36}}\hspace{0.4em}\hyperlink{ch-genesis-37}{\textit{\small 37}}\hspace{0.4em}\hyperlink{ch-genesis-38}{\textit{\small 38}}\hspace{0.4em}\hyperlink{ch-genesis-39}{\textit{\small 39}}\hspace{0.4em}\hyperlink{ch-genesis-40}{\textit{\small 40}}\hspace{0.4em}\hyperlink{ch-genesis-41}{\textit{\small 41}}\hspace{0.4em}\hyperlink{ch-genesis-42}{\textit{\small 42}}\hspace{0.4em}\hyperlink{ch-genesis-43}{\textit{\small 43}}\hspace{0.4em}\hyperlink{ch-genesis-44}{\textit{\small 44}}\hspace{0.4em}\hyperlink{ch-genesis-45}{\textit{\small 45}}\hspace{0.4em}\hyperlink{ch-genesis-46}{\textit{\small 46}}\hspace{0.4em}\hyperlink{ch-genesis-47}{\textit{\small 47}}\hspace{0.4em}\hyperlink{ch-genesis-48}{\textit{\small 48}}\hspace{0.4em}\hyperlink{ch-genesis-49}{\textit{\small 49}}\hspace{0.4em}\hyperlink{ch-genesis-50}{\textit{\small 50}}}%
\vspace{6pt}}
\bbook[images/genese_heading]{The First Book of Moses}{Called Genesis}{Moses in effect declares three things, which are in this book chiefly to be considered: First, that the world and all things in it were created by God, and to praise his Name for the infinite graces, with which he had endued him, fell willingly from God through disobedience, who yet for his own mercies sake restored him to life, and confirmed him in the same by his promise of Christ to come, by whom he should overcome Satan, death and hell. Secondly, that the wicked, unmindful of God's most excellent benefits, remained still in their wickedness, and so falling most horribly from sin to sin, provoked God (who by his preachers called them continually to repentance) at length to destroy the whole world. Thirdly, he assures us by the examples of Abraham, Isaac, Jacob and the rest of the patriarchs, that his mercies never fail those whom he chooses to be his Church, and to profess his Name in earth, but in all their afflictions and persecutions he assists them, sends comfort, and delivers them, so that the beginning, increase, preservation and success of it might be attributed to God only. Moses shows by the examples of Cain, Ishmael, Esau and others, who were noble in man's judgment, that this Church depends not on the estimation and nobility of the world: and also by the fewness of those, who have at all times worshipped him purely according to his word that it stands not in the multitude, but in the poor and despised, in the small flock and little number, that man in his wisdom might be confounded, and the name of God praised forever.}{genesis}

\begin{scripture}

\markboth{Genesis 1:1}{Genesis 1:1}
\ch{1} \allowchapbreak\hypertarget{ch-genesis-1}{}\lettrine[lines=8]{\color{pentateuch}I}{\textsc{n}} the beginning God created the heavens and the earth.\gva{a}\marginnote{\gva{a}\,First of all, and before any creature was, God made heaven and earth out of nothing.}\everypar{}
\\\indent\markboth{Genesis 1:2}{Genesis 1:2}\vs{2} Now the earth was without shape and empty, and darkness was over the surface of the watery deep, but the Spirit of God was moving over the surface of the water.\gva{b}\marginnote{\gva{b}\,As an unformed lump and without any creature in it: for the waters covered everything.}\gva{c}\marginnote{\gva{c}\,Darkness covered the deep waters, for the waters covered everything.}\gva{d}\marginnote{\gva{d}\,He maintained this disordered mass by his secret power.}
\markboth{Genesis 1:3}{Genesis 1:3}\vs{3} God said, ``Let there be light.'' And there was light!\gva{e}\marginnote{\gva{e}\,The light was made before either Sun or Moon was created: therefore we must not attribute that to the creatures that are God's instruments, which only belong to God.}
\markboth{Genesis 1:4}{Genesis 1:4}\vs{4} God saw that the light was good, so God separated the light from the darkness.
\markboth{Genesis 1:5}{Genesis 1:5}\vs{5} God called the light ``day'' and the darkness ``night.'' There was evening, and there was morning, marking the first day.
\\\indent\markboth{Genesis 1:6}{Genesis 1:6}\vs{6} God said, ``Let there be an expanse in the midst of the waters and let it separate water from water.''
\markboth{Genesis 1:7}{Genesis 1:7}\vs{7} So God made the expanse and separated the water under the expanse from the water above it. It was so.\gva{f}\marginnote{\gva{f}\,As the sea and rivers, from those waters that are in the clouds, which are upheld by God's power, least they should overwhelm the world.}
\markboth{Genesis 1:8}{Genesis 1:8}\vs{8} God called the expanse ``sky.'' There was evening, and there was morning, a second day.\gva{g}\marginnote{\gva{g}\,That is, the region of the air, and all that is above us.}
\everypar{}

\parshape=0
\markboth{Genesis 1:9}{Genesis 1:9}\vs{9} God said, ``Let the water under the sky be gathered to one place and let dry ground appear.'' It was so.
\markboth{Genesis 1:10}{Genesis 1:10}\vs{10} God called the dry ground ``land'' and the gathered waters he called ``seas.'' God saw that it was good.
\everypar{}

\parshape=0
\markboth{Genesis 1:11}{Genesis 1:11}\vs{11} God said, ``Let the land produce vegetation: plants yielding seeds and trees on the land bearing fruit with seed in it, according to their kinds.'' It was so.\gva{h}\marginnote{\gva{h}\,So that we see it is the only the power of God's word that makes the earth fruitful, which naturally is barren.}
\markboth{Genesis 1:12}{Genesis 1:12}\vs{12} The land produced vegetation---plants yielding seeds according to their kinds, and trees bearing fruit with seed in it according to their kinds. God saw that it was good.\gva{i}\marginnote{\gva{i}\,This sentence is often repeated, to signify that God made all his creatures to serve for his glory and for the profit of man: but because of sin they were cursed, yet the elect, by Christ are restored, and serve to their wealth.}
\markboth{Genesis 1:13}{Genesis 1:13}\vs{13} There was evening, and there was morning, a third day.
\everypar{}

\parshape=0
\markboth{Genesis 1:14}{Genesis 1:14}\vs{14} God said, ``Let there be lights in the expanse of the sky to separate the day from the night, and let them be signs to indicate seasons and days and years,\gva{k}\marginnote{\gva{k}\,By the lights be means the sun, the moon, and the stars.}\gva{l}\marginnote{\gva{l}\,Which is the artificial day, from the sun rising, to the going down.}\gva{m}\marginnote{\gva{m}\,Of things belonging to natural and political orders and seasons.}
\markboth{Genesis 1:15}{Genesis 1:15}\vs{15} and let them serve as lights in the expanse of the sky to give light on the earth.'' It was so.
\markboth{Genesis 1:16}{Genesis 1:16}\vs{16} God made two great lights---the greater light to rule over the day and the lesser light to rule over the night. He made the stars also.\gva{n}\marginnote{\gva{n}\,That is, the sun and the moon, and here he speaks as man judges by his eye: for else the moon is less than the planet Saturn.}\gva{o}\marginnote{\gva{o}\,To give it sufficient light, as instruments appointed for the same, to serve man's purposes.}
\markboth{Genesis 1:17}{Genesis 1:17}\vs{17} God placed the lights in the expanse of the sky to shine on the earth,
\markboth{Genesis 1:18}{Genesis 1:18}\vs{18} to preside over the day and the night, and to separate the light from the darkness. God saw that it was good.
\markboth{Genesis 1:19}{Genesis 1:19}\vs{19} There was evening, and there was morning, a fourth day.
\everypar{}

\parshape=0
\markboth{Genesis 1:20}{Genesis 1:20}\vs{20} God said, ``Let the water swarm with swarms of living creatures and let birds fly above the earth across the expanse of the sky.''\gva{p}\marginnote{\gva{p}\,As fish and worms which slide, swim or creep.}
\markboth{Genesis 1:21}{Genesis 1:21}\vs{21} God created the great sea creatures and every living and moving thing with which the water swarmed, according to their kinds, and every winged bird according to its kind. God saw that it was good.\gva{q}\marginnote{\gva{q}\,The fish and fowls had both one beginning, in which we see that nature gives place to God's will, in that the one sort is made to fly about in the air, and the other to swim beneath in the water.}
\markboth{Genesis 1:22}{Genesis 1:22}\vs{22} God blessed them and said, ``Be fruitful and multiply and fill the water in the seas, and let the birds multiply on the earth.''\gva{r}\marginnote{\gva{r}\,That is, by the virtue of his word he gave power to his creatures to reproduce.}
\markboth{Genesis 1:23}{Genesis 1:23}\vs{23} There was evening, and there was morning, a fifth day.
\everypar{}

\parshape=0
\markboth{Genesis 1:24}{Genesis 1:24}\vs{24} God said, ``Let the land produce living creatures according to their kinds: cattle, creeping things, and wild animals, each according to its kind.'' It was so.
\markboth{Genesis 1:25}{Genesis 1:25}\vs{25} God made the wild animals according to their kinds, the cattle according to their kinds, and all the creatures that creep along the ground according to their kinds. God saw that it was good.
\everypar{}

\parshape=0
\markboth{Genesis 1:26}{Genesis 1:26}\vs{26} Then God said, ``Let us make humankind in our image, after our likeness, so they may rule over the fish of the sea and the birds of the air, over the cattle, and over all the earth, and over all the creatures that move on the earth.''\gva{s}\marginnote{\gva{s}\,God commanded the water and the earth to bring forth other creatures: but of man he says, "Let us make..." signifying that God takes counsel with his wisdom and virtue purposing to make an excellent work above all the rest of his creation.}\gva{t}\marginnote{\gva{t}\,This image and likeness of God in man is expounded in Ephesians 4:24 where it is written that man was created after God in righteousness and true holiness meaning by these two words, all perfection, as wisdom, truth, innocency, power, etc.}
\everypar{}

\parshape=0
\markboth{Genesis 1:27}{Genesis 1:27}\vs{27} God created humankind in his own image, in the image of God he created them, male and female he created them.
\everypar{}

\parshape=0
\markboth{Genesis 1:28}{Genesis 1:28}\vs{28} God blessed them and said to them, ``Be fruitful and multiply! Fill the earth and subdue it! Rule over the fish of the sea and the birds of the air and every creature that moves on the ground.''\gva{u}\marginnote{\gva{u}\,The propagation.}
\markboth{Genesis 1:29}{Genesis 1:29}\vs{29} Then God said, ``I now give you every seed-bearing plant on the face of the entire earth and every tree that has fruit with seed in it. They will be yours for food.\gva{x}\marginnote{\gva{x}\,God's great.}
\markboth{Genesis 1:30}{Genesis 1:30}\vs{30} And to all the animals of the earth, and to every bird of the air, and to all the creatures that move on the ground---everything that has living breath in it---I give every green plant for food.'' It was so.
\everypar{}

\parshape=0
\markboth{Genesis 1:31}{Genesis 1:31}\vs{31} God saw all that he had made---and it was very good! There was evening, and there was morning, the sixth day.

\markboth{Genesis 2:1}{Genesis 2:1}
\Needspace*{8\baselineskip}
\ch{2} \allowchapbreak\hypertarget{ch-genesis-2}{}\lettrine[lines=5]{\color{pentateuch}T}{\textsc{he}} heavens and the earth were completed with everything that was in them.\gva{a}\marginnote{\gva{a}\,That is, the innumerable abundance of creatures in heaven and earth.}\everypar{}
\markboth{Genesis 2:2}{Genesis 2:2}\vs{2} By the seventh day God finished the work that he had been doing, and he ceased on the seventh day all the work that he had been doing.\gva{b}\marginnote{\gva{b}\,For he had now finished his creation, but his providence still watches over his creatures and governs them.}
\markboth{Genesis 2:3}{Genesis 2:3}\vs{3} God blessed the seventh day and made it holy because on it he ceased all the work that he had been doing in creation.\gva{c}\marginnote{\gva{c}\,Appointed it to be kept holy, that man might in it consider the excellency of his works and God's goodness toward him.}
\everypar{}

\parshape=0
\markboth{Genesis 2:4}{Genesis 2:4}\vs{4} This is the account of the heavens and the earth when they were created---when the Lord God made the earth and heavens.
\everypar{}

\parshape=0
\markboth{Genesis 2:5}{Genesis 2:5}\vs{5} Now no shrub of the field had yet grown on the earth, and no plant of the field had yet sprouted, for the Lord God had not caused it to rain on the earth, and there was no man to cultivate the ground.\gva{d}\marginnote{\gva{d}\,God only opens the heavens and shuts them, he sends drought and rain according to his good pleasure.}
\markboth{Genesis 2:6}{Genesis 2:6}\vs{6} Springs would well up from the earth and water the whole surface of the ground.
\markboth{Genesis 2:7}{Genesis 2:7}\vs{7} The Lord God formed the man from the soil of the ground and breathed into his nostrils the breath of life, and the man became a living being.\gva{e}\marginnote{\gva{e}\,He shows what man's body was created from, to the intent that man should not glory in the excellency of his own nature.}
\everypar{}

\parshape=0
\markboth{Genesis 2:8}{Genesis 2:8}\vs{8} The Lord God planted an orchard in the east, in Eden; and there he placed the man he had formed.\gva{f}\marginnote{\gva{f}\,This was the name of a place, as some think in Mesopotamia, most pleasant and abundant in all things.}
\markboth{Genesis 2:9}{Genesis 2:9}\vs{9} The Lord God made all kinds of trees grow from the soil, every tree that was pleasing to look at and good for food. (Now the tree of life and the tree of the knowledge of good and evil were in the middle of the orchard.)\gva{g}\marginnote{\gva{g}\,Who was a sign of the life received from God.}\gva{h}\marginnote{\gva{h}\,That is, of miserable experience, which came by disobeying God.}
\everypar{}

\parshape=0
\markboth{Genesis 2:10}{Genesis 2:10}\vs{10} Now a river flows from Eden to water the orchard, and from there it divides into four headstreams.
\markboth{Genesis 2:11}{Genesis 2:11}\vs{11} The name of the first is Pishon; it runs through the entire land of Havilah, where there is gold.\gva{i}\marginnote{\gva{i}\,Havilah is a country adjoining Persia to the east, and inclining towards the west.}
\markboth{Genesis 2:12}{Genesis 2:12}\vs{12} (The gold of that land is pure; pearls and lapis lazuli are also there.)
\markboth{Genesis 2:13}{Genesis 2:13}\vs{13} The name of the second river is Gihon; it runs through the entire land of Cush.
\markboth{Genesis 2:14}{Genesis 2:14}\vs{14} The name of the third river is Tigris; it runs along the east side of Assyria. The fourth river is the Euphrates.
\everypar{}

\parshape=0
\markboth{Genesis 2:15}{Genesis 2:15}\vs{15} The Lord God took the man and placed him in the orchard in Eden to care for it and to maintain it.\gva{k}\marginnote{\gva{k}\,God would not have man idle, though as yet there was no need to labour.}
\markboth{Genesis 2:16}{Genesis 2:16}\vs{16} Then the Lord God commanded the man, ``You may freely eat fruit from every tree of the orchard,\gva{l}\marginnote{\gva{l}\,So that man might know there was a sovereign Lord, to whom he owed obedience.}
\markboth{Genesis 2:17}{Genesis 2:17}\vs{17} but you must not eat from the tree of the knowledge of good and evil, for when you eat from it you will surely die.''\gva{m}\marginnote{\gva{m}\,By death he means the separation of man from God, who is our life and chief happiness: and also that our disobedience is the cause of it.}
\everypar{}

\parshape=0
\markboth{Genesis 2:18}{Genesis 2:18}\vs{18} The Lord God said, ``It is not good for the man to be alone. I will make a companion for him who corresponds to him.''
\markboth{Genesis 2:19}{Genesis 2:19}\vs{19} The Lord God formed out of the ground every living animal of the field and every bird of the air. He brought them to the man to see what he would name them, and whatever the man called each living creature, that was its name.\gva{n}\marginnote{\gva{n}\,By moving them to come and submit themselves to Adam.}
\markboth{Genesis 2:20}{Genesis 2:20}\vs{20} So the man named all the animals, the birds of the air, and the living creatures of the field, but for Adam no companion who corresponded to him was found.
\markboth{Genesis 2:21}{Genesis 2:21}\vs{21} So the Lord God caused the man to fall into a deep sleep, and while he was asleep, he took part of the man's side and closed up the place with flesh.
\markboth{Genesis 2:22}{Genesis 2:22}\vs{22} Then the Lord God made a woman from the part he had taken out of the man, and he brought her to the man.\gva{o}\marginnote{\gva{o}\,Signifying that mankind was perfect, when the woman was created, who before was like an imperfect building.}
\everypar{}

\parshape=0
\markboth{Genesis 2:23}{Genesis 2:23}\vs{23} Then the man said, ``This one at last is bone of my bones and flesh of my flesh; this one will be called `woman,' for she was taken out of man.''
\everypar{}

\parshape=0
\markboth{Genesis 2:24}{Genesis 2:24}\vs{24} That is why a man leaves his father and mother and unites with his wife, and they become one family.\gva{p}\marginnote{\gva{p}\,So marriage requires a greater duty of us toward our wives, than otherwise we are bound to show to our parents.}
\markboth{Genesis 2:25}{Genesis 2:25}\vs{25} The man and his wife were both naked, but they were not ashamed.\gva{q}\marginnote{\gva{q}\,For before sin entered, all things were honest and comely.}

\markboth{Genesis 3:1}{Genesis 3:1}
\Needspace*{8\baselineskip}
\ch{3} \allowchapbreak\hypertarget{ch-genesis-3}{}\lettrine[lines=5]{\color{pentateuch}N}{\textsc{ow}} the serpent was shrewder than any of the wild animals that the Lord God had made. He said to the woman, ``Is it really true that God said, `You must not eat from any tree of the orchard'?''\gva{a}\marginnote{\gva{a}\,As Satan can change himself into an angel of light, so did he abuse the wisdom of the serpent to deceive man.}\gva{b}\marginnote{\gva{b}\,God allowed Satan to make the serpent his instrument and to speak through him.}\everypar{}
\markboth{Genesis 3:2}{Genesis 3:2}\vs{2} The woman said to the serpent, ``We may eat of the fruit from the trees of the orchard;
\markboth{Genesis 3:3}{Genesis 3:3}\vs{3} but concerning the fruit of the tree that is in the middle of the orchard God said, `You must not eat from it, and you must not touch it, or else you will die.'''\gva{c}\marginnote{\gva{c}\,In doubting God's warnings she yielded to Satan.}
\markboth{Genesis 3:4}{Genesis 3:4}\vs{4} The serpent said to the woman, ``Surely you will not die,\gva{d}\marginnote{\gva{d}\,This is Satan's chiefest subtilty, to cause us not to fear God's warnings.}
\markboth{Genesis 3:5}{Genesis 3:5}\vs{5} for God knows that when you eat from it your eyes will open and you will be like God, knowing good and evil.''\gva{e}\marginnote{\gva{e}\,As though he said, God forbids you to eat of the fruit, only because he knows that if you eat of it, you will be like him.}
\everypar{}

\parshape=0
\markboth{Genesis 3:6}{Genesis 3:6}\vs{6} When the woman saw that the tree produced fruit that was good for food, was attractive to the eye, and was desirable for making one wise, she took some of its fruit and ate it. She also gave some of it to her husband who was with her, and he ate it.\gva{f}\marginnote{\gva{f}\,Not so much to please his wife, as moved by ambition at her persuasion.}
\markboth{Genesis 3:7}{Genesis 3:7}\vs{7} Then the eyes of both of them opened, and they knew they were naked; so they sewed fig leaves together and made coverings for themselves.\gva{g}\marginnote{\gva{g}\,They began to feel their misery, but they did not seek God for a remedy.}
\everypar{}

\parshape=0
\markboth{Genesis 3:8}{Genesis 3:8}\vs{8} Then the man and his wife heard the sound of the Lord God moving about in the orchard at the breezy time of the day, and they hid from the Lord God among the trees of the orchard.\gva{h}\marginnote{\gva{h}\,The sinful conscience flees God's presence.}
\markboth{Genesis 3:9}{Genesis 3:9}\vs{9} But the Lord God called to the man and said to him, ``Where are you?''
\markboth{Genesis 3:10}{Genesis 3:10}\vs{10} The man replied, ``I heard you moving about in the orchard, and I was afraid because I was naked, so I hid.''\gva{i}\marginnote{\gva{i}\,His hypocrisy appears in that he hid the cause of his nakedness, which was the transgression of God's commandment.}
\markboth{Genesis 3:11}{Genesis 3:11}\vs{11} And the Lord God said, ``Who told you that you were naked? Did you eat from the tree that I commanded you not to eat from?''
\markboth{Genesis 3:12}{Genesis 3:12}\vs{12} The man said, ``The woman whom you gave me, she gave me some fruit from the tree and I ate it.''\gva{k}\marginnote{\gva{k}\,His wickedness and lack of true repentance appears in this that he blamed God because he had given him a wife.}
\markboth{Genesis 3:13}{Genesis 3:13}\vs{13} So the Lord God said to the woman, ``What is this you have done?'' And the woman replied, ``The serpent tricked me, and I ate.''\gva{l}\marginnote{\gva{l}\,Instead of confessing her sin, she increases it by accusing the serpent.}
\everypar{}

\parshape=0
\markboth{Genesis 3:14}{Genesis 3:14}\vs{14} The Lord God said to the serpent, ``Because you have done this, cursed are you above all the cattle and all the living creatures of the field! On your belly you will crawl and dust you will eat all the days of your life.\gva{m}\marginnote{\gva{m}\,He asked the reason from Adam and his wife, because he would bring them to repentance, but he does not ask the serpent, because he would show him no mercy.}\gva{n}\marginnote{\gva{n}\,As a vile and contemptible beast, Isaiah 65:25 .}
\everypar{}

\parshape=0
\markboth{Genesis 3:15}{Genesis 3:15}\vs{15} And I will put hostility between you and the woman and between your offspring and her offspring; he will strike your head, and you will strike his heel.''\gva{o}\marginnote{\gva{o}\,He chiefly means Satan, by whose action and deceit the serpent deceived the woman.}\gva{p}\marginnote{\gva{p}\,That is, the power of sin and death.}\gva{q}\marginnote{\gva{q}\,Satan shall sting Christ and his members, but not overcome them.}
\everypar{}

\parshape=0
\markboth{Genesis 3:16}{Genesis 3:16}\vs{16} To the woman he said, ``I will greatly increase your labor pains; with pain you will give birth to children. You will want to control your husband, but he will dominate you.''\gva{r}\marginnote{\gva{r}\,The Lord comforts Adam by the promise of the blessed seed, and also punishes the body for the sin which the soul should have been punished for; that the spirit having conceived hope of forgiveness might live by faith. 1 Corinthians 14:34 .}
\everypar{}

\parshape=0
\markboth{Genesis 3:17}{Genesis 3:17}\vs{17} But to Adam he said, ``Because you obeyed your wife and ate from the tree about which I commanded you, `You must not eat from it,' the ground is cursed because of you; in painful toil you will eat of it all the days of your life.\gva{s}\marginnote{\gva{s}\,The transgression of God's commandment was the reason that both mankind and all other creatures were subject to the curse.}
\everypar{}

\parshape=0
\markboth{Genesis 3:18}{Genesis 3:18}\vs{18} It will produce thorns and thistles for you, but you will eat the grain of the field.\gva{t}\marginnote{\gva{t}\,These are not the natural fruit of the earth, but proceed from the corruption of sin.}
\everypar{}

\parshape=0
\markboth{Genesis 3:19}{Genesis 3:19}\vs{19} By the sweat of your brow you will eat food until you return to the ground, for out of it you were taken; for you are dust, and to dust you will return.''
\everypar{}

\parshape=0
\markboth{Genesis 3:20}{Genesis 3:20}\vs{20} The man named his wife Eve, because she was the mother of all the living.
\markboth{Genesis 3:21}{Genesis 3:21}\vs{21} The Lord God made garments from skin for Adam and his wife, and clothed them.\gva{u}\marginnote{\gva{u}\,Or, gave them knowledge to make themselves coats.}
\markboth{Genesis 3:22}{Genesis 3:22}\vs{22} And the Lord God said, ``Now that the man has become like one of us, knowing good and evil, he must not be allowed to stretch out his hand and take also from the tree of life and eat, and live forever.''\gva{x}\marginnote{\gva{x}\,By this derision by reproaches Adam's misery, into which he was fallen by ambition.}\gva{y}\marginnote{\gva{y}\,Adam deprived of life, lost also the sign of it.}
\markboth{Genesis 3:23}{Genesis 3:23}\vs{23} So the Lord God expelled him from the orchard in Eden to cultivate the ground from which he had been taken.
\markboth{Genesis 3:24}{Genesis 3:24}\vs{24} When he drove the man out, he placed on the eastern side of the orchard in Eden angelic sentries who used the flame of a whirling sword to guard the way to the tree of life.

\markboth{Genesis 4:1}{Genesis 4:1}
\Needspace*{8\baselineskip}
\ch{4} \allowchapbreak\hypertarget{ch-genesis-4}{}\lettrine[lines=5]{\color{pentateuch}N}{\textsc{ow}} the man was intimate with his wife Eve, and she became pregnant and gave birth to Cain. Then she said, ``I have created a man just as the Lord did!''\gva{a}\marginnote{\gva{a}\,Man's nature, the estate of marriage, and God's blessing were not utterly abolished through sin, but the quality or condition of it was changed.}\gva{b}\marginnote{\gva{b}\,That is, according to the Lord's promise, as some read Genesis 3:15 , "To the Lord" rejoicing for the son she had born, whom she would offer to the Lord as the first fruits of her birth.}\everypar{}
\markboth{Genesis 4:2}{Genesis 4:2}\vs{2} Then she gave birth to his brother Abel. Abel took care of the flocks, while Cain cultivated the ground.
\\\indent\markboth{Genesis 4:3}{Genesis 4:3}\vs{3} At the designated time Cain brought some of the fruit of the ground for an offering to the Lord.\gva{c}\marginnote{\gva{c}\,This declares that the father instructed his children in the knowledge of God, and also how God gave them sacrifices to signify their salvation, though they were destitute of the ordinance of the tree of life.}
\markboth{Genesis 4:4}{Genesis 4:4}\vs{4} But Abel brought some of the firstborn of his flock---even the fattest of them. And the Lord was pleased with Abel and his offering,
\markboth{Genesis 4:5}{Genesis 4:5}\vs{5} but with Cain and his offering he was not pleased. So Cain became very angry, and his expression was downcast.
\everypar{}

\parshape=0
\markboth{Genesis 4:6}{Genesis 4:6}\vs{6} Then the Lord said to Cain, ``Why are you angry, and why is your expression downcast?
\markboth{Genesis 4:7}{Genesis 4:7}\vs{7} Is it not true that if you do what is right, you will be fine? But if you do not do what is right, sin is crouching at the door. It desires to dominate you, but you must subdue it.''\gva{e}\marginnote{\gva{e}\,Both you and your sacrifice shall be acceptable to me.}\gva{f}\marginnote{\gva{f}\,Sin will still torment your conscience.}\gva{g}\marginnote{\gva{g}\,The dignity of the first born is given to Cain over Abel.}
\everypar{}

\parshape=0
\markboth{Genesis 4:8}{Genesis 4:8}\vs{8} Cain said to his brother Abel, ``Let's go out to the field.'' While they were in the field, Cain attacked his brother Abel and killed him.
\everypar{}

\parshape=0
\markboth{Genesis 4:9}{Genesis 4:9}\vs{9} Then the Lord said to Cain, ``Where is your brother Abel?'' And he replied, ``I don't know! Am I my brother's guardian?''\gva{h}\marginnote{\gva{h}\,This is the nature of the reprobate when they are rebuke for their hypocrisy, even to neglect God and outrage him.}
\markboth{Genesis 4:10}{Genesis 4:10}\vs{10} But the Lord said, ``What have you done? The voice of your brother's blood is crying out to me from the ground!
\markboth{Genesis 4:11}{Genesis 4:11}\vs{11} So now you are banished from the ground, which has opened its mouth to receive your brother's blood from your hand.\gva{k}\marginnote{\gva{k}\,The earth will be a witness against you, which mercifully received the blood you most cruelly shed.}
\markboth{Genesis 4:12}{Genesis 4:12}\vs{12} When you try to cultivate the ground it will no longer yield its best for you. You will be a homeless wanderer on the earth.''\gva{l}\marginnote{\gva{l}\,You will never have rest for your heart will be in continual fear and worry.}
\everypar{}

\parshape=0
\markboth{Genesis 4:13}{Genesis 4:13}\vs{13} Then Cain said to the Lord, ``My punishment is too great to endure!\gva{m}\marginnote{\gva{m}\,He burdens God as a cruel judge because he punished him so severely.}
\markboth{Genesis 4:14}{Genesis 4:14}\vs{14} Look, you are driving me off the land today, and I must hide from your presence. I will be a homeless wanderer on the earth; whoever finds me will kill me!''
\markboth{Genesis 4:15}{Genesis 4:15}\vs{15} But the Lord said to him, ``All right then, if anyone kills Cain, Cain will be avenged seven times as much.'' Then the Lord put a special mark on Cain so that no one who found him would strike him down.\gva{n}\marginnote{\gva{n}\,Not for the love he had for Cain, but to suppress murder.}\gva{o}\marginnote{\gva{o}\,Which was some visible sign of God's judgment, that others should fear by it.}
\markboth{Genesis 4:16}{Genesis 4:16}\vs{16} So Cain went out from the presence of the Lord and lived in the land of Nod, east of Eden.
\everypar{}

\parshape=0
\markboth{Genesis 4:17}{Genesis 4:17}\vs{17} Cain was intimate with his wife, and she became pregnant and gave birth to Enoch. Cain was building a city, and he named the city after his son Enoch.\gva{p}\marginnote{\gva{p}\,Thinking by this to be safe, and to have less reason to fear God's judgments against him.}
\markboth{Genesis 4:18}{Genesis 4:18}\vs{18} To Enoch was born Irad, and Irad was the father of Mehujael. Mehujael was the father of Methushael, and Methushael was the father of Lamech.
\everypar{}

\parshape=0
\markboth{Genesis 4:19}{Genesis 4:19}\vs{19} Lamech took two wives for himself; the name of the first was Adah, and the name of the second was Zillah.\gva{q}\marginnote{\gva{q}\,The lawful institution of marriage, which is, that two should be one flesh, was first corrupted in the house of Cain by Lamech.}
\markboth{Genesis 4:20}{Genesis 4:20}\vs{20} Adah gave birth to Jabal; he was the first of those who live in tents and keep livestock.
\markboth{Genesis 4:21}{Genesis 4:21}\vs{21} The name of his brother was Jubal; he was the first of all who play the harp and the flute.
\markboth{Genesis 4:22}{Genesis 4:22}\vs{22} Now Zillah also gave birth to Tubal-Cain, who heated metal and shaped all kinds of tools made of bronze and iron. The sister of Tubal-Cain was Naamah.
\everypar{}

\parshape=0
\markboth{Genesis 4:23}{Genesis 4:23}\vs{23} Lamech said to his wives, ``Adah and Zillah, listen to me! You wives of Lamech, hear my words! I have killed a man for wounding me, a young man for hurting me.\gva{r}\marginnote{\gva{r}\,His wives seeing that all men hated him for his cruelty, were afraid, therefore he brags that there is none strong enough to resist, even though he was already wounded.}
\everypar{}

\parshape=0
\markboth{Genesis 4:24}{Genesis 4:24}\vs{24} If Cain is to be avenged seven times as much, then Lamech seventy-seven times!''\gva{s}\marginnote{\gva{s}\,He mocked at God's tolerance in Cain jesting as though God would allow no one to punish him and yet give him permission to murder others.}
\everypar{}

\parshape=0
\markboth{Genesis 4:25}{Genesis 4:25}\vs{25} And Adam was intimate with his wife again, and she gave birth to a son. She named him Seth, saying, ``God has given me another child in place of Abel because Cain killed him.''
\markboth{Genesis 4:26}{Genesis 4:26}\vs{26} And a son was also born to Seth, whom he named Enosh. At that time people began to worship the Lord.\gva{t}\marginnote{\gva{t}\,In these days God began to move the hearts of the godly to restore religion, which had been suppressed by the wicked for a long time.}

\markboth{Genesis 5:1}{Genesis 5:1}
\Needspace*{8\baselineskip}
\ch{5} \allowchapbreak\hypertarget{ch-genesis-5}{}\lettrine[lines=5]{\color{pentateuch}T}{\textsc{his}} is the record of the family line of Adam. When God created humankind, he made them in the likeness of God.\gva{a}\marginnote{\gva{a}\,Read Genesis 1:26 .}\everypar{}
\markboth{Genesis 5:2}{Genesis 5:2}\vs{2} He created them male and female; when they were created, he blessed them and named them ``humankind.''\gva{b}\marginnote{\gva{b}\,By giving them both one name, he notes the inseparable conjunction of man and wife.}
\\\indent\markboth{Genesis 5:3}{Genesis 5:3}\vs{3} When Adam had lived 130 years he fathered a son in his own likeness, according to his image, and he named him Seth.\gva{c}\marginnote{\gva{c}\,As well, concerning his creation, as his corruption.}
\markboth{Genesis 5:4}{Genesis 5:4}\vs{4} The length of time Adam lived after he became the father of Seth was 800 years; during this time he had other sons and daughters.
\markboth{Genesis 5:5}{Genesis 5:5}\vs{5} The entire lifetime of Adam was 930 years, and then he died.
\everypar{}

\parshape=0
\markboth{Genesis 5:6}{Genesis 5:6}\vs{6} When Seth had lived 105 years, he became the father of Enosh.\gva{d}\marginnote{\gva{d}\,He proves Adam's generation by those who came from Seth, to show the true Church, and also what care God had over the same from the beginning, in that he continued his graces toward it by a continual succession.}
\markboth{Genesis 5:7}{Genesis 5:7}\vs{7} Seth lived 807 years after he became the father of Enosh, and he had other sons and daughters.
\markboth{Genesis 5:8}{Genesis 5:8}\vs{8} The entire lifetime of Seth was 912 years, and then he died.\gva{e}\marginnote{\gva{e}\,The main reason for long life in the first age, was the multiplication of mankind, that according to God's commandment at the beginning the world might be filled with people, who would universally praise him.}
\everypar{}

\parshape=0
\markboth{Genesis 5:9}{Genesis 5:9}\vs{9} When Enosh had lived 90 years, he became the father of Kenan.
\markboth{Genesis 5:10}{Genesis 5:10}\vs{10} Enosh lived 815 years after he became the father of Kenan, and he had other sons and daughters.
\markboth{Genesis 5:11}{Genesis 5:11}\vs{11} The entire lifetime of Enosh was 905 years, and then he died.
\everypar{}

\parshape=0
\markboth{Genesis 5:12}{Genesis 5:12}\vs{12} When Kenan had lived 70 years, he became the father of Mahalalel.
\markboth{Genesis 5:13}{Genesis 5:13}\vs{13} Kenan lived 840 years after he became the father of Mahalalel, and he had other sons and daughters.
\markboth{Genesis 5:14}{Genesis 5:14}\vs{14} The entire lifetime of Kenan was 910 years, and then he died.
\everypar{}

\parshape=0
\markboth{Genesis 5:15}{Genesis 5:15}\vs{15} When Mahalalel had lived 65 years, he became the father of Jared.
\markboth{Genesis 5:16}{Genesis 5:16}\vs{16} Mahalalel lived 830 years after he became the father of Jared, and he had other sons and daughters.
\markboth{Genesis 5:17}{Genesis 5:17}\vs{17} The entire lifetime of Mahalalel was 895 years, and then he died.
\everypar{}

\parshape=0
\markboth{Genesis 5:18}{Genesis 5:18}\vs{18} When Jared had lived 162 years, he became the father of Enoch.
\markboth{Genesis 5:19}{Genesis 5:19}\vs{19} Jared lived 800 years after he became the father of Enoch, and he had other sons and daughters.
\markboth{Genesis 5:20}{Genesis 5:20}\vs{20} The entire lifetime of Jared was 962 years, and then he died.
\everypar{}

\parshape=0
\markboth{Genesis 5:21}{Genesis 5:21}\vs{21} When Enoch had lived 65 years, he became the father of Methuselah.
\markboth{Genesis 5:22}{Genesis 5:22}\vs{22} After he became the father of Methuselah, Enoch walked with God for 300 years, and he had other sons and daughters.\gva{f}\marginnote{\gva{f}\,That is, he led an upright and godly life.}
\markboth{Genesis 5:23}{Genesis 5:23}\vs{23} The entire lifetime of Enoch was 365 years.
\markboth{Genesis 5:24}{Genesis 5:24}\vs{24} Enoch walked with God, and then he disappeared because God took him away.\gva{g}\marginnote{\gva{g}\,To show that there was a better life prepared and to be a testimony of the immortality of souls and bodies. To inquire where he went is mere curiosity.}
\everypar{}

\parshape=0
\markboth{Genesis 5:25}{Genesis 5:25}\vs{25} When Methuselah had lived 187 years, he became the father of Lamech.
\markboth{Genesis 5:26}{Genesis 5:26}\vs{26} Methuselah lived 782 years after he became the father of Lamech, and he had other sons and daughters.
\markboth{Genesis 5:27}{Genesis 5:27}\vs{27} The entire lifetime of Methuselah was 969 years, and then he died.
\everypar{}

\parshape=0
\markboth{Genesis 5:28}{Genesis 5:28}\vs{28} When Lamech had lived 182 years, he had a son.
\markboth{Genesis 5:29}{Genesis 5:29}\vs{29} He named him Noah, saying, ``This one will bring us comfort from our labor and from the painful toil of our hands because of the ground that the Lord has cursed.''\gva{h}\marginnote{\gva{h}\,Lamech had respect for the promise, Genesis 3:15 , and desired to see the deliverer who would be sent and yet saw but a figure of it. He spoke this by the spirit of prophecy because Noah delivered the Church and preserved it by his obedience.}
\markboth{Genesis 5:30}{Genesis 5:30}\vs{30} Lamech lived 595 years after he became the father of Noah, and he had other sons and daughters.
\markboth{Genesis 5:31}{Genesis 5:31}\vs{31} The entire lifetime of Lamech was 777 years, and then he died.
\everypar{}

\parshape=0
\markboth{Genesis 5:32}{Genesis 5:32}\vs{32} After Noah was 500 years old, he became the father of Shem, Ham, and Japheth.

\markboth{Genesis 6:1}{Genesis 6:1}
\Needspace*{8\baselineskip}
\ch{6} \allowchapbreak\hypertarget{ch-genesis-6}{}\lettrine[lines=5]{\color{pentateuch}W}{\textsc{hen}} humankind began to multiply on the face of the earth, and daughters were born to them,\everypar{}
\markboth{Genesis 6:2}{Genesis 6:2}\vs{2} the sons of God saw that the daughters of humankind were beautiful. Thus they took wives for themselves from any they chose.\gva{a}\marginnote{\gva{a}\,The children of the godly who began to degenerate.}\gva{b}\marginnote{\gva{b}\,Those that had wicked parents, as if from Cain.}\gva{c}\marginnote{\gva{c}\,Having more respect for their beauty and worldly considerations than for their manners and godliness.}
\markboth{Genesis 6:3}{Genesis 6:3}\vs{3} So the Lord said, ``My Spirit will not remain in humankind indefinitely, since they are mortal. They will remain for 120 more years.''\gva{d}\marginnote{\gva{d}\,Because man could not by won by God's leniency and patience by which he tried to win him, he would no longer withhold his vengeance.}\gva{e}\marginnote{\gva{e}\,Which time span God gave man to repent before he would destroy the earth, 1 Peter 3:20 .}
\everypar{}

\parshape=0
\markboth{Genesis 6:4}{Genesis 6:4}\vs{4} The Nephilim were on the earth in those days (and also after this) when the sons of God would sleep with the daughters of humankind, who gave birth to their children. They were the mighty heroes of old, the famous men.\gva{f}\marginnote{\gva{f}\,Who usurped authority over others, and degenerated from that simplicity, in which their father's lived.}
\everypar{}

\parshape=0
\markboth{Genesis 6:5}{Genesis 6:5}\vs{5} But the Lord saw that the wickedness of humankind had become great on the earth. Every inclination of the thoughts of their minds was only evil all the time.
\markboth{Genesis 6:6}{Genesis 6:6}\vs{6} The Lord regretted that he had made humankind on the earth, and he was highly offended.\gva{g}\marginnote{\gva{g}\,God never repents, but he speaks in human terms, because he destroyed him, and in a way denied him as his creature.}
\markboth{Genesis 6:7}{Genesis 6:7}\vs{7} So the Lord said, ``I will wipe humankind, whom I have created, from the face of the earth---everything from humankind to animals, including creatures that move on the ground and birds of the air, for I regret that I have made them.''\gva{h}\marginnote{\gva{h}\,God declares how much he detests sin, seeing the punishment of it extends to the brute beasts.}
\everypar{}

\parshape=0
\markboth{Genesis 6:8}{Genesis 6:8}\vs{8} But Noah found favor in the sight of the Lord.\gva{i}\marginnote{\gva{i}\,God was merciful to him.}
\everypar{}

\parshape=0
\markboth{Genesis 6:9}{Genesis 6:9}\vs{9} This is the account of Noah. Noah was a godly man; he was blameless among his contemporaries. He walked with God.
\markboth{Genesis 6:10}{Genesis 6:10}\vs{10} Noah had three sons: Shem, Ham, and Japheth.
\everypar{}

\parshape=0
\markboth{Genesis 6:11}{Genesis 6:11}\vs{11} The earth was ruined in the sight of God; the earth was filled with violence.\gva{k}\marginnote{\gva{k}\,Meaning, that all were given to the contempt of God, and oppression of their neighbours.}
\markboth{Genesis 6:12}{Genesis 6:12}\vs{12} God saw the earth, and indeed it was ruined, for all living creatures on the earth were sinful.
\markboth{Genesis 6:13}{Genesis 6:13}\vs{13} So God said to Noah, ``I have decided that all living creatures must die, for the earth is filled with violence because of them. Now I am about to destroy them and the earth.
\markboth{Genesis 6:14}{Genesis 6:14}\vs{14} Make for yourself an ark of cypress wood. Make rooms in the ark, and cover it with pitch inside and out.
\markboth{Genesis 6:15}{Genesis 6:15}\vs{15} This is how you should make it: The ark is to be 450 feet long, 75 feet wide, and 45 feet high.
\markboth{Genesis 6:16}{Genesis 6:16}\vs{16} Make a roof for the ark and finish it, leaving 18 inches from the top. Put a door in the side of the ark, and make lower, middle, and upper decks.
\markboth{Genesis 6:17}{Genesis 6:17}\vs{17} I am about to bring floodwaters on the earth to destroy from under the sky all the living creatures that have the breath of life in them. Everything that is on the earth will die,
\markboth{Genesis 6:18}{Genesis 6:18}\vs{18} but I will confirm my covenant with you. You will enter the ark---you, your sons, your wife, and your sons' wives with you.\gva{m}\marginnote{\gva{m}\,So that in this great undertaking and mocking of the whole world, you may be confirmed so your faith does not fail.}
\markboth{Genesis 6:19}{Genesis 6:19}\vs{19} You must bring into the ark two of every kind of living creature from all flesh, male and female, to keep them alive with you.
\markboth{Genesis 6:20}{Genesis 6:20}\vs{20} Of the birds after their kinds, and of the cattle after their kinds, and of every creeping thing of the ground after its kind, two of every kind will come to you so you can keep them alive.
\markboth{Genesis 6:21}{Genesis 6:21}\vs{21} And you must take for yourself every kind of food that is eaten, and gather it together. It will be food for you and for them.''
\everypar{}

\parshape=0
\markboth{Genesis 6:22}{Genesis 6:22}\vs{22} And Noah did all that God commanded him---he did indeed.\gva{n}\marginnote{\gva{n}\,That is, he obeyed God's commandment in all points without adding or taking away.}

\markboth{Genesis 7:1}{Genesis 7:1}
\Needspace*{8\baselineskip}
\ch{7} \allowchapbreak\hypertarget{ch-genesis-7}{}\lettrine[lines=5]{\color{pentateuch}T}{\textsc{he}} Lord said to Noah, ``Come into the ark, you and all your household, for I consider you godly among this generation.\gva{a}\marginnote{\gva{a}\,In respect to the rest of the world, and because he had a desire to serve God and live uprightly.}\everypar{}
\markboth{Genesis 7:2}{Genesis 7:2}\vs{2} You must take with you seven pairs of every kind of clean animal, the male and its mate, two of every kind of unclean animal, the male and its mate,\gva{b}\marginnote{\gva{b}\,Which might be offered in sacrifice, of which six were for breeding and the seventh for sacrifice.}
\markboth{Genesis 7:3}{Genesis 7:3}\vs{3} and also seven pairs of every kind of bird in the sky, male and female, to preserve their offspring on the face of the entire earth.
\markboth{Genesis 7:4}{Genesis 7:4}\vs{4} For in seven days I will cause it to rain on the earth for 40 days and 40 nights, and I will wipe from the face of the ground every living thing that I have made.''
\everypar{}

\parshape=0
\markboth{Genesis 7:5}{Genesis 7:5}\vs{5} And Noah did all that the Lord commanded him.
\everypar{}

\parshape=0
\markboth{Genesis 7:6}{Genesis 7:6}\vs{6} Noah was 600 years old when the floodwaters engulfed the earth.
\markboth{Genesis 7:7}{Genesis 7:7}\vs{7} Noah entered the ark along with his sons, his wife, and his sons' wives because of the floodwaters.
\markboth{Genesis 7:8}{Genesis 7:8}\vs{8} Pairs of clean animals, of unclean animals, of birds, and of everything that creeps along the ground,
\markboth{Genesis 7:9}{Genesis 7:9}\vs{9} male and female, came into the ark to Noah, just as God had commanded him.\gva{c}\marginnote{\gva{c}\,God compelled them to present themselves to Noah, as they did before to Adam, when he gave them names, Genesis 2:19 .}
\markboth{Genesis 7:10}{Genesis 7:10}\vs{10} And after seven days the floodwaters engulfed the earth.
\everypar{}

\parshape=0
\markboth{Genesis 7:11}{Genesis 7:11}\vs{11} In the six hundredth year of Noah's life, in the second month, on the seventeenth day of the month---on that day all the fountains of the great deep burst open and the floodgates of the heavens were opened.\gva{e}\marginnote{\gva{e}\,Both the waters in the earth overflowed and also the clouds poured down.}
\markboth{Genesis 7:12}{Genesis 7:12}\vs{12} And the rain fell on the earth 40 days and 40 nights.
\everypar{}

\parshape=0
\markboth{Genesis 7:13}{Genesis 7:13}\vs{13} On that very day Noah entered the ark, accompanied by his sons Shem, Ham, and Japheth, along with his wife and his sons' three wives.
\markboth{Genesis 7:14}{Genesis 7:14}\vs{14} They entered, along with every living creature after its kind, every animal after its kind, every creeping thing that creeps on the earth after its kind, and every bird after its kind, everything with wings.
\markboth{Genesis 7:15}{Genesis 7:15}\vs{15} Pairs of all creatures that have the breath of life came into the ark to Noah.\gva{f}\marginnote{\gva{f}\,Every living thing that God would have be preserved on earth, came into the ark to Noah.}
\markboth{Genesis 7:16}{Genesis 7:16}\vs{16} Those that entered were male and female, just as God commanded him. Then the Lord shut him in.\gva{g}\marginnote{\gva{g}\,So that God's secret power defended him against the rage of the mighty waters.}
\everypar{}

\parshape=0
\markboth{Genesis 7:17}{Genesis 7:17}\vs{17} The flood engulfed the earth for 40 days. As the waters increased, they lifted the ark and raised it above the earth.
\markboth{Genesis 7:18}{Genesis 7:18}\vs{18} The waters completely overwhelmed the earth, and the ark floated on the surface of the waters.
\markboth{Genesis 7:19}{Genesis 7:19}\vs{19} The waters completely inundated the earth so that even all the high mountains under the entire sky were covered.
\markboth{Genesis 7:20}{Genesis 7:20}\vs{20} The waters rose more than 20 feet above the mountains.
\markboth{Genesis 7:21}{Genesis 7:21}\vs{21} And all living things that moved on the earth died, including the birds, domestic animals, wild animals, all the creatures that swarm over the earth, and all humankind.
\markboth{Genesis 7:22}{Genesis 7:22}\vs{22} Everything on dry land that had the breath of life in its nostrils died.
\markboth{Genesis 7:23}{Genesis 7:23}\vs{23} So the Lord destroyed every living thing that was on the surface of the ground, including people, animals, creatures that creep along the ground, and birds of the sky. They were wiped off the earth. Only Noah and those who were with him in the ark survived.\gva{h}\marginnote{\gva{h}\,That is, by God.}\gva{i}\marginnote{\gva{i}\,Learn what it is to obey God only, and to forsake the multitude, 1 Peter 3:20 .}
\markboth{Genesis 7:24}{Genesis 7:24}\vs{24} The waters prevailed over the earth for 150 days.

\markboth{Genesis 8:1}{Genesis 8:1}
\Needspace*{8\baselineskip}
\ch{8} \allowchapbreak\hypertarget{ch-genesis-8}{}\lettrine[lines=5]{\color{pentateuch}B}{\textsc{ut}} God remembered Noah and all the wild animals and domestic animals that were with him in the ark. God caused a wind to blow over the earth and the waters receded.\gva{a}\marginnote{\gva{a}\,Not that God forgets his at any time, but when he sends comfort then he shows that he remembers them.}\gva{b}\marginnote{\gva{b}\,If God remembered every brute beast, that ought also to assure his children.}\everypar{}
\markboth{Genesis 8:2}{Genesis 8:2}\vs{2} The fountains of the deep and the floodgates of heaven were closed, and the rain stopped falling from the sky.
\markboth{Genesis 8:3}{Genesis 8:3}\vs{3} The waters kept receding steadily from the earth, so that they had gone down by the end of the 150 days.
\markboth{Genesis 8:4}{Genesis 8:4}\vs{4} On the seventeenth day of the seventh month, the ark came to rest on one of the mountains of Ararat.\gva{c}\marginnote{\gva{c}\,Part of September and part of October.}
\markboth{Genesis 8:5}{Genesis 8:5}\vs{5} The waters kept on receding until the tenth month. On the first day of the tenth month, the tops of the mountains became visible.\gva{d}\marginnote{\gva{d}\,Which was the month of December.}
\everypar{}

\parshape=0
\markboth{Genesis 8:6}{Genesis 8:6}\vs{6} At the end of 40 days, Noah opened the window he had made in the ark
\markboth{Genesis 8:7}{Genesis 8:7}\vs{7} and sent out a raven; it kept flying back and forth until the waters had dried up on the earth.
\everypar{}

\parshape=0
\markboth{Genesis 8:8}{Genesis 8:8}\vs{8} Then Noah sent out a dove to see if the waters had receded from the surface of the ground.
\markboth{Genesis 8:9}{Genesis 8:9}\vs{9} The dove could not find a resting place for its feet because water still covered the surface of the entire earth, and so it returned to Noah in the ark. He stretched out his hand, took the dove, and brought it back into the ark.\gva{e}\marginnote{\gva{e}\,The raven flew to and fro, resting on the ark, but did not come into it, unlike the dove that was taken in.}
\markboth{Genesis 8:10}{Genesis 8:10}\vs{10} He waited seven more days and then sent out the dove again from the ark.
\markboth{Genesis 8:11}{Genesis 8:11}\vs{11} When the dove returned to him in the evening, there was a freshly plucked olive leaf in its beak! Noah knew that the waters had receded from the earth.\gva{f}\marginnote{\gva{f}\,Which was a sign that the waters were much diminished: for the olives do not grow on the high mountains.}
\markboth{Genesis 8:12}{Genesis 8:12}\vs{12} He waited another seven days and sent the dove out again, but it did not return to him this time.
\everypar{}

\parshape=0
\markboth{Genesis 8:13}{Genesis 8:13}\vs{13} In Noah's six hundred and first year, in the first day of the first month, the waters had dried up from the earth, and Noah removed the covering from the ark and saw that the surface of the ground was dry.\gva{g}\marginnote{\gva{g}\,Called in Hebrew Abib, part of March and part of April.}
\markboth{Genesis 8:14}{Genesis 8:14}\vs{14} And by the twenty-seventh day of the second month the earth was dry.
\everypar{}

\parshape=0
\markboth{Genesis 8:15}{Genesis 8:15}\vs{15} Then God spoke to Noah and said,
\markboth{Genesis 8:16}{Genesis 8:16}\vs{16} ``Come out of the ark, you, your wife, your sons, and your sons' wives with you.\gva{h}\marginnote{\gva{h}\,Noah declares his obedience, in that he would not leave the ark without God's express commandment, as he did not enter in without the same: the ark being a figure of the Church, in which nothing must be done outside the word of God.}
\markboth{Genesis 8:17}{Genesis 8:17}\vs{17} Bring out with you all the living creatures that are with you. Bring out every living thing, including the birds, animals, and every creeping thing that creeps on the earth. Let them increase and be fruitful and multiply on the earth!''
\everypar{}

\parshape=0
\markboth{Genesis 8:18}{Genesis 8:18}\vs{18} Noah went out along with his sons, his wife, and his sons' wives.
\markboth{Genesis 8:19}{Genesis 8:19}\vs{19} Every living creature, every creeping thing, every bird, and everything that moves on the earth went out of the ark in their groups.
\everypar{}

\parshape=0
\markboth{Genesis 8:20}{Genesis 8:20}\vs{20} Noah built an altar to the Lord. He then took some of every kind of clean animal and clean bird and offered burnt offerings on the altar.\gva{i}\marginnote{\gva{i}\,For sacrifices which were as an exercise of their faith, by which they used to give thanks to God for his benefits.}
\markboth{Genesis 8:21}{Genesis 8:21}\vs{21} And the Lord smelled the soothing aroma and said to himself, ``I will never again curse the ground because of humankind, even though the inclination of their minds is evil from childhood on. I will never again destroy everything that lives, as I have just done.\gva{k}\marginnote{\gva{k}\,That is, by it he showed himself appeased and his anger at rest.}
\everypar{}

\parshape=0
\markboth{Genesis 8:22}{Genesis 8:22}\vs{22} ``While the earth continues to exist, planting time and harvest, cold and heat, summer and winter, and day and night will not cease.''\gva{l}\marginnote{\gva{l}\,The order of nature destroyed by the flood, is restored by God's promise.}

\markboth{Genesis 9:1}{Genesis 9:1}
\Needspace*{8\baselineskip}
\ch{9} \allowchapbreak\hypertarget{ch-genesis-9}{}\lettrine[lines=5]{\color{pentateuch}T}{\textsc{hen}} God blessed Noah and his sons and said to them, ``Be fruitful and multiply and fill the earth.\gva{a}\marginnote{\gva{a}\,God increased them with fruit, and declared to them his counsel as concerning the replenishing of the earth.}\everypar{}
\markboth{Genesis 9:2}{Genesis 9:2}\vs{2} Every living creature of the earth and every bird of the sky will be terrified of you. Everything that creeps on the ground and all the fish of the sea are under your authority.\gva{b}\marginnote{\gva{b}\,By the virtue of this commandment, beasts do not rage as much against man as they would, yea and many serve his purposes by it.}
\markboth{Genesis 9:3}{Genesis 9:3}\vs{3} You may eat any moving thing that lives. As I gave you the green plants, I now give you everything.\gva{c}\marginnote{\gva{c}\,By this permission man may with a good conscience use the creatures of God for his needs.}
\everypar{}

\parshape=0
\markboth{Genesis 9:4}{Genesis 9:4}\vs{4} ``But you must not eat meat with its life (that is, its blood) in it.\gva{d}\marginnote{\gva{d}\,That is, living creatures, and the flesh of beasts that are strangled: and by this all cruelty is forbidden.}
\markboth{Genesis 9:5}{Genesis 9:5}\vs{5} For your lifeblood I will surely exact punishment, from every living creature I will exact punishment. From each person I will exact punishment for the life of the individual since the man was his relative.\gva{e}\marginnote{\gva{e}\,That is, I will take vengeance for your blood.}
\everypar{}

\parshape=0
\markboth{Genesis 9:6}{Genesis 9:6}\vs{6} ``Whoever sheds human blood, by other humans must his blood be shed; for in God's image God has made humankind.\gva{f}\marginnote{\gva{f}\,Not only by the magistrate, but often God raises up one murderer to kill another.}\gva{g}\marginnote{\gva{g}\,Therefore to kill man is to deface God's image, and so injury is not only done to man, but also to God.}
\everypar{}

\parshape=0
\markboth{Genesis 9:7}{Genesis 9:7}\vs{7} ``But as for you, be fruitful and multiply; increase abundantly on the earth and multiply on it.''
\everypar{}

\parshape=0
\markboth{Genesis 9:8}{Genesis 9:8}\vs{8} God said to Noah and his sons,
\markboth{Genesis 9:9}{Genesis 9:9}\vs{9} ``Look. I now confirm my covenant with you and your descendants after you\gva{h}\marginnote{\gva{h}\,To assure you that the world will never again be destroyed by a flood.}\gva{i}\marginnote{\gva{i}\,The children which are not yet born, are comprehended in God's covenant with their fathers.}
\markboth{Genesis 9:10}{Genesis 9:10}\vs{10} and with every living creature that is with you, including the birds, the domestic animals, and every living creature of the earth with you, all those that came out of the ark with you---every living creature of the earth.
\markboth{Genesis 9:11}{Genesis 9:11}\vs{11} I confirm my covenant with you: Never again will all living things be wiped out by the waters of a flood; never again will a flood destroy the earth.''
\everypar{}

\parshape=0
\markboth{Genesis 9:12}{Genesis 9:12}\vs{12} And God said, ``This is the guarantee of the covenant I am making with you and every living creature with you, a covenant for all subsequent generations:
\markboth{Genesis 9:13}{Genesis 9:13}\vs{13} I will place my rainbow in the clouds, and it will become a guarantee of the covenant between me and the earth.\gva{k}\marginnote{\gva{k}\,By this we see that signs or ordinances should not be separate from the word.}
\markboth{Genesis 9:14}{Genesis 9:14}\vs{14} Whenever I bring clouds over the earth and the rainbow appears in the clouds,
\markboth{Genesis 9:15}{Genesis 9:15}\vs{15} then I will remember my covenant with you and with all living creatures of all kinds. Never again will the waters become a flood and destroy all living things.\gva{l}\marginnote{\gva{l}\,When men see my bow in the sky, they will know that I have not forgotten my covenant with them.}
\markboth{Genesis 9:16}{Genesis 9:16}\vs{16} When the rainbow is in the clouds, I will notice it and remember the perpetual covenant between God and all living creatures of all kinds that are on the earth.''
\everypar{}

\parshape=0
\markboth{Genesis 9:17}{Genesis 9:17}\vs{17} So God said to Noah, ``This is the guarantee of the covenant that I am confirming between me and all living things that are on the earth.''\gva{m}\marginnote{\gva{m}\,God repeats this often to confirm Noah's faith even more.}
\everypar{}

\parshape=0
\markboth{Genesis 9:18}{Genesis 9:18}\vs{18} The sons of Noah who came out of the ark were Shem, Ham, and Japheth. (Now Ham was the father of Canaan.)
\markboth{Genesis 9:19}{Genesis 9:19}\vs{19} These were the three sons of Noah, and from them the whole earth was populated.\gva{n}\marginnote{\gva{n}\,This declares what the virtue of God's blessing was, when he said, increase and bring forth in Genesis 1:28 .}
\everypar{}

\parshape=0
\markboth{Genesis 9:20}{Genesis 9:20}\vs{20} Noah, a man of the soil, began to plant a vineyard.
\markboth{Genesis 9:21}{Genesis 9:21}\vs{21} When he drank some of the wine, he got drunk and uncovered himself inside his tent.\gva{o}\marginnote{\gva{o}\,This is set before us to show what a horrible thing drunkenness is.}
\markboth{Genesis 9:22}{Genesis 9:22}\vs{22} Ham, the father of Canaan, saw his father's nakedness and told his two brothers who were outside.\gva{p}\marginnote{\gva{p}\,Of whom came the Canaanites that wicked nation, who were also cursed by God.}\gva{q}\marginnote{\gva{q}\,In derision and contempt of his father.}
\markboth{Genesis 9:23}{Genesis 9:23}\vs{23} Shem and Japheth took the garment and placed it on their shoulders. Then they walked in backwards and covered up their father's nakedness. Their faces were turned the other way so they did not see their father's nakedness.
\everypar{}

\parshape=0
\markboth{Genesis 9:24}{Genesis 9:24}\vs{24} When Noah awoke from his drunken stupor he learned what his youngest son had done to him.
\everypar{}

\parshape=0
\markboth{Genesis 9:25}{Genesis 9:25}\vs{25} So he said, ``Cursed be Canaan! The lowest of slaves he will be to his brothers.''\gva{r}\marginnote{\gva{r}\,He pronounces as a prophet the curse of God against all those who do not honour their parents: for Ham and his posterity were cursed.}\gva{s}\marginnote{\gva{s}\,That is, a most vile slave.}
\everypar{}

\parshape=0
\markboth{Genesis 9:26}{Genesis 9:26}\vs{26} He also said, ``Worthy of praise is the Lord, the God of Shem! May Canaan be the slave of Shem!
\everypar{}

\parshape=0
\markboth{Genesis 9:27}{Genesis 9:27}\vs{27} May God enlarge Japheth's territory and numbers! May he live in the tents of Shem and may Canaan be the slave of Japheth!''\gva{t}\marginnote{\gva{t}\,He declares that the Gentiles, who came from Japheth, and were separated from the Church, should be joined to the same by the persuasion of God's Spirit, and preaching of the gospel.}
\everypar{}

\parshape=0
\markboth{Genesis 9:28}{Genesis 9:28}\vs{28} After the flood Noah lived 350 years.
\markboth{Genesis 9:29}{Genesis 9:29}\vs{29} The entire lifetime of Noah was 950 years, and then he died.

\markboth{Genesis 10:1}{Genesis 10:1}
\Needspace*{8\baselineskip}
\ch{10} \allowchapbreak\hypertarget{ch-genesis-10}{}\lettrine[lines=5]{\color{pentateuch}T}{\textsc{his}} is the account of Noah's sons: Shem, Ham, and Japheth. Sons were born to them after the flood.\gva{a}\marginnote{\gva{a}\,These generations are here recited, partly to declare the marvellous increase, and also to set forth their great forgetfulness of God's grace towards their fathers.}\everypar{}
\\\indent\markboth{Genesis 10:2}{Genesis 10:2}\vs{2} The sons of Japheth were Gomer, Magog, Madai, Javan, Tubal, Meshech, and Tiras.\gva{b}\marginnote{\gva{b}\,Of Madai and Javan came the Medes and Greeks.}
\markboth{Genesis 10:3}{Genesis 10:3}\vs{3} The sons of Gomer were Ashkenaz, Riphath, and Togarmah.
\markboth{Genesis 10:4}{Genesis 10:4}\vs{4} The sons of Javan were Elishah, Tarshish, the Kittim, and the Dodanim.
\markboth{Genesis 10:5}{Genesis 10:5}\vs{5} From these the coastlands of the nations were separated into their lands, every one according to its language, according to their families, by their nations.\gva{c}\marginnote{\gva{c}\,So do the Jews call all countries which are separated from them by sea, as Greece, Italy etc, which were given to the children of Japheth, of whom came the Gentiles.}
\everypar{}

\parshape=0
\markboth{Genesis 10:6}{Genesis 10:6}\vs{6} The sons of Ham were Cush, Mizraim, Put, and Canaan.\gva{d}\marginnote{\gva{d}\,Of Cush and Mizraim came the Ethiopians and Egyptians.}
\markboth{Genesis 10:7}{Genesis 10:7}\vs{7} The sons of Cush were Seba, Havilah, Sabtah, Raamah, and Sabteca. The sons of Raamah were Sheba and Dedan.
\everypar{}

\parshape=0
\markboth{Genesis 10:8}{Genesis 10:8}\vs{8} Cush was the father of Nimrod; he began to be a valiant warrior on the earth.\gva{e}\marginnote{\gva{e}\,Meaning, a cruel oppressor and tyrant.}
\markboth{Genesis 10:9}{Genesis 10:9}\vs{9} He was a mighty hunter before the Lord. (That is why it is said, ``Like Nimrod, a mighty hunter before the Lord.'')\gva{f}\marginnote{\gva{f}\,His tyranny came into a proverb as hated both by God and man: for he did not cease to commit cruelty even in God's presence.}
\markboth{Genesis 10:10}{Genesis 10:10}\vs{10} The primary regions of his kingdom were Babel, Erech, Akkad, and Calneh in the land of Shinar.\gva{g}\marginnote{\gva{g}\,For there was another city in Egypt, called Babel.}
\markboth{Genesis 10:11}{Genesis 10:11}\vs{11} From that land he went to Assyria, where he built Nineveh, Rehoboth Ir, Calah,
\markboth{Genesis 10:12}{Genesis 10:12}\vs{12} and Resen, which is between Nineveh and the great city Calah.
\everypar{}

\parshape=0
\markboth{Genesis 10:13}{Genesis 10:13}\vs{13} Mizraim was the father of the Ludites, Anamites, Lehabites, Naphtuhites,\gva{h}\marginnote{\gva{h}\,Of Lud came the Lydians.}
\markboth{Genesis 10:14}{Genesis 10:14}\vs{14} Pathrusites, Casluhites (from whom the Philistines came), and Caphtorites.
\everypar{}

\parshape=0
\markboth{Genesis 10:15}{Genesis 10:15}\vs{15} Canaan was the father of Sidon his firstborn, Heth,
\markboth{Genesis 10:16}{Genesis 10:16}\vs{16} the Jebusites, Amorites, Girgashites,
\markboth{Genesis 10:17}{Genesis 10:17}\vs{17} Hivites, Arkites, Sinites,
\markboth{Genesis 10:18}{Genesis 10:18}\vs{18} Arvadites, Zemarites, and Hamathites. Eventually the families of the Canaanites were scattered
\markboth{Genesis 10:19}{Genesis 10:19}\vs{19} and the borders of Canaan extended from Sidon all the way to Gerar as far as Gaza, and all the way to Sodom, Gomorrah, Admah, and Zeboyim, as far as Lasha.
\markboth{Genesis 10:20}{Genesis 10:20}\vs{20} These are the sons of Ham, according to their families, according to their languages, by their lands, and by their nations.
\everypar{}

\parshape=0
\markboth{Genesis 10:21}{Genesis 10:21}\vs{21} And sons were also born to Shem (the older brother of Japheth), the father of all the sons of Eber.\gva{i}\marginnote{\gva{i}\,In his stock the Church was preserved: therefore Moses stops speaking of Japheth and Ham, and speaks of Shem extensively.}\gva{k}\marginnote{\gva{k}\,Of whom came the Hebrews or Jews.}
\everypar{}

\parshape=0
\markboth{Genesis 10:22}{Genesis 10:22}\vs{22} The sons of Shem were Elam, Asshur, Arphaxad, Lud, and Aram.
\markboth{Genesis 10:23}{Genesis 10:23}\vs{23} The sons of Aram were Uz, Hul, Gether, and Mash.
\markboth{Genesis 10:24}{Genesis 10:24}\vs{24} Arphaxad was the father of Shelah, and Shelah was the father of Eber.
\markboth{Genesis 10:25}{Genesis 10:25}\vs{25} Two sons were born to Eber: One was named Peleg because in his days the earth was divided, and his brother's name was Joktan.\gva{l}\marginnote{\gva{l}\,This division came by the diversity of language, as appears in Genesis 11:9 .}
\markboth{Genesis 10:26}{Genesis 10:26}\vs{26} Joktan was the father of Almodad, Sheleph, Hazarmaveth, Jerah,
\markboth{Genesis 10:27}{Genesis 10:27}\vs{27} Hadoram, Uzal, Diklah,
\markboth{Genesis 10:28}{Genesis 10:28}\vs{28} Obal, Abimael, Sheba,
\markboth{Genesis 10:29}{Genesis 10:29}\vs{29} Ophir, Havilah, and Jobab. All these were sons of Joktan.
\markboth{Genesis 10:30}{Genesis 10:30}\vs{30} Their dwelling place was from Mesha all the way to Sephar in the eastern hills.
\markboth{Genesis 10:31}{Genesis 10:31}\vs{31} These are the sons of Shem according to their families, according to their languages, by their lands, and according to their nations.
\everypar{}

\parshape=0
\markboth{Genesis 10:32}{Genesis 10:32}\vs{32} These are the families of the sons of Noah, according to their genealogies, by their nations, and from these the nations spread over the earth after the flood.

\markboth{Genesis 11:1}{Genesis 11:1}
\Needspace*{8\baselineskip}
\ch{11} \allowchapbreak\hypertarget{ch-genesis-11}{}\lettrine[lines=5]{\color{pentateuch}T}{\textsc{he}} whole earth had a common language and a common vocabulary.\everypar{}
\markboth{Genesis 11:2}{Genesis 11:2}\vs{2} When the people moved eastward, they found a plain in Shinar and settled there.\gva{a}\marginnote{\gva{a}\,One hundred and thirty years after the flood.}\gva{b}\marginnote{\gva{b}\,That is, Nimrod and his company.}\gva{c}\marginnote{\gva{c}\,That is, from Armenia where the ark stayed.}\gva{d}\marginnote{\gva{d}\,Which was afterward called Chaldea.}
\markboth{Genesis 11:3}{Genesis 11:3}\vs{3} Then they said to one another, ``Come, let's make bricks and bake them thoroughly.'' (They had brick instead of stone and tar instead of mortar.)
\markboth{Genesis 11:4}{Genesis 11:4}\vs{4} Then they said, ``Come, let's build ourselves a city and a tower with its top in the heavens so that we may make a name for ourselves. Otherwise we will be scattered across the face of the entire earth.''\gva{e}\marginnote{\gva{e}\,They were moved with pride and ambition, preferring their own glory to God's honour.}
\everypar{}

\parshape=0
\markboth{Genesis 11:5}{Genesis 11:5}\vs{5} But the Lord came down to see the city and the tower that the people had started building.\gva{f}\marginnote{\gva{f}\,Meaning, that he declared by effect, that he knew their wicked enterprise; for God's power is everywhere, and neither ascends nor descends.}
\markboth{Genesis 11:6}{Genesis 11:6}\vs{6} And the Lord said, ``If as one people all sharing a common language they have begun to do this, then nothing they plan to do will be beyond them.\gva{g}\marginnote{\gva{g}\,God speaks this in derision, because of their foolish persuasion and enterprise.}
\markboth{Genesis 11:7}{Genesis 11:7}\vs{7} Come, let's go down and confuse their language so they won't be able to understand each other.''\gva{h}\marginnote{\gva{h}\,He speaks as though he took counsel with his own wisdom and power: that is, with the Son and holy Spirit: signifying the greatness and certainty of the punishment.}\gva{i}\marginnote{\gva{i}\,By this great plague of the confusion of tongues appears God's horrible judgment against man's pride and vain glory.}
\everypar{}

\parshape=0
\markboth{Genesis 11:8}{Genesis 11:8}\vs{8} So the Lord scattered them from there across the face of the entire earth, and they stopped building the city.
\markboth{Genesis 11:9}{Genesis 11:9}\vs{9} That is why its name was called Babel---because there the Lord confused the language of the entire world, and from there the Lord scattered them across the face of the entire earth.
\everypar{}

\parshape=0
\markboth{Genesis 11:10}{Genesis 11:10}\vs{10} This is the account of Shem. Shem was 100 years old when he became the father of Arphaxad, two years after the flood.\gva{k}\marginnote{\gva{k}\,He returns to the genealogy of Shem, to come to the history of Abram, in which the Church of God is described, which is Moses' principle purpose.}
\markboth{Genesis 11:11}{Genesis 11:11}\vs{11} And after becoming the father of Arphaxad, Shem lived 500 years and had other sons and daughters.
\everypar{}

\parshape=0
\markboth{Genesis 11:12}{Genesis 11:12}\vs{12} When Arphaxad had lived 35 years, he became the father of Shelah.
\markboth{Genesis 11:13}{Genesis 11:13}\vs{13} And after he became the father of Shelah, Arphaxad lived 403 years and had other sons and daughters.
\everypar{}

\parshape=0
\markboth{Genesis 11:14}{Genesis 11:14}\vs{14} When Shelah had lived 30 years, he became the father of Eber.
\markboth{Genesis 11:15}{Genesis 11:15}\vs{15} And after he became the father of Eber, Shelah lived 403 years and had other sons and daughters.
\everypar{}

\parshape=0
\markboth{Genesis 11:16}{Genesis 11:16}\vs{16} When Eber had lived 34 years, he became the father of Peleg.
\markboth{Genesis 11:17}{Genesis 11:17}\vs{17} And after he became the father of Peleg, Eber lived 430 years and had other sons and daughters.
\everypar{}

\parshape=0
\markboth{Genesis 11:18}{Genesis 11:18}\vs{18} When Peleg had lived 30 years, he became the father of Reu.
\markboth{Genesis 11:19}{Genesis 11:19}\vs{19} And after he became the father of Reu, Peleg lived 209 years and had other sons and daughters.
\everypar{}

\parshape=0
\markboth{Genesis 11:20}{Genesis 11:20}\vs{20} When Reu had lived 32 years, he became the father of Serug.
\markboth{Genesis 11:21}{Genesis 11:21}\vs{21} And after he became the father of Serug, Reu lived 207 years and had other sons and daughters.
\everypar{}

\parshape=0
\markboth{Genesis 11:22}{Genesis 11:22}\vs{22} When Serug had lived 30 years, he became the father of Nahor.
\markboth{Genesis 11:23}{Genesis 11:23}\vs{23} And after he became the father of Nahor, Serug lived 200 years and had other sons and daughters.
\everypar{}

\parshape=0
\markboth{Genesis 11:24}{Genesis 11:24}\vs{24} When Nahor had lived 29 years, he became the father of Terah.
\markboth{Genesis 11:25}{Genesis 11:25}\vs{25} And after he became the father of Terah, Nahor lived 119 years and had other sons and daughters.
\everypar{}

\parshape=0
\markboth{Genesis 11:26}{Genesis 11:26}\vs{26} When Terah had lived 70 years, he became the father of Abram, Nahor, and Haran.
\everypar{}

\parshape=0
\markboth{Genesis 11:27}{Genesis 11:27}\vs{27} This is the account of Terah. Terah became the father of Abram, Nahor, and Haran. And Haran became the father of Lot.
\markboth{Genesis 11:28}{Genesis 11:28}\vs{28} Haran died in the land of his birth, in Ur of the Chaldeans, while his father Terah was still alive.
\markboth{Genesis 11:29}{Genesis 11:29}\vs{29} And Abram and Nahor took wives for themselves. The name of Abram's wife was Sarai. And the name of Nahor's wife was Milcah; she was the daughter of Haran, who was the father of both Milcah and Iscah.\gva{m}\marginnote{\gva{m}\,Some think that this Iscah was Sarai.}
\markboth{Genesis 11:30}{Genesis 11:30}\vs{30} But Sarai was barren; she had no children.
\everypar{}

\parshape=0
\markboth{Genesis 11:31}{Genesis 11:31}\vs{31} Terah took his son Abram, his grandson Lot (the son of Haran), and his daughter-in-law Sarai, his son Abram's wife, and with them he set out from Ur of the Chaldeans to go to Canaan. When they came to Haran, they settled there.\gva{n}\marginnote{\gva{n}\,Though the oracle of God came to Abram, yet the honour is given to Terah, because he was the father.}\gva{o}\marginnote{\gva{o}\,Which was a city of Mesopotamia.}
\markboth{Genesis 11:32}{Genesis 11:32}\vs{32} The lifetime of Terah was 205 years, and he died in Haran.

\markboth{Genesis 12:1}{Genesis 12:1}
\Needspace*{8\baselineskip}
\ch{12} \allowchapbreak\hypertarget{ch-genesis-12}{}\lettrine[lines=5]{\color{pentateuch}N}{\textsc{ow}} the Lord said to Abram, ``Go out from your country, your relatives, and your father's household to the land that I will show you.\gva{a}\marginnote{\gva{a}\,From the flood to this time were four hundred and twenty-three years.}\gva{b}\marginnote{\gva{b}\,In appointing him no certain place, he proves so much more his faith and obedience.}\everypar{}
\\\indent\markboth{Genesis 12:2}{Genesis 12:2}\vs{2} Then I will make you into a great nation, and I will bless you, and I will make your name great, so that you will exemplify divine blessing.\gva{c}\marginnote{\gva{c}\,The world shall recover by your seed, which is Christ, the blessing which they lost in Adam.}
\everypar{}

\parshape=0
\markboth{Genesis 12:3}{Genesis 12:3}\vs{3} I will bless those who bless you, but the one who treats you lightly I must curse, so that all the families of the earth may receive blessing through you.''
\everypar{}

\parshape=0
\markboth{Genesis 12:4}{Genesis 12:4}\vs{4} So Abram left, just as the Lord had told him to do, and Lot went with him. (Now Abram was 75 years old when he departed from Haran.)
\markboth{Genesis 12:5}{Genesis 12:5}\vs{5} And Abram took his wife Sarai, his nephew Lot, and all the possessions they had accumulated and the people they had acquired in Haran, and they left for the land of Canaan. They entered the land of Canaan.\gva{d}\marginnote{\gva{d}\,Meaning servants as well as cattle.}
\everypar{}

\parshape=0
\markboth{Genesis 12:6}{Genesis 12:6}\vs{6} Abram traveled through the land as far as the oak tree of Moreh at Shechem. (At that time the Canaanites were in the land.)\gva{e}\marginnote{\gva{e}\,He wandered to and fro in the land before he could find a settling place: thus God exercises the faith of his children.}\gva{f}\marginnote{\gva{f}\,Which was a cruel and rebellious nation, by whom God kept his in continual exercise.}
\markboth{Genesis 12:7}{Genesis 12:7}\vs{7} The Lord appeared to Abram and said, ``To your descendants I will give this land.'' So Abram built an altar there to the Lord, who had appeared to him.\gva{g}\marginnote{\gva{g}\,It was not enough for him to worship God in his heart, but it was expedient to declare by outward profession his faith before men, of which this altar was a sign.}
\everypar{}

\parshape=0
\markboth{Genesis 12:8}{Genesis 12:8}\vs{8} Then he moved from there to the hill country east of Bethel and pitched his tent, with Bethel on the west and Ai on the east. There he built an altar to the Lord and worshiped the Lord.\gva{h}\marginnote{\gva{h}\,Because of the troubles that he had among that wicked people.}\gva{i}\marginnote{\gva{i}\,And so served the true God, and renounced all idolatry.}
\markboth{Genesis 12:9}{Genesis 12:9}\vs{9} Abram continually journeyed by stages down to the Negev.\gva{k}\marginnote{\gva{k}\,Thus the children of God may look for no rest in this world, but must wait for the heavenly rest and quietness.}
\everypar{}

\parshape=0
\markboth{Genesis 12:10}{Genesis 12:10}\vs{10} There was a famine in the land, so Abram went down to Egypt to stay for a while because the famine was severe.\gva{l}\marginnote{\gva{l}\,This was a new trial of Abram's faith: by which we see that the end of one affliction is the beginning of another.}
\markboth{Genesis 12:11}{Genesis 12:11}\vs{11} As he approached Egypt, he said to his wife Sarai, ``Look, I know that you are a beautiful woman.
\markboth{Genesis 12:12}{Genesis 12:12}\vs{12} When the Egyptians see you they will say, `This is his wife.' Then they will kill me but will keep you alive.
\markboth{Genesis 12:13}{Genesis 12:13}\vs{13} So tell them you are my sister so that it may go well for me because of you and my life will be spared on account of you.''\gva{m}\marginnote{\gva{m}\,By this we learn not to use unlawful means nor to put others in danger to save ourselves, Genesis 12:20 . Though it may appear that Abram did not fear death, so much as dying without children, he acts as though God's promise had not taken place; in which appeared a weak faith.}
\everypar{}

\parshape=0
\markboth{Genesis 12:14}{Genesis 12:14}\vs{14} When Abram entered Egypt, the Egyptians saw that the woman was very beautiful.
\markboth{Genesis 12:15}{Genesis 12:15}\vs{15} When Pharaoh's officials saw her, they praised her to Pharaoh. So Abram's wife was taken into the household of Pharaoh,\gva{n}\marginnote{\gva{n}\,To be his wife.}
\markboth{Genesis 12:16}{Genesis 12:16}\vs{16} and he did treat Abram well on account of her. Abram received sheep and cattle, male donkeys, male servants, female servants, female donkeys, and camels.
\everypar{}

\parshape=0
\markboth{Genesis 12:17}{Genesis 12:17}\vs{17} But the Lord struck Pharaoh and his household with severe diseases because of Sarai, Abram's wife.\gva{o}\marginnote{\gva{o}\,The Lord took the defence of this poor stranger against a mighty king: and as he is ever careful over his, so did he preserve Sarai.}
\markboth{Genesis 12:18}{Genesis 12:18}\vs{18} So Pharaoh summoned Abram and said, ``What is this you have done to me? Why didn't you tell me that she was your wife?
\markboth{Genesis 12:19}{Genesis 12:19}\vs{19} Why did you say, `She is my sister,' so that I took her to be my wife? Now, here is your wife. Take her and go!''
\markboth{Genesis 12:20}{Genesis 12:20}\vs{20} Pharaoh gave his men orders about Abram, and so they expelled him, along with his wife and all his possessions.\gva{p}\marginnote{\gva{p}\,To the intent that none should hurt him either in his person or goods.}

\markboth{Genesis 13:1}{Genesis 13:1}
\Needspace*{8\baselineskip}
\ch{13} \allowchapbreak\hypertarget{ch-genesis-13}{}\lettrine[lines=5]{\color{pentateuch}S}{\textsc{o}} Abram went up from Egypt into the Negev. He took his wife and all his possessions with him, as well as Lot.\gva{a}\marginnote{\gva{a}\,His great riches gotten in Egypt, did not hinder him in following his vocation.}\everypar{}
\markboth{Genesis 13:2}{Genesis 13:2}\vs{2} (Now Abram was very wealthy in livestock, silver, and gold.)
\\\indent\markboth{Genesis 13:3}{Genesis 13:3}\vs{3} And he journeyed from place to place from the Negev as far as Bethel. He returned to the place where he had pitched his tent at the beginning, between Bethel and Ai.\gva{b}\marginnote{\gva{b}\,He calls the place by the name which was later given to it, Genesis 23:19 .}
\markboth{Genesis 13:4}{Genesis 13:4}\vs{4} This was the place where he had first built the altar, and there Abram worshiped the Lord.
\everypar{}

\parshape=0
\markboth{Genesis 13:5}{Genesis 13:5}\vs{5} Now Lot, who was traveling with Abram, also had flocks, herds, and tents.
\markboth{Genesis 13:6}{Genesis 13:6}\vs{6} But the land could not support them while they were living side by side. Because their possessions were so great, they were not able to live alongside one another.\gva{c}\marginnote{\gva{c}\,This inconvenience came by their riches, which break friendships and the bounds of nature.}
\markboth{Genesis 13:7}{Genesis 13:7}\vs{7} So there were quarrels between Abram's herdsmen and Lot's herdsmen. (Now the Canaanites and the Perizzites were living in the land at that time.)\gva{d}\marginnote{\gva{d}\,Who seeing their contention, might blaspheme God and destroy them.}
\everypar{}

\parshape=0
\markboth{Genesis 13:8}{Genesis 13:8}\vs{8} Abram said to Lot, ``Let there be no quarreling between me and you, and between my herdsmen and your herdsmen, for we are close relatives.\gva{e}\marginnote{\gva{e}\,He cuts off the opportunity for contention: therefore the evil ceases.}
\markboth{Genesis 13:9}{Genesis 13:9}\vs{9} Is not the whole land before you? Separate yourself now from me. If you go to the left, then I'll go to the right, but if you go to the right, then I'll go to the left.''\gva{f}\marginnote{\gva{f}\,Abram resigns his own right to buy peace.}
\everypar{}

\parshape=0
\markboth{Genesis 13:10}{Genesis 13:10}\vs{10} Lot looked up and saw the whole region of the Jordan. He noticed that all of it was well watered (this was before the Lord obliterated Sodom and Gomorrah) like the garden of the Lord, like the land of Egypt, all the way to Zoar.\gva{g}\marginnote{\gva{g}\,Which was in Eden, Genesis 2:10 .}
\everypar{}

\parshape=0
\markboth{Genesis 13:11}{Genesis 13:11}\vs{11} Lot chose for himself the whole region of the Jordan and traveled toward the east. So the relatives separated from each other.\gva{h}\marginnote{\gva{h}\,This was done by God's providence, that only Abram and his seed might dwell in the land of Canaan.}
\markboth{Genesis 13:12}{Genesis 13:12}\vs{12} Abram settled in the land of Canaan, but Lot settled among the cities of the Jordan plain and pitched his tents next to Sodom.
\markboth{Genesis 13:13}{Genesis 13:13}\vs{13} (Now the people of Sodom were extremely wicked rebels against the Lord.)\gva{i}\marginnote{\gva{i}\,Lot thinking to get paradise, found hell.}
\everypar{}

\parshape=0
\markboth{Genesis 13:14}{Genesis 13:14}\vs{14} After Lot had departed, the Lord said to Abram, ``Look from the place where you stand to the north, south, east, and west.\gva{k}\marginnote{\gva{k}\,The Lord comforted him, lest he should have taken thought for the departure of his nephew.}
\markboth{Genesis 13:15}{Genesis 13:15}\vs{15} I will give all the land that you see to you and your descendants forever.\gva{l}\marginnote{\gva{l}\,Meaning a long time, and till the coming of Christ as in Exodus 12:14 ; Exodus 21:6 , De 15:17 and spiritually this refers to the true children of Abram born according to the promise, and not according to the flesh, which are heirs of the true land of Canaan.}
\markboth{Genesis 13:16}{Genesis 13:16}\vs{16} And I will make your descendants like the dust of the earth, so that if anyone is able to count the dust of the earth, then your descendants also can be counted.
\markboth{Genesis 13:17}{Genesis 13:17}\vs{17} Get up and walk throughout the land, for I will give it to you.''
\everypar{}

\parshape=0
\markboth{Genesis 13:18}{Genesis 13:18}\vs{18} So Abram moved his tents and went to live by the oaks of Mamre in Hebron, and he built an altar to the Lord there.

\markboth{Genesis 14:1}{Genesis 14:1}
\Needspace*{8\baselineskip}
\ch{14} \allowchapbreak\hypertarget{ch-genesis-14}{}\lettrine[lines=5]{\color{pentateuch}A}{\textsc{t}} that time Amraphel king of Shinar, Arioch king of Ellasar, Kedorlaomer king of Elam, and Tidal king of nations\gva{a}\marginnote{\gva{a}\,That is, of Babylon: by kings here, meaning, them that were governors of cities.}\gva{b}\marginnote{\gva{b}\,Of a people gathered from various countries.}\everypar{}
\markboth{Genesis 14:2}{Genesis 14:2}\vs{2} went to war against Bera king of Sodom, Birsha king of Gomorrah, Shinab king of Admah, Shemeber king of Zeboyim, and the king of Bela (that is, Zoar).
\markboth{Genesis 14:3}{Genesis 14:3}\vs{3} These last five kings joined forces in the Valley of Siddim (that is, the Salt Sea).\gva{c}\marginnote{\gva{c}\,Ambition is the chief cause of wars among princes.}\gva{d}\marginnote{\gva{d}\,Called also the dead sea, or the lake Asphaltite, near Sodom and Gomorrah.}
\markboth{Genesis 14:4}{Genesis 14:4}\vs{4} For twelve years they had served Kedorlaomer, but in the thirteenth year they rebelled.
\markboth{Genesis 14:5}{Genesis 14:5}\vs{5} In the fourteenth year, Kedorlaomer and the kings who were his allies came and defeated the Rephaites in Ashteroth Karnaim, the Zuzites in Ham, the Emites in Shaveh Kiriathaim,
\markboth{Genesis 14:6}{Genesis 14:6}\vs{6} and the Horites in their hill country of Seir, as far as El Paran, which is near the desert.
\markboth{Genesis 14:7}{Genesis 14:7}\vs{7} Then they attacked En Mishpat (that is, Kadesh) again, and they conquered all the territory of the Amalekites, as well as the Amorites who were living in Hazezon Tamar.
\everypar{}

\parshape=0
\markboth{Genesis 14:8}{Genesis 14:8}\vs{8} Then the king of Sodom, the king of Gomorrah, the king of Admah, the king of Zeboyim, and the king of Bela (that is, Zoar) went out and prepared for battle. In the Valley of Siddim they met
\markboth{Genesis 14:9}{Genesis 14:9}\vs{9} Kedorlaomer king of Elam, Tidal king of nations, Amraphel king of Shinar, and Arioch king of Ellasar. Four kings fought against five.
\markboth{Genesis 14:10}{Genesis 14:10}\vs{10} Now the Valley of Siddim was full of tar pits. When the kings of Sodom and Gomorrah fled, they fell into them, but some survivors fled to the hills.\gva{e}\marginnote{\gva{e}\,And afterward was overwhelmed with water, and so was called the salt sea.}
\markboth{Genesis 14:11}{Genesis 14:11}\vs{11} The four victorious kings took all the possessions and food of Sodom and Gomorrah and left.
\markboth{Genesis 14:12}{Genesis 14:12}\vs{12} They also took Abram's nephew Lot and his possessions when they left, for Lot was living in Sodom.\gva{f}\marginnote{\gva{f}\,The godly are plagued many times with the wicked: therefore their company is dangerous.}
\everypar{}

\parshape=0
\markboth{Genesis 14:13}{Genesis 14:13}\vs{13} A fugitive came and told Abram the Hebrew. Now Abram was living by the oaks of Mamre the Amorite, the brother of Eshcol and Aner. (All these were allied by treaty with Abram.)\gva{g}\marginnote{\gva{g}\,God removed them to join Abram, and preserves him from their idolatry and superstitions.}
\markboth{Genesis 14:14}{Genesis 14:14}\vs{14} When Abram heard that his nephew had been taken captive, he mobilized his 318 trained men who had been born in his household, and he pursued the invaders as far as Dan.
\markboth{Genesis 14:15}{Genesis 14:15}\vs{15} Then, during the night, Abram divided his forces against them and defeated them. He chased them as far as Hobah, which is north of Damascus.
\markboth{Genesis 14:16}{Genesis 14:16}\vs{16} He retrieved all the stolen property. He also brought back his nephew Lot and his possessions, as well as the women and the rest of the people.
\everypar{}

\parshape=0
\markboth{Genesis 14:17}{Genesis 14:17}\vs{17} After Abram returned from defeating Kedorlaomer and the kings who were with him, the king of Sodom went out to meet Abram in the Valley of Shaveh (known as the King's Valley).
\markboth{Genesis 14:18}{Genesis 14:18}\vs{18} Melchizedek king of Salem brought out bread and wine. (Now he was the priest of the Most High God.)\gva{h}\marginnote{\gva{h}\,For Abram and his soldiers refreshment, not to offer sacrifice.}
\everypar{}

\parshape=0
\markboth{Genesis 14:19}{Genesis 14:19}\vs{19} He blessed Abram, saying, ``Blessed be Abram by the Most High God, Creator of heaven and earth.\gva{i}\marginnote{\gva{i}\,Melchizedek fed Abram, declared himself to represent a king, and he blessed him as the high priest.}
\everypar{}

\parshape=0
\markboth{Genesis 14:20}{Genesis 14:20}\vs{20} Worthy of praise is the Most High God, who delivered your enemies into your hand.'' Abram gave Melchizedek a tenth of everything.
\everypar{}

\parshape=0
\markboth{Genesis 14:21}{Genesis 14:21}\vs{21} Then the king of Sodom said to Abram, ``Give me the people and take the possessions for yourself.''
\markboth{Genesis 14:22}{Genesis 14:22}\vs{22} But Abram replied to the king of Sodom, ``I raise my hand to the Lord, the Most High God, Creator of heaven and earth, and vow
\markboth{Genesis 14:23}{Genesis 14:23}\vs{23} that I will take nothing belonging to you, not even a thread or the strap of a sandal. That way you can never say, `It is I who made Abram rich.'
\markboth{Genesis 14:24}{Genesis 14:24}\vs{24} I will take nothing except compensation for what the young men have eaten. As for the share of the men who went with me---Aner, Eshcol, and Mamre---let them take their share.''\gva{k}\marginnote{\gva{k}\,He did not want his liberality to be hurtful to others.}

\markboth{Genesis 15:1}{Genesis 15:1}
\Needspace*{8\baselineskip}
\ch{15} \allowchapbreak\hypertarget{ch-genesis-15}{}\lettrine[lines=5]{\color{pentateuch}A}{\textsc{fter}} these things the Lord's message came to Abram in a vision: ``Fear not, Abram! I am your shield and the one who will reward you in great abundance.''\everypar{}
\\\indent\markboth{Genesis 15:2}{Genesis 15:2}\vs{2} But Abram said, ``O Sovereign Lord, what will you give me since I continue to be childless, and my heir is Eliezer of Damascus?''\gva{a}\marginnote{\gva{a}\,His fear was not only lest he should not have children, but lest the promise of the blessed seed should not be accomplished in him.}
\markboth{Genesis 15:3}{Genesis 15:3}\vs{3} Abram added, ``Since you have not given me a descendant, then look, one born in my house will be my heir!''
\everypar{}

\parshape=0
\markboth{Genesis 15:4}{Genesis 15:4}\vs{4} But look, the Lord's message came to him: ``This man will not be your heir, but instead a son who comes from your own body will be your heir.''
\markboth{Genesis 15:5}{Genesis 15:5}\vs{5} The Lord took him outside and said, ``Gaze into the sky and count the stars---if you are able to count them!'' Then he said to him, ``So will your descendants be.''
\everypar{}

\parshape=0
\markboth{Genesis 15:6}{Genesis 15:6}\vs{6} Abram believed the Lord, and the Lord credited it as righteousness to him.
\everypar{}

\parshape=0
\markboth{Genesis 15:7}{Genesis 15:7}\vs{7} The Lord said to him, ``I am the Lord who brought you out from Ur of the Chaldeans to give you this land to possess.''
\markboth{Genesis 15:8}{Genesis 15:8}\vs{8} But Abram said, ``O Sovereign Lord, by what can I know that I am to possess it?''\gva{b}\marginnote{\gva{b}\,This is a particular motion of God's Spirit, which is not lawful for all to follow, in asking signs: but was permitted for some by a peculiar motion, as to Gideon and Ezekiel.}
\everypar{}

\parshape=0
\markboth{Genesis 15:9}{Genesis 15:9}\vs{9} The Lord said to him, ``Take for me a heifer, a goat, and a ram, each three years old, along with a dove and a young pigeon.''
\markboth{Genesis 15:10}{Genesis 15:10}\vs{10} So Abram took all these for him and then cut them in two and placed each half opposite the other, but he did not cut the birds in half.\gva{c}\marginnote{\gva{c}\,This was the old custom in making covenants, Jeremiah 39:18 , to which God added these conditions, that Abram's posterity would be as torn in pieces, but after they would be rejoined: also that it would be assaulted, but yet delivered.}
\markboth{Genesis 15:11}{Genesis 15:11}\vs{11} When birds of prey came down on the carcasses, Abram drove them away.
\everypar{}

\parshape=0
\markboth{Genesis 15:12}{Genesis 15:12}\vs{12} When the sun went down, Abram fell sound asleep, and great terror overwhelmed him.
\markboth{Genesis 15:13}{Genesis 15:13}\vs{13} Then the Lord said to Abram, ``Know for certain that your descendants will be strangers in a foreign country. They will be enslaved and oppressed for 400 years.\gva{d}\marginnote{\gva{d}\,Counting from the birth of Isaac to their departure of Egypt: Which declares that God will allow his to be afflicted in this world.}
\markboth{Genesis 15:14}{Genesis 15:14}\vs{14} But I will execute judgment on the nation that they will serve. Afterward they will come out with many possessions.
\markboth{Genesis 15:15}{Genesis 15:15}\vs{15} But as for you, you will go to your ancestors in peace and be buried at a good old age.
\markboth{Genesis 15:16}{Genesis 15:16}\vs{16} In the fourth generation your descendants will return here, for the sin of the Amorites has not yet reached its limit.''\gva{e}\marginnote{\gva{e}\,Though God tolerates the wicked for a time, yet his vengeance falls on them when the measure of their wickedness is full.}
\everypar{}

\parshape=0
\markboth{Genesis 15:17}{Genesis 15:17}\vs{17} When the sun had gone down and it was dark, a smoking firepot with a flaming torch passed between the animal parts.
\markboth{Genesis 15:18}{Genesis 15:18}\vs{18} That day the Lord made a covenant with Abram: ``To your descendants I give this land, from the river of Egypt to the great river, the Euphrates River---
\markboth{Genesis 15:19}{Genesis 15:19}\vs{19} the land of the Kenites, Kenizzites, Kadmonites,
\markboth{Genesis 15:20}{Genesis 15:20}\vs{20} Hittites, Perizzites, Rephaites,
\markboth{Genesis 15:21}{Genesis 15:21}\vs{21} Amorites, Canaanites, Girgashites, and Jebusites.''

\markboth{Genesis 16:1}{Genesis 16:1}
\Needspace*{8\baselineskip}
\ch{16} \allowchapbreak\hypertarget{ch-genesis-16}{}\lettrine[lines=5]{\color{pentateuch}N}{\textsc{ow}} Sarai, Abram's wife, had not given birth to any children, but she had an Egyptian servant named Hagar.\gva{a}\marginnote{\gva{a}\,It seems that she had respect for God's promise, which could not be accomplished without issue.}\everypar{}
\markboth{Genesis 16:2}{Genesis 16:2}\vs{2} So Sarai said to Abram, ``Since the Lord has prevented me from having children, please sleep with my servant. Perhaps I can have a family by her.'' Abram did what Sarai told him.\gva{b}\marginnote{\gva{b}\,She fails by limiting God's power to the common order of nature, as though God could not give her children in her old age.}
\everypar{}

\parshape=0
\markboth{Genesis 16:3}{Genesis 16:3}\vs{3} So after Abram had lived in Canaan for ten years, Sarai, Abram's wife, gave Hagar, her Egyptian servant, to her husband to be his wife.
\markboth{Genesis 16:4}{Genesis 16:4}\vs{4} He slept with Hagar, and she became pregnant. Once Hagar realized she was pregnant, she despised Sarai.\gva{c}\marginnote{\gva{c}\,This punishment declares what they gain if they attempt any thing against the word of God.}
\markboth{Genesis 16:5}{Genesis 16:5}\vs{5} Then Sarai said to Abram, ``You have brought this wrong on me! I gave my servant into your embrace, but when she realized that she was pregnant, she despised me. May the Lord judge between you and me!''
\everypar{}

\parshape=0
\markboth{Genesis 16:6}{Genesis 16:6}\vs{6} Abram said to Sarai, ``Since your servant is under your authority, do to her whatever you think best.'' Then Sarai treated Hagar harshly, so she ran away from Sarai.
\everypar{}

\parshape=0
\markboth{Genesis 16:7}{Genesis 16:7}\vs{7} The angel of the Lord found Hagar near a spring of water in the wilderness---the spring that is along the road to Shur.\gva{d}\marginnote{\gva{d}\,Which was Christ, as appears in Genesis 16:13 ; Genesis 18:17 .}
\markboth{Genesis 16:8}{Genesis 16:8}\vs{8} He said, ``Hagar, servant of Sarai, where have you come from, and where are you going?'' She replied, ``I'm running away from my mistress, Sarai.''
\everypar{}

\parshape=0
\markboth{Genesis 16:9}{Genesis 16:9}\vs{9} Then the angel of the Lord said to her, ``Return to your mistress and submit to her authority.\gva{e}\marginnote{\gva{e}\,God rejects no estate of people in their misery, but sends them comfort.}
\markboth{Genesis 16:10}{Genesis 16:10}\vs{10} I will greatly multiply your descendants,'' the angel of the Lord added, ``so that they will be too numerous to count.''
\everypar{}

\parshape=0
\markboth{Genesis 16:11}{Genesis 16:11}\vs{11} Then the angel of the Lord said to her, ``You are now pregnant and are about to give birth to a son. You are to name him Ishmael, for the Lord has heard your painful groans.
\everypar{}

\parshape=0
\markboth{Genesis 16:12}{Genesis 16:12}\vs{12} He will be a wild donkey of a man. He will be hostile to everyone, and everyone will be hostile to him. He will live away from his brothers.''\gva{f}\marginnote{\gva{f}\,That is, the Ishmaelites will be a separate people by themselves and not part of another people.}
\everypar{}

\parshape=0
\markboth{Genesis 16:13}{Genesis 16:13}\vs{13} So Hagar named the Lord who spoke to her, ``You are the God who sees me,'' for she said, ``Here I have seen one who sees me!''\gva{g}\marginnote{\gva{g}\,She rebukes her own dullness and acknowledges God's graces, who was present with her everywhere.}
\markboth{Genesis 16:14}{Genesis 16:14}\vs{14} That is why the well was called Beer Lahai Roi. (It is located between Kadesh and Bered.)
\everypar{}

\parshape=0
\markboth{Genesis 16:15}{Genesis 16:15}\vs{15} So Hagar gave birth to Abram's son, whom Abram named Ishmael.
\markboth{Genesis 16:16}{Genesis 16:16}\vs{16} (Now Abram was eighty-six years old when Hagar gave birth to Ishmael.)

\markboth{Genesis 17:1}{Genesis 17:1}
\Needspace*{8\baselineskip}
\ch{17} \allowchapbreak\hypertarget{ch-genesis-17}{}\lettrine[lines=5]{\color{pentateuch}W}{\textsc{hen}} Abram was ninety-nine years old, the Lord appeared to him and said, ``I am the Sovereign God. Walk before me and be blameless.\everypar{}
\markboth{Genesis 17:2}{Genesis 17:2}\vs{2} Then I will confirm my covenant between me and you, and I will give you a multitude of descendants.''
\\\indent\markboth{Genesis 17:3}{Genesis 17:3}\vs{3} Abram bowed down with his face to the ground, and God said to him,
\markboth{Genesis 17:4}{Genesis 17:4}\vs{4} ``As for me, this is my covenant with you: You will be the father of a multitude of nations.\gva{a}\marginnote{\gva{a}\,Not only physical descendants, but of a far greater multitude by faith, Romans 4:17 .}
\markboth{Genesis 17:5}{Genesis 17:5}\vs{5} No longer will your name be Abram. Instead, your name will be Abraham because I will make you the father of a multitude of nations.\gva{b}\marginnote{\gva{b}\,The changing of his name is a seal to confirm God's promise to him.}
\markboth{Genesis 17:6}{Genesis 17:6}\vs{6} I will make you extremely fruitful. I will make nations of you, and kings will descend from you.
\markboth{Genesis 17:7}{Genesis 17:7}\vs{7} I will confirm my covenant as a perpetual covenant between me and you. It will extend to your descendants after you throughout their generations. I will be your God and the God of your descendants after you.
\markboth{Genesis 17:8}{Genesis 17:8}\vs{8} I will give the whole land of Canaan---the land where you are now residing---to you and your descendants after you as a permanent possession. I will be their God.''
\everypar{}

\parshape=0
\markboth{Genesis 17:9}{Genesis 17:9}\vs{9} Then God said to Abraham, ``As for you, you must keep the covenantal requirement I am imposing on you and your descendants after you throughout their generations.
\markboth{Genesis 17:10}{Genesis 17:10}\vs{10} This is my requirement that you and your descendants after you must keep: Every male among you must be circumcised.\gva{c}\marginnote{\gva{c}\,Circumcision is called the covenant, because it signifies the covenant and has the promise of grace joined to it: a phrase that is common to all ordinances.}
\markboth{Genesis 17:11}{Genesis 17:11}\vs{11} You must circumcise the flesh of your foreskins. This will be a reminder of the covenant between me and you.\gva{d}\marginnote{\gva{d}\,That private part is circumcised, to show that all that is begotten by man is corrupt, and must die.}
\markboth{Genesis 17:12}{Genesis 17:12}\vs{12} Throughout your generations every male among you who is eight days old must be circumcised, whether born in your house or bought with money from any foreigner who is not one of your descendants.
\markboth{Genesis 17:13}{Genesis 17:13}\vs{13} They must indeed be circumcised, whether born in your house or bought with money. The sign of my covenant will be visible in your flesh as a permanent reminder.
\markboth{Genesis 17:14}{Genesis 17:14}\vs{14} Any uncircumcised male who has not been circumcised in the flesh of his foreskin will be cut off from his people---he has failed to carry out my requirement.''\gva{e}\marginnote{\gva{e}\,Though women were not circumcised, they still partook of God's promise: for under mankind all was consecrated. Here it is declared, that whoever condemns the sign, also despises the promise.}
\everypar{}

\parshape=0
\markboth{Genesis 17:15}{Genesis 17:15}\vs{15} Then God said to Abraham, ``As for your wife, you must no longer call her Sarai; Sarah will be her name.
\markboth{Genesis 17:16}{Genesis 17:16}\vs{16} I will bless her and will give you a son through her. I will bless her and she will become a mother of nations. Kings of countries will come from her!''
\everypar{}

\parshape=0
\markboth{Genesis 17:17}{Genesis 17:17}\vs{17} Then Abraham bowed down with his face to the ground and laughed as he said to himself, ``Can a son be born to a man who is a hundred years old? Can Sarah bear a child at the age of ninety?''\gva{f}\marginnote{\gva{f}\,Which proceeded from a sudden joy, and not from lack of faith.}
\markboth{Genesis 17:18}{Genesis 17:18}\vs{18} Abraham said to God, ``O that Ishmael might live before you!''
\everypar{}

\parshape=0
\markboth{Genesis 17:19}{Genesis 17:19}\vs{19} God said, ``No, Sarah your wife is going to bear you a son, and you will name him Isaac. I will confirm my covenant with him as a perpetual covenant for his descendants after him.\gva{g}\marginnote{\gva{g}\,The everlasting covenant is made with the children of the Spirit. A temporary promise is made with the children of the flesh, as was promised to Ishmael.}
\markboth{Genesis 17:20}{Genesis 17:20}\vs{20} As for Ishmael, I have heard you. I will indeed bless him, make him fruitful, and give him a multitude of descendants. He will become the father of twelve princes; I will make him into a great nation.
\markboth{Genesis 17:21}{Genesis 17:21}\vs{21} But I will establish my covenant with Isaac, whom Sarah will bear to you at this set time next year.''
\markboth{Genesis 17:22}{Genesis 17:22}\vs{22} When he finished speaking with Abraham, God went up from him.
\everypar{}

\parshape=0
\markboth{Genesis 17:23}{Genesis 17:23}\vs{23} Abraham took his son Ishmael and every male in his household (whether born in his house or bought with money) and circumcised them on that very same day, just as God had told him to do.\gva{h}\marginnote{\gva{h}\,They were well taught if they obeyed and were circumcised without resistance. This teaches that masters in their houses ought to be as preachers to their families, that from the highest to the lowest they may obey the will of God.}
\markboth{Genesis 17:24}{Genesis 17:24}\vs{24} Now Abraham was ninety-nine years old when he was circumcised;
\markboth{Genesis 17:25}{Genesis 17:25}\vs{25} his son Ishmael was thirteen years old when he was circumcised.
\markboth{Genesis 17:26}{Genesis 17:26}\vs{26} Abraham and his son Ishmael were circumcised on the very same day.
\markboth{Genesis 17:27}{Genesis 17:27}\vs{27} All the men of his household, whether born in his household or bought with money from a foreigner, were circumcised with him.

\markboth{Genesis 18:1}{Genesis 18:1}
\Needspace*{8\baselineskip}
\ch{18} \allowchapbreak\hypertarget{ch-genesis-18}{}\lettrine[lines=5]{\color{pentateuch}T}{\textsc{he}} Lord appeared to Abraham by the oaks of Mamre while he was sitting at the entrance to his tent during the hottest time of the day.\everypar{}
\markboth{Genesis 18:2}{Genesis 18:2}\vs{2} Abraham looked up and saw three men standing across from him. When he saw them he ran from the entrance of the tent to meet them and bowed low to the ground.\gva{a}\marginnote{\gva{a}\,That is, three angels in the shape of men.}
\\\indent\markboth{Genesis 18:3}{Genesis 18:3}\vs{3} He said, ``My lord, if I have found favor in your sight, do not pass by and leave your servant.\gva{b}\marginnote{\gva{b}\,Speaking to the one who appeared to be most majestic, for he thought they were men.}
\markboth{Genesis 18:4}{Genesis 18:4}\vs{4} Let a little water be brought so that you may all wash your feet and rest under the tree.\gva{c}\marginnote{\gva{c}\,For men used to go bare footed in those parts because of the heat.}
\markboth{Genesis 18:5}{Genesis 18:5}\vs{5} And let me get a bit of food so that you may refresh yourselves since you have passed by your servant's home. After that you may be on your way.'' ``All right,'' they replied, ``you may do as you say.''\gva{d}\marginnote{\gva{d}\,As sent by God that I should do my duty to you.}
\everypar{}

\parshape=0
\markboth{Genesis 18:6}{Genesis 18:6}\vs{6} So Abraham hurried into the tent and said to Sarah, ``Quick! Take three measures of fine flour, knead it, and make bread.''
\markboth{Genesis 18:7}{Genesis 18:7}\vs{7} Then Abraham ran to the herd and chose a fine, tender calf, and gave it to a servant, who quickly prepared it.
\markboth{Genesis 18:8}{Genesis 18:8}\vs{8} Abraham then took some curds and milk, along with the calf that had been prepared, and placed the food before them. They ate while he was standing near them under a tree.\gva{e}\marginnote{\gva{e}\,For as God gave them bodies for a time, so he gave them the abilities of them, to walk, to eat and drink, and such like.}
\everypar{}

\parshape=0
\markboth{Genesis 18:9}{Genesis 18:9}\vs{9} Then they asked him, ``Where is Sarah your wife?'' He replied, ``There, in the tent.''
\markboth{Genesis 18:10}{Genesis 18:10}\vs{10} One of them said, ``I will surely return to you when the season comes round again, and your wife Sarah will have a son!'' (Now Sarah was listening at the entrance to the tent, not far behind him.\gva{f}\marginnote{\gva{f}\,That is, about this time when she shall be alive, or when the child shall come into this life.}
\markboth{Genesis 18:11}{Genesis 18:11}\vs{11} Abraham and Sarah were old and advancing in years; Sarah had long since passed menopause.)
\markboth{Genesis 18:12}{Genesis 18:12}\vs{12} So Sarah laughed to herself, thinking, ``After I am worn out will I have pleasure, especially when my husband is old too?''\gva{g}\marginnote{\gva{g}\,For she believed the order of nature, rather than believing the promise of God.}
\everypar{}

\parshape=0
\markboth{Genesis 18:13}{Genesis 18:13}\vs{13} The Lord said to Abraham, ``Why did Sarah laugh and say, `Will I really have a child when I am old?'
\markboth{Genesis 18:14}{Genesis 18:14}\vs{14} Is anything impossible for the Lord? I will return to you when the season comes round again and Sarah will have a son.''
\markboth{Genesis 18:15}{Genesis 18:15}\vs{15} Then Sarah lied, saying, ``I did not laugh,'' because she was afraid. But the Lord said, ``No! You did laugh.''
\everypar{}

\parshape=0
\markboth{Genesis 18:16}{Genesis 18:16}\vs{16} When the men got up to leave, they looked out over Sodom. (Now Abraham was walking with them to see them on their way.)
\markboth{Genesis 18:17}{Genesis 18:17}\vs{17} Then the Lord said, ``Should I hide from Abraham what I am about to do?\gva{h}\marginnote{\gva{h}\,Jehovah the Hebrew word we call Lord, shows that this angel was Christ: for this word is only applied to God.}
\markboth{Genesis 18:18}{Genesis 18:18}\vs{18} After all, Abraham will surely become a great and powerful nation, and all the nations on the earth may receive blessing through him.
\markboth{Genesis 18:19}{Genesis 18:19}\vs{19} I have chosen him so that he may command his children and his household after him to keep the way of the Lord by doing what is right and just. Then the Lord will give to Abraham what he promised him.''\gva{i}\marginnote{\gva{i}\,He shows that fathers ought both to know God's judgments, and to declare them to their children.}
\everypar{}

\parshape=0
\markboth{Genesis 18:20}{Genesis 18:20}\vs{20} So the Lord said, ``The outcry against Sodom and Gomorrah is so great and their sin so blatant
\markboth{Genesis 18:21}{Genesis 18:21}\vs{21} that I must go down and see if they are as wicked as the outcry suggests. If not, I want to know.''\gva{k}\marginnote{\gva{k}\,God speaks after the fashion of men: that is, I will enter into judgment with good advise.}\gva{l}\marginnote{\gva{l}\,For our sins cry for vengeance, though no one accuses us.}
\everypar{}

\parshape=0
\markboth{Genesis 18:22}{Genesis 18:22}\vs{22} The two men turned and headed toward Sodom, but Abraham was still standing before the Lord.
\markboth{Genesis 18:23}{Genesis 18:23}\vs{23} Abraham approached and said, ``Will you really sweep away the godly along with the wicked?
\markboth{Genesis 18:24}{Genesis 18:24}\vs{24} What if there are 50 godly people in the city? Will you really wipe it out and not spare the place for the sake of the 50 godly people who are in it?
\markboth{Genesis 18:25}{Genesis 18:25}\vs{25} Far be it from you to do such a thing---to kill the godly with the wicked, treating the godly and the wicked alike! Far be it from you! Will not the judge of the whole earth do what is right?''
\everypar{}

\parshape=0
\markboth{Genesis 18:26}{Genesis 18:26}\vs{26} So the Lord replied, ``If I find in the city of Sodom 50 godly people, I will spare the whole place for their sake.''\gva{m}\marginnote{\gva{m}\,God declares that his judgments were done with great mercy, even though all were so corrupt that not only fifty but ten righteous men could not be found there, and also that the wicked are spared for the sake of the righteous.}
\everypar{}

\parshape=0
\markboth{Genesis 18:27}{Genesis 18:27}\vs{27} Then Abraham asked, ``Since I have undertaken to speak to the Lord (although I am but dust and ashes),\gva{n}\marginnote{\gva{n}\,By this we learn, that the nearer we approach to God, the more our miserable estate appears, and the more we are humbled.}
\markboth{Genesis 18:28}{Genesis 18:28}\vs{28} what if there are five less than the 50 godly people? Will you destroy the whole city because five are lacking?'' He replied, ``I will not destroy it if I find 45 there.''
\everypar{}

\parshape=0
\markboth{Genesis 18:29}{Genesis 18:29}\vs{29} Abraham spoke to him again, ``What if 40 are found there?'' He replied, ``I will not do it for the sake of the 40.''
\everypar{}

\parshape=0
\markboth{Genesis 18:30}{Genesis 18:30}\vs{30} Then Abraham said, ``May the Lord not be angry so that I may speak! What if 30 are found there?'' He replied, ``I will not do it if I find 30 there.''
\everypar{}

\parshape=0
\markboth{Genesis 18:31}{Genesis 18:31}\vs{31} Abraham said, ``Since I have undertaken to speak to the Lord, what if only 20 are found there?'' He replied, ``I will not destroy it for the sake of the 20.''
\everypar{}

\parshape=0
\markboth{Genesis 18:32}{Genesis 18:32}\vs{32} Finally Abraham said, ``May the Lord not be angry so that I may speak just once more. What if 10 are found there?'' He replied, ``I will not destroy it for the sake of the 10.''\gva{o}\marginnote{\gva{o}\,If God did not refuse the prayer for the wicked Sodomites, even to the sixth request, how much more will he grant the prayers of the godly for the afflicted Church?}
\everypar{}

\parshape=0
\markboth{Genesis 18:33}{Genesis 18:33}\vs{33} The Lord went on his way when he had finished speaking to Abraham. Then Abraham returned home.

\markboth{Genesis 19:1}{Genesis 19:1}
\Needspace*{8\baselineskip}
\ch{19} \allowchapbreak\hypertarget{ch-genesis-19}{}\lettrine[lines=5]{\color{pentateuch}T}{\textsc{he}} two angels came to Sodom in the evening while Lot was sitting in the city's gateway. When Lot saw them, he got up to meet them and bowed down with his face toward the ground.\gva{a}\marginnote{\gva{a}\,In which we see God's provident care in preserving his: even though he does not reveal himself to all alike: for Lot had but two angels, and Abraham three.}\everypar{}
\everypar{}

\parshape=0
\markboth{Genesis 19:2}{Genesis 19:2}\vs{2} He said, ``Here, my lords, please turn aside to your servant's house. Stay the night and wash your feet. Then you can be on your way early in the morning.'' ``No,'' they replied, ``we'll spend the night in the town square.''
\everypar{}

\parshape=0
\markboth{Genesis 19:3}{Genesis 19:3}\vs{3} But he urged them persistently, so they turned aside with him and entered his house. He prepared a feast for them, including bread baked without yeast, and they ate.\gva{b}\marginnote{\gva{b}\,That is, he begged them so insistently.}\gva{c}\marginnote{\gva{c}\,Not because they had need, but because the time was not yet come for them to reveal themselves.}
\markboth{Genesis 19:4}{Genesis 19:4}\vs{4} Before they could lie down to sleep, all the men---both young and old, from every part of the city of Sodom---surrounded the house.\gva{d}\marginnote{\gva{d}\,Nothing is more dangerous than to live where sin reigns: for it corrupts all.}
\markboth{Genesis 19:5}{Genesis 19:5}\vs{5} They shouted to Lot, ``Where are the men who came to you tonight? Bring them out to us so we can take carnal knowledge of them!''
\everypar{}

\parshape=0
\markboth{Genesis 19:6}{Genesis 19:6}\vs{6} Lot went outside to them, shutting the door behind him.
\markboth{Genesis 19:7}{Genesis 19:7}\vs{7} He said, ``No, my brothers! Don't act so wickedly!
\markboth{Genesis 19:8}{Genesis 19:8}\vs{8} Look, I have two daughters who have never been intimate with a man. Let me bring them out to you, and you can do to them whatever you please. Only don't do anything to these men, for they have come under the protection of my roof.''\gva{e}\marginnote{\gva{e}\,He deserves praise for defending his guests, but he is to be blamed for seeking unlawful means.}\gva{f}\marginnote{\gva{f}\,That I should preserve them from all injury.}
\everypar{}

\parshape=0
\markboth{Genesis 19:9}{Genesis 19:9}\vs{9} ``Out of our way!'' they cried, ``This man came to live here as a foreigner, and now he dares to judge us! We'll do more harm to you than to them!'' They kept pressing in on Lot until they were close enough to break down the door.
\everypar{}

\parshape=0
\markboth{Genesis 19:10}{Genesis 19:10}\vs{10} So the men inside reached out and pulled Lot back into the house as they shut the door.
\markboth{Genesis 19:11}{Genesis 19:11}\vs{11} Then they struck the men who were at the door of the house, from the youngest to the oldest, with blindness. The men outside wore themselves out trying to find the door.
\markboth{Genesis 19:12}{Genesis 19:12}\vs{12} Then the two visitors said to Lot, ``Who else do you have here? Do you have any sons-in-law, sons, daughters, or other relatives in the city? Get them out of this place
\markboth{Genesis 19:13}{Genesis 19:13}\vs{13} because we are about to destroy it. The outcry against this place is so great before the Lord that he has sent us to destroy it.''\gva{g}\marginnote{\gva{g}\,This proves that the angels are ministers, both to execute God's wrath and to declare his favour.}
\everypar{}

\parshape=0
\markboth{Genesis 19:14}{Genesis 19:14}\vs{14} Then Lot went out and spoke to his sons-in-law who were going to marry his daughters. He said, ``Quick, get out of this place because the Lord is about to destroy the city!'' But his sons-in-law thought he was ridiculing them.
\everypar{}

\parshape=0
\markboth{Genesis 19:15}{Genesis 19:15}\vs{15} At dawn the angels hurried Lot along, saying, ``Get going! Take your wife and your two daughters who are here, or else you will be destroyed when the city is judged!''
\markboth{Genesis 19:16}{Genesis 19:16}\vs{16} When Lot hesitated, the men grabbed his hand and the hands of his wife and two daughters because the Lord had compassion on them. They led them away and placed them outside the city.\gva{h}\marginnote{\gva{h}\,The mercy of God strives to overcome man's slowness in following God's calling.}
\markboth{Genesis 19:17}{Genesis 19:17}\vs{17} When they had brought them outside, they said, ``Run for your lives! Don't look behind you or stop anywhere in the valley! Escape to the mountains or you will be destroyed!''\gva{i}\marginnote{\gva{i}\,He willed him to flee God's judgments and not to be sorry to leave that rich country, full of vain pleasures.}
\everypar{}

\parshape=0
\markboth{Genesis 19:18}{Genesis 19:18}\vs{18} But Lot said to them, ``No, please, Lord!
\markboth{Genesis 19:19}{Genesis 19:19}\vs{19} Your servant has found favor with you, and you have shown me great kindness by sparing my life. But I am not able to escape to the mountains because this disaster will overtake me and I'll die.
\markboth{Genesis 19:20}{Genesis 19:20}\vs{20} Look, this town over here is close enough to escape to, and it's just a little one. Let me go there. It's just a little place, isn't it? Then I'll survive.''\gva{k}\marginnote{\gva{k}\,Though it is little, yet it is great enough to save my life: in which he errs by choosing another place than the angel had appointed him.}
\everypar{}

\parshape=0
\markboth{Genesis 19:21}{Genesis 19:21}\vs{21} ``Very well,'' he replied, ``I will grant this request too and will not overthrow the town you mentioned.
\markboth{Genesis 19:22}{Genesis 19:22}\vs{22} Run there quickly, for I cannot do anything until you arrive there.'' (This incident explains why the town was called Zoar.)\gva{l}\marginnote{\gva{l}\,Because God's commandment was to destroy the city and to save Lot.}\gva{m}\marginnote{\gva{m}\,Which before was called Belah, in Genesis 14:2 .}
\everypar{}

\parshape=0
\markboth{Genesis 19:23}{Genesis 19:23}\vs{23} The sun had just risen over the land as Lot reached Zoar.
\markboth{Genesis 19:24}{Genesis 19:24}\vs{24} Then the Lord rained down sulfur and fire on Sodom and Gomorrah. It was sent down from the sky by the Lord.
\markboth{Genesis 19:25}{Genesis 19:25}\vs{25} So he overthrew those cities and all that region, including all the inhabitants of the cities and the vegetation that grew from the ground.
\markboth{Genesis 19:26}{Genesis 19:26}\vs{26} But Lot's wife looked back longingly and was turned into a pillar of salt.\gva{n}\marginnote{\gva{n}\,Concerning the body only: this was a notable monument of God's vengeance to all who passed that way.}
\everypar{}

\parshape=0
\markboth{Genesis 19:27}{Genesis 19:27}\vs{27} Abraham got up early in the morning and went to the place where he had stood before the Lord.
\markboth{Genesis 19:28}{Genesis 19:28}\vs{28} He looked out toward Sodom and Gomorrah and all the land of that region. As he did so, he saw the smoke rising up from the land like smoke from a furnace.
\everypar{}

\parshape=0
\markboth{Genesis 19:29}{Genesis 19:29}\vs{29} So when God destroyed the cities of the region, God honored Abraham's request. He removed Lot from the midst of the destruction when he destroyed the cities Lot had lived in.
\everypar{}

\parshape=0
\markboth{Genesis 19:30}{Genesis 19:30}\vs{30} Lot went up from Zoar with his two daughters and settled in the mountains because he was afraid to live in Zoar. So he lived in a cave with his two daughters.\gva{o}\marginnote{\gva{o}\,Having felt God's mercy, he did not dare provoke him again by continuing among the wicked.}
\markboth{Genesis 19:31}{Genesis 19:31}\vs{31} Later the older daughter said to the younger, ``Our father is old, and there is no man in the country to sleep with us, the way everyone does.\gva{p}\marginnote{\gva{p}\,Meaning in the country which the Lord had now destroyed.}
\markboth{Genesis 19:32}{Genesis 19:32}\vs{32} Come, let's make our father drunk with wine so we can go to bed with him and preserve our family line through our father.''\gva{q}\marginnote{\gva{q}\,For unless he had been drunk, he would never have done that abominable act.}
\everypar{}

\parshape=0
\markboth{Genesis 19:33}{Genesis 19:33}\vs{33} So that night they made their father drunk with wine, and the older daughter came in and went to bed with her father. But he was not aware of when she lay down with him or when she got up.
\markboth{Genesis 19:34}{Genesis 19:34}\vs{34} So in the morning the older daughter said to the younger, ``Since I went to bed with my father last night, let's make him drunk again tonight. Then you go in and go to bed with him so we can preserve our family line through our father.''
\markboth{Genesis 19:35}{Genesis 19:35}\vs{35} So they made their father drunk that night as well, and the younger one came and went to bed with him. But he was not aware of when she lay down with him or when she got up.
\everypar{}

\parshape=0
\markboth{Genesis 19:36}{Genesis 19:36}\vs{36} In this way both of Lot's daughters became pregnant by their father.\gva{r}\marginnote{\gva{r}\,Thus God permitted him to fall most horribly in the solitary mountains, whom the wickedness of Sodom could not overcome.}
\markboth{Genesis 19:37}{Genesis 19:37}\vs{37} The older daughter gave birth to a son and named him Moab. He is the ancestor of the Moabites of today.\gva{s}\marginnote{\gva{s}\,Who as they were born in most horrible incest, so were they and their posterity vile and wicked.}
\markboth{Genesis 19:38}{Genesis 19:38}\vs{38} The younger daughter also gave birth to a son and named him Ben Ammi. He is the ancestor of the Ammonites of today.\gva{t}\marginnote{\gva{t}\,That is, son of my people: signifying that they rejoiced in their sin, rather than repenting of it.}

\markboth{Genesis 20:1}{Genesis 20:1}
\Needspace*{8\baselineskip}
\ch{20} \allowchapbreak\hypertarget{ch-genesis-20}{}\lettrine[lines=5]{\color{pentateuch}A}{\textsc{braham}} journeyed from there to the Negev region and settled between Kadesh and Shur. While he lived as a temporary resident in Gerar,\gva{a}\marginnote{\gva{a}\,Which was toward Egypt.}\everypar{}
\markboth{Genesis 20:2}{Genesis 20:2}\vs{2} Abraham said about his wife Sarah, ``She is my sister.'' So Abimelech, king of Gerar, sent for Sarah and took her.\gva{b}\marginnote{\gva{b}\,Abraham had now twice fallen into this sin: such is man's frailty.}
\everypar{}

\parshape=0
\markboth{Genesis 20:3}{Genesis 20:3}\vs{3} But God appeared to Abimelech in a dream at night and said to him, ``You are as good as dead because of the woman you have taken, for she is someone else's wife.''\gva{c}\marginnote{\gva{c}\,So greatly God detests the breach of marriage.}
\everypar{}

\parshape=0
\markboth{Genesis 20:4}{Genesis 20:4}\vs{4} Now Abimelech had not gone near her. He said, ``Lord, would you really slaughter an innocent nation?\gva{d}\marginnote{\gva{d}\,The infidels confessed that God would not punish but for just occasion: therefore, when he punishes, the occasion is just.}
\markboth{Genesis 20:5}{Genesis 20:5}\vs{5} Did Abraham not say to me, `She is my sister'? And she herself said, `He is my brother.' I have done this with a clear conscience and with innocent hands!''\gva{e}\marginnote{\gva{e}\,As one falling by ignorance, and not doing evil on purpose.}\gva{f}\marginnote{\gva{f}\,Not thinking to do any man harm.}
\everypar{}

\parshape=0
\markboth{Genesis 20:6}{Genesis 20:6}\vs{6} Then in the dream God replied to him, ``Yes, I know that you have done this with a clear conscience. That is why I have kept you from sinning against me and why I did not allow you to touch her.\gva{g}\marginnote{\gva{g}\,God by his holy Spirit restrains those who offend in ignorance, that they not fall into greater offence..}
\markboth{Genesis 20:7}{Genesis 20:7}\vs{7} But now give back the man's wife. Indeed he is a prophet and he will pray for you; thus you will live. But if you don't give her back, know that you will surely die along with all who belong to you.''\gva{h}\marginnote{\gva{h}\,That is, one to whom God reveals himself familiarly.}\gva{i}\marginnote{\gva{i}\,For the prayer of the godly is of force towards God.}
\everypar{}

\parshape=0
\markboth{Genesis 20:8}{Genesis 20:8}\vs{8} Early in the morning Abimelech summoned all his servants. When he told them about all these things, they were terrified.
\markboth{Genesis 20:9}{Genesis 20:9}\vs{9} Abimelech summoned Abraham and said to him, ``What have you done to us? What sin did I commit against you that would cause you to bring such great guilt on me and my kingdom? You have done things to me that should not be done!''\gva{k}\marginnote{\gva{k}\,The wickedness of the king brings God's wrath on the whole realm.}
\markboth{Genesis 20:10}{Genesis 20:10}\vs{10} Then Abimelech asked Abraham, ``What prompted you to do this thing?''
\everypar{}

\parshape=0
\markboth{Genesis 20:11}{Genesis 20:11}\vs{11} Abraham replied, ``Because I thought, `Surely no one fears God in this place. They will kill me because of my wife.'\gva{l}\marginnote{\gva{l}\,He shows that no honesty can be hoped for, where there is no fear of God.}
\markboth{Genesis 20:12}{Genesis 20:12}\vs{12} What's more, she is indeed my sister, my father's daughter, but not my mother's daughter. She became my wife.\gva{m}\marginnote{\gva{m}\,By sister, he means his full cousin, and by daughter Abraham's niece, Genesis 11:29 for so the Hebrews use these words.}
\markboth{Genesis 20:13}{Genesis 20:13}\vs{13} When God made me wander from my father's house, I told her, `This is what you can do to show your loyalty to me: Every place we go, say about me, ``He is my brother.'''''
\everypar{}

\parshape=0
\markboth{Genesis 20:14}{Genesis 20:14}\vs{14} So Abimelech gave sheep, cattle, and male and female servants to Abraham. He also gave his wife Sarah back to him.
\markboth{Genesis 20:15}{Genesis 20:15}\vs{15} Then Abimelech said, ``Look, my land is before you; live wherever you please.''
\everypar{}

\parshape=0
\markboth{Genesis 20:16}{Genesis 20:16}\vs{16} To Sarah he said, ``Look, I have given 1,000 pieces of silver to your `brother.' This is compensation for you so that you will stand vindicated before all who are with you.''\gva{n}\marginnote{\gva{n}\,Such a head as with whom you may be preserved from all dangers.}\gva{o}\marginnote{\gva{o}\,God caused this heathen king to reprove her because she concealed her identity, seeing that God had given her a husband as her veil and defence.}
\everypar{}

\parshape=0
\markboth{Genesis 20:17}{Genesis 20:17}\vs{17} Abraham prayed to God, and God healed Abimelech, as well as his wife and female slaves so that they were able to have children.
\markboth{Genesis 20:18}{Genesis 20:18}\vs{18} For the Lord had caused infertility to strike every woman in the household of Abimelech because he took Sarah, Abraham's wife.\gva{p}\marginnote{\gva{p}\,Had taken away from them the gift of conceiving.}

\markboth{Genesis 21:1}{Genesis 21:1}
\Needspace*{8\baselineskip}
\ch{21} \allowchapbreak\hypertarget{ch-genesis-21}{}\lettrine[lines=5]{\color{pentateuch}T}{\textsc{he}} Lord visited Sarah just as he had said he would and did for Sarah what he had promised.\everypar{}
\markboth{Genesis 21:2}{Genesis 21:2}\vs{2} So Sarah became pregnant and bore Abraham a son in his old age at the appointed time that God had told him.\gva{a}\marginnote{\gva{a}\,Therefore the miracle was greater.}
\markboth{Genesis 21:3}{Genesis 21:3}\vs{3} Abraham named his son---whom Sarah bore to him---Isaac.
\markboth{Genesis 21:4}{Genesis 21:4}\vs{4} When his son Isaac was eight days old, Abraham circumcised him just as God had commanded him to do.
\markboth{Genesis 21:5}{Genesis 21:5}\vs{5} (Now Abraham was one hundred years old when his son Isaac was born to him.)
\everypar{}

\parshape=0
\markboth{Genesis 21:6}{Genesis 21:6}\vs{6} Sarah said, ``God has made me laugh. Everyone who hears about this will laugh with me.''
\markboth{Genesis 21:7}{Genesis 21:7}\vs{7} She went on to say, ``Who would have said to Abraham that Sarah would nurse children? Yet I have given birth to a son for him in his old age!''\gva{b}\marginnote{\gva{b}\,She accuses herself of ingratitude, that she did not believe the angel.}
\everypar{}

\parshape=0
\markboth{Genesis 21:8}{Genesis 21:8}\vs{8} The child grew and was weaned. Abraham prepared a great feast on the day that Isaac was weaned.
\markboth{Genesis 21:9}{Genesis 21:9}\vs{9} But Sarah noticed the son of Hagar the Egyptian---the son whom Hagar had borne to Abraham---mocking.\gva{c}\marginnote{\gva{c}\,He derided God's promise made to Isaac which the apostle calls persecution Galatians 4:29 .}
\markboth{Genesis 21:10}{Genesis 21:10}\vs{10} So she said to Abraham, ``Banish that slave woman and her son, for the son of that slave woman will not be an heir along with my son Isaac!''
\everypar{}

\parshape=0
\markboth{Genesis 21:11}{Genesis 21:11}\vs{11} Sarah's demand displeased Abraham greatly because Ishmael was his son.
\markboth{Genesis 21:12}{Genesis 21:12}\vs{12} But God said to Abraham, ``Do not be upset about the boy or your slave wife. Do all that Sarah is telling you because through Isaac your descendants will be counted.\gva{d}\marginnote{\gva{d}\,The promised seed will be from Isaac, and not from Ishmael, Romans 9:7 , Hebrews 11:18 .}
\markboth{Genesis 21:13}{Genesis 21:13}\vs{13} But I will also make the son of the slave wife into a great nation, for he is your descendant too.''\gva{e}\marginnote{\gva{e}\,The Ishmaelites will come from him.}
\everypar{}

\parshape=0
\markboth{Genesis 21:14}{Genesis 21:14}\vs{14} Early in the morning Abraham took some food and a skin of water and gave them to Hagar. He put them on her shoulders, gave her the child, and sent her away. So she went wandering aimlessly through the wilderness of Beer Sheba.\gva{f}\marginnote{\gva{f}\,True faith renounces all natural affections to obey God's commandment.}
\markboth{Genesis 21:15}{Genesis 21:15}\vs{15} When the water in the skin was gone, she shoved the child under one of the shrubs.
\markboth{Genesis 21:16}{Genesis 21:16}\vs{16} Then she went and sat down by herself across from him at quite a distance, about a bowshot, away; for she thought, ``I refuse to watch the child die.'' So she sat across from him and wept uncontrollably.
\everypar{}

\parshape=0
\markboth{Genesis 21:17}{Genesis 21:17}\vs{17} But God heard the boy's voice. The angel of God called to Hagar from heaven and asked her, ``What is the matter, Hagar? Don't be afraid, for God has heard the boy's voice right where he is crying.\gva{g}\marginnote{\gva{g}\,For his promise sake made to Abraham; and not because the child had discretion and judgment to pray.}
\markboth{Genesis 21:18}{Genesis 21:18}\vs{18} Get up! Help the boy up and hold him by the hand, for I will make him into a great nation.''
\markboth{Genesis 21:19}{Genesis 21:19}\vs{19} Then God enabled Hagar to see a well of water. She went over and filled the skin with water, and then gave the boy a drink.\gva{h}\marginnote{\gva{h}\,Unless God opens our eyes, we can neither see, nor use the means which are before us.}
\everypar{}

\parshape=0
\markboth{Genesis 21:20}{Genesis 21:20}\vs{20} God was with the boy as he grew. He lived in the wilderness and became an archer.\gva{i}\marginnote{\gva{i}\,Concerning outward things God caused him to prosper.}
\markboth{Genesis 21:21}{Genesis 21:21}\vs{21} He lived in the wilderness of Paran. His mother found a wife for him from the land of Egypt.
\everypar{}

\parshape=0
\markboth{Genesis 21:22}{Genesis 21:22}\vs{22} At that time Abimelech and Phicol, the commander of his army, said to Abraham, ``God is with you in all that you do.
\markboth{Genesis 21:23}{Genesis 21:23}\vs{23} Now swear to me right here in God's name that you will not deceive me, my children, or my descendants. Show me, and the land where you are staying, the same loyalty that I have shown you.''
\everypar{}

\parshape=0
\markboth{Genesis 21:24}{Genesis 21:24}\vs{24} Abraham said, ``I swear to do this.''\gva{k}\marginnote{\gva{k}\,So that it is a lawful thing to take an oath in matters of importance, to justify the truth, and to assure others of our sincerity.}
\markboth{Genesis 21:25}{Genesis 21:25}\vs{25} But Abraham lodged a complaint against Abimelech concerning a well that Abimelech's servants had seized.
\markboth{Genesis 21:26}{Genesis 21:26}\vs{26} ``I do not know who has done this thing,'' Abimelech replied. ``Moreover, you did not tell me. I did not hear about it until today.''\gva{l}\marginnote{\gva{l}\,Wicked servants do many evils unknown to their masters.}
\everypar{}

\parshape=0
\markboth{Genesis 21:27}{Genesis 21:27}\vs{27} Abraham took some sheep and cattle and gave them to Abimelech. The two of them made a treaty.
\markboth{Genesis 21:28}{Genesis 21:28}\vs{28} Then Abraham set seven ewe lambs apart from the flock by themselves.
\markboth{Genesis 21:29}{Genesis 21:29}\vs{29} Abimelech asked Abraham, ``What is the meaning of these seven ewe lambs that you have set apart?''
\markboth{Genesis 21:30}{Genesis 21:30}\vs{30} He replied, ``You must take these seven ewe lambs from my hand as legal proof that I dug this well.''
\markboth{Genesis 21:31}{Genesis 21:31}\vs{31} That is why he named that place Beer Sheba, because the two of them swore an oath there.
\everypar{}

\parshape=0
\markboth{Genesis 21:32}{Genesis 21:32}\vs{32} So they made a treaty at Beer Sheba; then Abimelech and Phicol, the commander of his army, returned to the land of the Philistines.\gva{m}\marginnote{\gva{m}\,Thus we see that the godly, concerning outward things may make peace with the wicked that do not know the true God.}
\markboth{Genesis 21:33}{Genesis 21:33}\vs{33} Abraham planted a tamarisk tree in Beer Sheba. There he worshiped the Lord, the eternal God.\gva{n}\marginnote{\gva{n}\,That is, he worshipped God in all points of true religion.}
\markboth{Genesis 21:34}{Genesis 21:34}\vs{34} So Abraham stayed in the land of the Philistines for quite some time.

\markboth{Genesis 22:1}{Genesis 22:1}
\Needspace*{8\baselineskip}
\ch{22} \allowchapbreak\hypertarget{ch-genesis-22}{}\lettrine[lines=5]{\color{pentateuch}S}{\textsc{ome}} time after these things God tested Abraham. He said to him, ``Abraham!'' ``Here I am!'' Abraham replied.\everypar{}
\markboth{Genesis 22:2}{Genesis 22:2}\vs{2} God said, ``Take your son---your only son, whom you love, Isaac---and go to the land of Moriah! Offer him up there as a burnt offering on one of the mountains which I will indicate to you.''\gva{a}\marginnote{\gva{a}\,Signifying the fear of God, in which place he was also honoured, Solomon later building the temple there.}\gva{b}\marginnote{\gva{b}\,This was the main point of his temptation, seeing that he was commanded to offer up him in whom God had promised to bless all the nations of the world.}
\everypar{}

\parshape=0
\markboth{Genesis 22:3}{Genesis 22:3}\vs{3} Early in the morning Abraham got up and saddled his donkey. He took two of his young servants with him, along with his son Isaac. When he had cut the wood for the burnt offering, he started out for the place God had spoken to him about.
\everypar{}

\parshape=0
\markboth{Genesis 22:4}{Genesis 22:4}\vs{4} On the third day Abraham caught sight of the place in the distance.
\markboth{Genesis 22:5}{Genesis 22:5}\vs{5} So he said to his servants, ``You two stay here with the donkey while the boy and I go up there. We will worship and then return to you.''\gva{e}\marginnote{\gva{e}\,He did not doubt that God would accomplish his promise, even if he should sacrifice his son.}
\everypar{}

\parshape=0
\markboth{Genesis 22:6}{Genesis 22:6}\vs{6} Abraham took the wood for the burnt offering and put it on his son Isaac. Then he took the fire and the knife in his hand, and the two of them walked on together.
\markboth{Genesis 22:7}{Genesis 22:7}\vs{7} Isaac said to his father Abraham, ``My father?'' ``What is it, my son?'' he replied. ``Here is the fire and the wood,'' Isaac said, ``but where is the lamb for the burnt offering?''
\markboth{Genesis 22:8}{Genesis 22:8}\vs{8} ``God will provide for himself the lamb for the burnt offering, my son,'' Abraham replied. The two of them continued on together.\gva{d}\marginnote{\gva{d}\,The only way to overcome all temptation is to rest on God's providence.}
\everypar{}

\parshape=0
\markboth{Genesis 22:9}{Genesis 22:9}\vs{9} When they came to the place God had told him about, Abraham built the altar there and arranged the wood on it. Next he tied up his son Isaac and placed him on the altar on top of the wood.\gva{e}\marginnote{\gva{e}\,For it is likely that his father had told him God's commandment, to which he showed himself obedient.}
\markboth{Genesis 22:10}{Genesis 22:10}\vs{10} Then Abraham reached out his hand, took the knife, and prepared to slaughter his son.
\markboth{Genesis 22:11}{Genesis 22:11}\vs{11} But the angel of the Lord called to him from heaven, ``Abraham! Abraham!'' ``Here I am!'' he answered.
\markboth{Genesis 22:12}{Genesis 22:12}\vs{12} ``Do not harm the boy!'' the angel said. ``Do not do anything to him, for now I know that you fear God because you did not withhold your son, your only son, from me.''\gva{f}\marginnote{\gva{f}\,That is, by your true obedience you have declared your living faith.}
\everypar{}

\parshape=0
\markboth{Genesis 22:13}{Genesis 22:13}\vs{13} Abraham looked up and saw behind him a ram caught in the bushes by its horns. So he went over and got the ram and offered it up as a burnt offering instead of his son.
\markboth{Genesis 22:14}{Genesis 22:14}\vs{14} And Abraham called the name of that place ``The Lord provides.'' It is said to this day, ``In the mountain of the Lord provision will be made.''\gva{g}\marginnote{\gva{g}\,The name is changed to show that God both sees and provides secretly for his and also evidently is seen, and felt in the right time.}
\everypar{}

\parshape=0
\markboth{Genesis 22:15}{Genesis 22:15}\vs{15} The angel of the Lord called to Abraham a second time from heaven
\markboth{Genesis 22:16}{Genesis 22:16}\vs{16} and said, ``I solemnly swear by my own name, decrees the Lord, that because you have done this and have not withheld your son, your only son,\gva{h}\marginnote{\gva{h}\,Signifying, that there is none greater then he.}
\markboth{Genesis 22:17}{Genesis 22:17}\vs{17} I will indeed bless you, and I will greatly multiply your descendants so that they will be as countless as the stars in the sky or the grains of sand on the seashore. Your descendants will take possession of the strongholds of their enemies.
\markboth{Genesis 22:18}{Genesis 22:18}\vs{18} Because you have obeyed me, all the nations of the earth will pronounce blessings on one another using the name of your descendants.''
\everypar{}

\parshape=0
\markboth{Genesis 22:19}{Genesis 22:19}\vs{19} Then Abraham returned to his servants, and they set out together for Beer Sheba where Abraham stayed.
\everypar{}

\parshape=0
\markboth{Genesis 22:20}{Genesis 22:20}\vs{20} After these things Abraham was told, ``Milcah also has borne children to your brother Nahor---
\markboth{Genesis 22:21}{Genesis 22:21}\vs{21} Uz the firstborn, his brother Buz, Kemuel (the father of Aram),
\markboth{Genesis 22:22}{Genesis 22:22}\vs{22} Kesed, Hazo, Pildash, Jidlaph, and Bethuel.''
\markboth{Genesis 22:23}{Genesis 22:23}\vs{23} (Now Bethuel became the father of Rebekah.) These were the eight sons Milcah bore to Abraham's brother Nahor.
\markboth{Genesis 22:24}{Genesis 22:24}\vs{24} His concubine, whose name was Reumah, also bore him children---Tebah, Gaham, Tahash, and Maacah.\gva{i}\marginnote{\gva{i}\,Concubine is often used to refer to those women who were inferior to the wives.}

\markboth{Genesis 23:1}{Genesis 23:1}
\Needspace*{8\baselineskip}
\ch{23} \allowchapbreak\hypertarget{ch-genesis-23}{}\lettrine[lines=5]{\color{pentateuch}S}{\textsc{arah}} lived 127 years.\everypar{}
\markboth{Genesis 23:2}{Genesis 23:2}\vs{2} Then she died in Kiriath Arba (that is, Hebron) in the land of Canaan. Abraham went to mourn for Sarah and to weep for her.
\\\indent\markboth{Genesis 23:3}{Genesis 23:3}\vs{3} Then Abraham got up from mourning his dead wife and said to the sons of Heth,\gva{a}\marginnote{\gva{a}\,That is, when he had mourned: so the godly may mourn if they do not pass measure, and the natural affection is commendable.}
\markboth{Genesis 23:4}{Genesis 23:4}\vs{4} ``I am a foreign resident, a temporary settler, among you. Grant me ownership of a burial site among you so that I may bury my dead.''
\\\indent\markboth{Genesis 23:5}{Genesis 23:5}\vs{5} The sons of Heth answered Abraham,
\markboth{Genesis 23:6}{Genesis 23:6}\vs{6} ``Listen, sir, you are a mighty prince among us! You may bury your dead in the choicest of our tombs. None of us will refuse you his tomb to prevent you from burying your dead.''
\everypar{}

\parshape=0
\markboth{Genesis 23:7}{Genesis 23:7}\vs{7} Abraham got up and bowed down to the local people, the sons of Heth.
\markboth{Genesis 23:8}{Genesis 23:8}\vs{8} Then he said to them, ``If you agree that I may bury my dead, then hear me out. Ask Ephron the son of Zohar
\markboth{Genesis 23:9}{Genesis 23:9}\vs{9} if he will sell me the cave of Machpelah that belongs to him; it is at the end of his field. Let him sell it to me publicly for the full price, so that I may own it as a burial site.''
\everypar{}

\parshape=0
\markboth{Genesis 23:10}{Genesis 23:10}\vs{10} (Now Ephron was sitting among the sons of Heth.) Ephron the Hittite replied to Abraham in the hearing of the sons of Heth---before all who entered the gate of his city---\gva{c}\marginnote{\gva{c}\,Meaning all the citizens and inhabitants.}
\markboth{Genesis 23:11}{Genesis 23:11}\vs{11} ``No, my lord! Hear me out. I sell you both the field and the cave that is in it. In the presence of my people I sell it to you. Bury your dead.''
\everypar{}

\parshape=0
\markboth{Genesis 23:12}{Genesis 23:12}\vs{12} Abraham bowed before the local people\gva{d}\marginnote{\gva{d}\,To show that he had them in good estimation and reverence.}
\markboth{Genesis 23:13}{Genesis 23:13}\vs{13} and said to Ephron in their hearing, ``Hear me, if you will. I pay to you the price of the field. Take it from me so that I may bury my dead there.''
\everypar{}

\parshape=0
\markboth{Genesis 23:14}{Genesis 23:14}\vs{14} Ephron answered Abraham, saying to him,
\markboth{Genesis 23:15}{Genesis 23:15}\vs{15} ``Hear me, my lord. The land is worth 400 pieces of silver, but what is that between me and you? So bury your dead.''\gva{e}\marginnote{\gva{e}\,The common shekel is about 20 pence, so then 400 shekels is equal to 33 pounds, 6 shillings and 8 pence at 5 shilling sterling to the ounce.}
\everypar{}

\parshape=0
\markboth{Genesis 23:16}{Genesis 23:16}\vs{16} So Abraham agreed to Ephron's price and weighed out for him the price that Ephron had quoted in the hearing of the sons of Heth---400 pieces of silver, according to the standard measurement at the time.
\everypar{}

\parshape=0
\markboth{Genesis 23:17}{Genesis 23:17}\vs{17} So Abraham secured Ephron's field in Machpelah, next to Mamre, including the field, the cave that was in it, and all the trees that were in the field and all around its border,
\markboth{Genesis 23:18}{Genesis 23:18}\vs{18} as his property in the presence of the sons of Heth before all who entered the gate of Ephron's city.
\everypar{}

\parshape=0
\markboth{Genesis 23:19}{Genesis 23:19}\vs{19} After this Abraham buried his wife Sarah in the cave in the field of Machpelah next to Mamre (that is, Hebron) in the land of Canaan.
\markboth{Genesis 23:20}{Genesis 23:20}\vs{20} So Abraham secured the field and the cave that was in it as a burial site from the sons of Heth.\gva{f}\marginnote{\gva{f}\,That is, all the people confirmed the sale.}

\markboth{Genesis 24:1}{Genesis 24:1}
\Needspace*{8\baselineskip}
\ch{24} \allowchapbreak\hypertarget{ch-genesis-24}{}\lettrine[lines=5]{\color{pentateuch}N}{\textsc{ow}} Abraham was old, well advanced in years, and the Lord had blessed him in everything.\everypar{}
\markboth{Genesis 24:2}{Genesis 24:2}\vs{2} Abraham said to his servant, the senior one in his household who was in charge of everything he had, ``Put your hand under my thigh\gva{a}\marginnote{\gva{a}\,Which ceremony declared, the servants obedience towards his master, and the master's power over the servant.}
\markboth{Genesis 24:3}{Genesis 24:3}\vs{3} so that I may make you solemnly promise by the Lord, the God of heaven and the God of the earth: You must not acquire a wife for my son from the daughters of the Canaanites, among whom I am living.\gva{b}\marginnote{\gva{b}\,This shows that an oath may be required in a lawful cause.}
\markboth{Genesis 24:4}{Genesis 24:4}\vs{4} You must go instead to my country and to my relatives to find a wife for my son Isaac.''\gva{c}\marginnote{\gva{c}\,He did not want his son to marry out of the godly family: for the problems that come from marrying the ungodly are set forth in various places throughout the scriptures.}
\everypar{}

\parshape=0
\markboth{Genesis 24:5}{Genesis 24:5}\vs{5} The servant asked him, ``What if the woman is not willing to come back with me to this land? Must I then take your son back to the land from which you came?''
\everypar{}

\parshape=0
\markboth{Genesis 24:6}{Genesis 24:6}\vs{6} ``Be careful never to take my son back there!'' Abraham told him.\gva{d}\marginnote{\gva{d}\,Lest he should love the inheritance promised.}
\markboth{Genesis 24:7}{Genesis 24:7}\vs{7} ``The Lord, the God of heaven, who took me from my father's house and the land of my relatives, promised me with a solemn oath, `To your descendants I will give this land.' He will send his angel before you so that you may find a wife for my son from there.
\markboth{Genesis 24:8}{Genesis 24:8}\vs{8} But if the woman is not willing to come back with you, you will be free from this oath of mine. But you must not take my son back there!''
\markboth{Genesis 24:9}{Genesis 24:9}\vs{9} So the servant placed his hand under the thigh of his master Abraham and gave his solemn promise he would carry out his wishes.
\everypar{}

\parshape=0
\markboth{Genesis 24:10}{Genesis 24:10}\vs{10} Then the servant took ten of his master's camels and departed with all kinds of gifts from his master at his disposal. He journeyed to the region of Aram Naharaim and the city of Nahor.\gva{e}\marginnote{\gva{e}\,That is, to Charran.}
\markboth{Genesis 24:11}{Genesis 24:11}\vs{11} He made the camels kneel down by the well outside the city. It was evening, the time when the women would go out to draw water.
\markboth{Genesis 24:12}{Genesis 24:12}\vs{12} He prayed, ``O Lord, God of my master Abraham, guide me today. Be faithful to my master Abraham.\gva{f}\marginnote{\gva{f}\,He grounds his prayer on God's promise made to his master.}
\markboth{Genesis 24:13}{Genesis 24:13}\vs{13} Here I am, standing by the spring, and the daughters of the people who live in the town are coming out to draw water.
\markboth{Genesis 24:14}{Genesis 24:14}\vs{14} I will say to a young woman, `Please lower your jar so I may drink.' May the one you have chosen for your servant Isaac reply, `Drink, and I'll give your camels water too.' In this way I will know that you have been faithful to my master.''\gva{g}\marginnote{\gva{g}\,The servant moved by God's spirit sought assurance by a sign, as to whether or not God would prosper his journey.}
\everypar{}

\parshape=0
\markboth{Genesis 24:15}{Genesis 24:15}\vs{15} Before he had finished praying, there came Rebekah with her water jug on her shoulder. She was the daughter of Bethuel son of Milcah (Milcah was the wife of Abraham's brother Nahor).\gva{h}\marginnote{\gva{h}\,God gives success to all things that are done for the glory of his name and according to his word.}
\markboth{Genesis 24:16}{Genesis 24:16}\vs{16} Now the young woman was very beautiful. She was a virgin; no man had ever been physically intimate with her. She went down to the spring, filled her jug, and came back up.\gva{i}\marginnote{\gva{i}\,Here is declared that God hears the prayers of his own, and grants their requests.}
\markboth{Genesis 24:17}{Genesis 24:17}\vs{17} Abraham's servant ran to meet her and said, ``Please give me a sip of water from your jug.''
\markboth{Genesis 24:18}{Genesis 24:18}\vs{18} ``Drink, my lord,'' she replied, and quickly lowering her jug to her hands, she gave him a drink.
\markboth{Genesis 24:19}{Genesis 24:19}\vs{19} When she had done so, she said, ``I'll draw water for your camels too, until they have drunk as much as they want.''
\markboth{Genesis 24:20}{Genesis 24:20}\vs{20} She quickly emptied her jug into the watering trough and ran back to the well to draw more water until she had drawn enough for all his camels.
\markboth{Genesis 24:21}{Genesis 24:21}\vs{21} Silently the man watched her with interest to determine if the Lord had made his journey successful or not.
\everypar{}

\parshape=0
\markboth{Genesis 24:22}{Genesis 24:22}\vs{22} After the camels had finished drinking, the man took out a gold nose ring weighing a beka and two gold wrist bracelets weighing ten shekels and gave them to her.\gva{k}\marginnote{\gva{k}\,God permitted many things both in apparel and other things which are now forbidden especially when they do not suit our humble estate.}\gva{l}\marginnote{\gva{l}\,The golden shekel is meant here, not silver.}
\markboth{Genesis 24:23}{Genesis 24:23}\vs{23} ``Whose daughter are you?'' he asked. ``Tell me, is there room in your father's house for us to spend the night?''
\everypar{}

\parshape=0
\markboth{Genesis 24:24}{Genesis 24:24}\vs{24} She said to him, ``I am the daughter of Bethuel the son of Milcah, whom Milcah bore to Nahor.
\markboth{Genesis 24:25}{Genesis 24:25}\vs{25} We have plenty of straw and feed,'' she added, ``and room for you to spend the night.''
\everypar{}

\parshape=0
\markboth{Genesis 24:26}{Genesis 24:26}\vs{26} The man bowed his head and worshiped the Lord,
\markboth{Genesis 24:27}{Genesis 24:27}\vs{27} saying, ``Praised be the Lord, the God of my master Abraham, who has not abandoned his faithful love for my master! The Lord has led me to the house of my master's relatives!''\gva{m}\marginnote{\gva{m}\,He does not boast in his good fortune}
\everypar{}

\parshape=0
\markboth{Genesis 24:28}{Genesis 24:28}\vs{28} The young woman ran and told her mother's household all about these things.
\markboth{Genesis 24:29}{Genesis 24:29}\vs{29} (Now Rebekah had a brother named Laban.) Laban rushed out to meet the man at the spring.
\markboth{Genesis 24:30}{Genesis 24:30}\vs{30} When he saw the bracelets on his sister's wrists and the nose ring and heard his sister Rebekah say, ``This is what the man said to me,'' he went out to meet the man. There he was, standing by the camels near the spring.\gva{n}\marginnote{\gva{n}\,For he waited on God's hand, who heard his prayer.}
\markboth{Genesis 24:31}{Genesis 24:31}\vs{31} Laban said to him, ``Come, you who are blessed by the Lord! Why are you standing out here when I have prepared the house and a place for the camels?''
\everypar{}

\parshape=0
\markboth{Genesis 24:32}{Genesis 24:32}\vs{32} So Abraham's servant went to the house and unloaded the camels. Straw and feed were given to the camels, and water was provided so that he and the men who were with him could wash their feet.\gva{o}\marginnote{\gva{o}\,That is, Laban.}\gva{p}\marginnote{\gva{p}\,The gentle entertainment of strangers practised by the godly fathers.}
\markboth{Genesis 24:33}{Genesis 24:33}\vs{33} When food was served, he said, ``I will not eat until I have said what I want to say.'' ``Tell us,'' Laban said.\gva{q}\marginnote{\gva{q}\,The faithfulness that servants owe to their masters, causes them to prefer their masters business before their own needs.}
\everypar{}

\parshape=0
\markboth{Genesis 24:34}{Genesis 24:34}\vs{34} ``I am the servant of Abraham,'' he began.
\markboth{Genesis 24:35}{Genesis 24:35}\vs{35} ``The Lord has richly blessed my master and he has become very wealthy. The Lord has given him sheep and cattle, silver and gold, male and female servants, and camels and donkeys.\gva{r}\marginnote{\gva{r}\,To bless signifies here to enrich, or increase with substance as the text in the same verse declares.}
\markboth{Genesis 24:36}{Genesis 24:36}\vs{36} My master's wife Sarah bore a son to him when she was old, and my master has given him everything he owns.
\markboth{Genesis 24:37}{Genesis 24:37}\vs{37} My master made me swear an oath. He said, `You must not acquire a wife for my son from the daughters of the Canaanites, among whom I am living,\gva{s}\marginnote{\gva{s}\,The Canaanites were cursed, and therefore the godly could not join with them in marriage.}
\markboth{Genesis 24:38}{Genesis 24:38}\vs{38} but you must go to the family of my father and to my relatives to find a wife for my son.'\gva{t}\marginnote{\gva{t}\,Meaning among his relatives, as in Genesis 24:40 .}
\markboth{Genesis 24:39}{Genesis 24:39}\vs{39} But I said to my master, `What if the woman does not want to go with me?'
\markboth{Genesis 24:40}{Genesis 24:40}\vs{40} He answered, `The Lord, before whom I have walked, will send his angel with you. He will make your journey a success and you will find a wife for my son from among my relatives, from my father's family.
\markboth{Genesis 24:41}{Genesis 24:41}\vs{41} You will be free from your oath if you go to my relatives and they will not give her to you. Then you will be free from your oath.'\gva{u}\marginnote{\gva{u}\,Which by my authority I caused you to make.}
\markboth{Genesis 24:42}{Genesis 24:42}\vs{42} When I came to the spring today, I prayed, `O Lord, God of my master Abraham, if you have decided to make my journey successful, may events unfold as follows:
\markboth{Genesis 24:43}{Genesis 24:43}\vs{43} Here I am, standing by the spring. When the young woman goes out to draw water, I'll say, ``Please give me a little water to drink from your jug.''
\markboth{Genesis 24:44}{Genesis 24:44}\vs{44} Then she will reply to me, ``Drink, and I'll draw water for your camels too.'' May that woman be the one whom the Lord has chosen for my master's son.'
\everypar{}

\parshape=0
\markboth{Genesis 24:45}{Genesis 24:45}\vs{45} ``Before I finished praying in my heart, along came Rebekah with her water jug on her shoulder! She went down to the spring and drew water. So I said to her, `Please give me a drink.'\gva{x}\marginnote{\gva{x}\,Signifying that this prayer was not spoken by the mouth, but only in his heart.}
\markboth{Genesis 24:46}{Genesis 24:46}\vs{46} She quickly lowered her jug from her shoulder and said, `Drink, and I'll give your camels water too.' So I drank, and she also gave the camels water.
\markboth{Genesis 24:47}{Genesis 24:47}\vs{47} Then I asked her, `Whose daughter are you?' She replied, `The daughter of Bethuel the son of Nahor, whom Milcah bore to Nahor.' I put the ring in her nose and the bracelets on her wrists.
\markboth{Genesis 24:48}{Genesis 24:48}\vs{48} Then I bowed down and worshiped the Lord. I praised the Lord, the God of my master Abraham, who had led me on the right path to find the granddaughter of my master's brother for his son.\gva{y}\marginnote{\gva{y}\,He shows our duty when we have received any benefit from the Lord.}
\markboth{Genesis 24:49}{Genesis 24:49}\vs{49} Now, if you will show faithful love to my master, tell me. But if not, tell me as well, so that I may go on my way.''\gva{z}\marginnote{\gva{z}\,If you will freely and faithfully give your daughter to my master's son.}\gva{a}\marginnote{\gva{a}\,That is, that I may look elsewhere.}
\everypar{}

\parshape=0
\markboth{Genesis 24:50}{Genesis 24:50}\vs{50} Then Laban and Bethuel replied, ``This is the Lord's doing. Our wishes are of no concern.\gva{b}\marginnote{\gva{b}\,As soon as they perceived that it was God's ordinance they yielded.}
\markboth{Genesis 24:51}{Genesis 24:51}\vs{51} Rebekah stands here before you. Take her and go so that she may become the wife of your master's son, just as the Lord has decided.''
\everypar{}

\parshape=0
\markboth{Genesis 24:52}{Genesis 24:52}\vs{52} When Abraham's servant heard their words, he bowed down to the ground before the Lord.
\markboth{Genesis 24:53}{Genesis 24:53}\vs{53} Then he brought out gold, silver jewelry, and clothing and gave them to Rebekah. He also gave valuable gifts to her brother and to her mother.
\everypar{}

\parshape=0
\markboth{Genesis 24:54}{Genesis 24:54}\vs{54} After this, he and the men who were with him ate a meal and stayed there overnight. When they got up in the morning, he said, ``Let me leave now so I can return to my master.''
\markboth{Genesis 24:55}{Genesis 24:55}\vs{55} But Rebekah's brother and her mother replied, ``Let the girl stay with us a few more days, perhaps ten. Then she can go.''
\markboth{Genesis 24:56}{Genesis 24:56}\vs{56} But he said to them, ``Don't detain me---the Lord has granted me success on my journey. Let me leave now so I may return to my master.''
\markboth{Genesis 24:57}{Genesis 24:57}\vs{57} Then they said, ``We'll call the girl and find out what she wants to do.''\gva{c}\marginnote{\gva{c}\,This shows that parents do not have the authority to marry their children without the consent of both parties.}
\markboth{Genesis 24:58}{Genesis 24:58}\vs{58} So they called Rebekah and asked her, ``Do you want to go with this man?'' She replied, ``I want to go.''
\everypar{}

\parshape=0
\markboth{Genesis 24:59}{Genesis 24:59}\vs{59} So they sent their sister Rebekah on her way, accompanied by her female attendant, with Abraham's servant and his men.
\everypar{}

\parshape=0
\markboth{Genesis 24:60}{Genesis 24:60}\vs{60} They blessed Rebekah with these words: ``Our sister, may you become the mother of thousands of ten thousands! May your descendants possess the strongholds of their enemies.''\gva{d}\marginnote{\gva{d}\,That is, let them be victorious over their enemies: a blessing that is fully accomplished through Jesus Christ.}
\everypar{}

\parshape=0
\markboth{Genesis 24:61}{Genesis 24:61}\vs{61} Then Rebekah and her female servants mounted the camels and rode away with the man. So Abraham's servant took Rebekah and left.
\everypar{}

\parshape=0
\markboth{Genesis 24:62}{Genesis 24:62}\vs{62} Now Isaac came from Beer Lahai Roi, for he was living in the Negev.
\markboth{Genesis 24:63}{Genesis 24:63}\vs{63} He went out to relax in the field in the early evening. Then he looked up and saw that there were camels approaching.\gva{e}\marginnote{\gva{e}\,This was the habit of the godly fathers to meditate on God's promises, and to pray for the accomplishment of it. The custom was that the bride was brought to her husband, her head covered, a token of humbleness and purity.}
\markboth{Genesis 24:64}{Genesis 24:64}\vs{64} Rebekah looked up and saw Isaac. She got down from her camel
\markboth{Genesis 24:65}{Genesis 24:65}\vs{65} and asked Abraham's servant, ``Who is that man walking in the field toward us?'' ``That is my master,'' the servant replied. So she took her veil and covered herself.
\everypar{}

\parshape=0
\markboth{Genesis 24:66}{Genesis 24:66}\vs{66} The servant told Isaac everything that had happened.
\markboth{Genesis 24:67}{Genesis 24:67}\vs{67} Then Isaac brought Rebekah into his mother Sarah's tent. He took her as his wife and loved her. So Isaac was comforted after his mother's death.

\markboth{Genesis 25:1}{Genesis 25:1}
\Needspace*{8\baselineskip}
\ch{25} \allowchapbreak\hypertarget{ch-genesis-25}{}\lettrine[lines=5]{\color{pentateuch}A}{\textsc{braham}} had taken another wife, named Keturah.\gva{a}\marginnote{\gva{a}\,While Sarah was yet alive.}\everypar{}
\markboth{Genesis 25:2}{Genesis 25:2}\vs{2} She bore him Zimran, Jokshan, Medan, Midian, Ishbak, and Shuah.
\markboth{Genesis 25:3}{Genesis 25:3}\vs{3} Jokshan became the father of Sheba and Dedan. The descendants of Dedan were the Asshurites, Letushites, and Leummites.
\markboth{Genesis 25:4}{Genesis 25:4}\vs{4} The sons of Midian were Ephah, Epher, Hanoch, Abida, and Eldaah. All these were descendants of Keturah.
\\\indent\markboth{Genesis 25:5}{Genesis 25:5}\vs{5} Everything he owned Abraham left to his son Isaac.
\markboth{Genesis 25:6}{Genesis 25:6}\vs{6} But while he was still alive, Abraham gave gifts to the sons of his concubines and sent them off to the east, away from his son Isaac.\gva{b}\marginnote{\gva{b}\,For by virtue of God's word he not only had Isaac, but begat many more.}\gva{c}\marginnote{\gva{c}\,See Genesis 22:24 .}\gva{d}\marginnote{\gva{d}\,To avoid the disputing that otherwise might have come because of the heritage.}
\everypar{}

\parshape=0
\markboth{Genesis 25:7}{Genesis 25:7}\vs{7} Abraham lived a total of 175 years.
\markboth{Genesis 25:8}{Genesis 25:8}\vs{8} Then Abraham breathed his last and died at a good old age, an old man who had lived a full life. He joined his ancestors.
\markboth{Genesis 25:9}{Genesis 25:9}\vs{9} His sons Isaac and Ishmael buried him in the cave of Machpelah near Mamre, in the field of Ephron the son of Zohar, the Hittite.
\markboth{Genesis 25:10}{Genesis 25:10}\vs{10} This was the field Abraham had purchased from the sons of Heth. There Abraham was buried with his wife Sarah.
\markboth{Genesis 25:11}{Genesis 25:11}\vs{11} After Abraham's death, God blessed his son Isaac. Isaac lived near Beer Lahai Roi.
\everypar{}

\parshape=0
\markboth{Genesis 25:12}{Genesis 25:12}\vs{12} This is the account of Abraham's son Ishmael, whom Hagar the Egyptian, Sarah's servant, bore to Abraham.
\everypar{}

\parshape=0
\markboth{Genesis 25:13}{Genesis 25:13}\vs{13} These are the names of Ishmael's sons, by their names according to their records: Nebaioth (Ishmael's firstborn), Kedar, Adbeel, Mibsam,
\markboth{Genesis 25:14}{Genesis 25:14}\vs{14} Mishma, Dumah, Massa,
\markboth{Genesis 25:15}{Genesis 25:15}\vs{15} Hadad, Tema, Jetur, Naphish, and Kedemah.
\markboth{Genesis 25:16}{Genesis 25:16}\vs{16} These are the sons of Ishmael, and these are their names by their settlements and their camps---twelve princes according to their clans.
\everypar{}

\parshape=0
\markboth{Genesis 25:17}{Genesis 25:17}\vs{17} Ishmael lived a total of 137 years. He breathed his last and died; then he joined his ancestors.\gva{f}\marginnote{\gva{f}\,Who dwelt among the Arabians, and were separate from the blessed seed.}
\markboth{Genesis 25:18}{Genesis 25:18}\vs{18} His descendants settled from Havilah to Shur, which runs next to Egypt all the way to Asshur. They settled away from all their relatives.\gva{g}\marginnote{\gva{g}\,He means that his lot fell to dwell among his brethren as the angel promised.}
\everypar{}

\parshape=0
\markboth{Genesis 25:19}{Genesis 25:19}\vs{19} This is the account of Isaac, the son of Abraham. Abraham became the father of Isaac.
\markboth{Genesis 25:20}{Genesis 25:20}\vs{20} When Isaac was forty years old, he married Rebekah, the daughter of Bethuel the Aramean from Paddan Aram and sister of Laban the Aramean.
\everypar{}

\parshape=0
\markboth{Genesis 25:21}{Genesis 25:21}\vs{21} Isaac prayed to the Lord on behalf of his wife because she was childless. The Lord answered his prayer, and his wife Rebekah became pregnant.
\markboth{Genesis 25:22}{Genesis 25:22}\vs{22} But the children struggled inside her, and she said, ``Why is this happening to me?'' So she asked the Lord,\gva{h}\marginnote{\gva{h}\,That is, with child, seeing one shall destroy another.}\gva{i}\marginnote{\gva{i}\,For that is the only refuge in all our miseries.}
\everypar{}

\parshape=0
\markboth{Genesis 25:23}{Genesis 25:23}\vs{23} and the Lord said to her, ``Two nations are in your womb, and two peoples will be separated from within you. One people will be stronger than the other, and the older will serve the younger.''
\everypar{}

\parshape=0
\markboth{Genesis 25:24}{Genesis 25:24}\vs{24} When the time came for Rebekah to give birth, there were twins in her womb.
\markboth{Genesis 25:25}{Genesis 25:25}\vs{25} The first came out reddish all over, like a hairy garment, so they named him Esau.
\markboth{Genesis 25:26}{Genesis 25:26}\vs{26} When his brother came out with his hand clutching Esau's heel, they named him Jacob. Isaac was sixty years old when they were born.
\everypar{}

\parshape=0
\markboth{Genesis 25:27}{Genesis 25:27}\vs{27} When the boys grew up, Esau became a skilled hunter, a man of the open fields, but Jacob was an even-tempered man, living in tents.
\markboth{Genesis 25:28}{Genesis 25:28}\vs{28} Isaac loved Esau because he had a taste for fresh game, but Rebekah loved Jacob.
\everypar{}

\parshape=0
\markboth{Genesis 25:29}{Genesis 25:29}\vs{29} Now Jacob cooked some stew, and when Esau came in from the open fields, he was famished.
\markboth{Genesis 25:30}{Genesis 25:30}\vs{30} So Esau said to Jacob, ``Feed me some of the red stuff---yes, this red stuff---because I'm starving!'' (That is why he was also called Edom.)
\everypar{}

\parshape=0
\markboth{Genesis 25:31}{Genesis 25:31}\vs{31} But Jacob replied, ``First sell me your birthright.''
\markboth{Genesis 25:32}{Genesis 25:32}\vs{32} ``Look,'' said Esau, ``I'm about to die! What use is the birthright to me?''\gva{k}\marginnote{\gva{k}\,The reprobate do not value God's benefits unless they feel them presently, and therefore they prefer present pleasures.}
\markboth{Genesis 25:33}{Genesis 25:33}\vs{33} But Jacob said, ``Swear an oath to me now.'' So Esau swore an oath to him and sold his birthright to Jacob.\gva{l}\marginnote{\gva{l}\,Thus the wicked prefer their worldly conveniences over God's spiritual graces: but the children of God do the opposite.}
\everypar{}

\parshape=0
\markboth{Genesis 25:34}{Genesis 25:34}\vs{34} Then Jacob gave Esau some bread and lentil stew; Esau ate and drank, then got up and went out. So Esau despised his birthright.

\markboth{Genesis 26:1}{Genesis 26:1}
\Needspace*{8\baselineskip}
\ch{26} \allowchapbreak\hypertarget{ch-genesis-26}{}\lettrine[lines=5]{\color{pentateuch}T}{\textsc{here}} was a famine in the land, subsequent to the earlier famine that occurred in the days of Abraham. Isaac went to Abimelech king of the Philistines at Gerar.\gva{a}\marginnote{\gva{a}\,In the land of Canaan.}\everypar{}
\markboth{Genesis 26:2}{Genesis 26:2}\vs{2} The Lord appeared to Isaac and said, ``Do not go down to Egypt; settle down in the land that I will point out to you.\gva{b}\marginnote{\gva{b}\,God's providence always watches to direct the ways of his children.}
\markboth{Genesis 26:3}{Genesis 26:3}\vs{3} Stay in this land. Then I will be with you and will bless you, for I will give all these lands to you and to your descendants, and I will fulfill the solemn promise I made to your father Abraham.
\markboth{Genesis 26:4}{Genesis 26:4}\vs{4} I will multiply your descendants so they will be as numerous as the stars in the sky, and I will give them all these lands. All the nations of the earth will pronounce blessings on one another using the name of your descendants.
\markboth{Genesis 26:5}{Genesis 26:5}\vs{5} All this will come to pass because Abraham obeyed me and kept my charge, my commandments, my statutes, and my laws.''\gva{c}\marginnote{\gva{c}\,He commends Abraham's obedience, because Isaac should be even more ready to follow the same: for as God made this promise of his free mercy, so does the confirmation of it proceed from the same fountain.}
\markboth{Genesis 26:6}{Genesis 26:6}\vs{6} So Isaac settled in Gerar.
\everypar{}

\parshape=0
\markboth{Genesis 26:7}{Genesis 26:7}\vs{7} When the men of that place asked him about his wife, he replied, ``She is my sister.'' He was afraid to say, ``She is my wife,'' for he thought to himself, ``The men of this place will kill me to get Rebekah because she is very beautiful.''\gva{d}\marginnote{\gva{d}\,By which we see that fear and distrust is found in the most faithful.}
\everypar{}

\parshape=0
\markboth{Genesis 26:8}{Genesis 26:8}\vs{8} After Isaac had been there a long time, Abimelech king of the Philistines happened to look out a window and observed Isaac caressing his wife Rebekah.\gva{e}\marginnote{\gva{e}\,Or showing some familiar sign of love, by which it might be known that she was his wife.}
\markboth{Genesis 26:9}{Genesis 26:9}\vs{9} So Abimelech summoned Isaac and said, ``She is really your wife! Why did you say, `She is my sister'?'' Isaac replied, ``Because I thought someone might kill me to get her.''
\everypar{}

\parshape=0
\markboth{Genesis 26:10}{Genesis 26:10}\vs{10} Then Abimelech exclaimed, ``What in the world have you done to us? One of the men nearly took your wife to bed, and you would have brought guilt on us!''\gva{f}\marginnote{\gva{f}\,In all ages men were persuaded that God's vengeance would come on adulterers.}
\markboth{Genesis 26:11}{Genesis 26:11}\vs{11} So Abimelech commanded all the people, ``Whoever touches this man or his wife will surely be put to death.''
\everypar{}

\parshape=0
\markboth{Genesis 26:12}{Genesis 26:12}\vs{12} When Isaac planted in that land, he reaped in the same year a hundred times what he had sown, because the Lord blessed him.
\markboth{Genesis 26:13}{Genesis 26:13}\vs{13} The man became wealthy. His influence continued to grow until he became very prominent.
\markboth{Genesis 26:14}{Genesis 26:14}\vs{14} He had so many sheep and cattle and such a great household of servants that the Philistines became jealous of him.\gva{g}\marginnote{\gva{g}\,The malicious always envy the graces of God in others.}
\markboth{Genesis 26:15}{Genesis 26:15}\vs{15} So the Philistines took dirt and filled up all the wells that his father's servants had dug back in the days of his father Abraham.
\everypar{}

\parshape=0
\markboth{Genesis 26:16}{Genesis 26:16}\vs{16} Then Abimelech said to Isaac, ``Leave us and go elsewhere, for you have become much more powerful than we are.''
\markboth{Genesis 26:17}{Genesis 26:17}\vs{17} So Isaac left there and settled in the Gerar Valley.\gva{k}\marginnote{\gva{k}\,The Hebrew word signifies a flood, or valley, where water at any time runs.}
\markboth{Genesis 26:18}{Genesis 26:18}\vs{18} Isaac reopened the wells that had been dug back in the days of his father Abraham, for the Philistines had stopped them up after Abraham died. Isaac gave these wells the same names his father had given them.
\everypar{}

\parshape=0
\markboth{Genesis 26:19}{Genesis 26:19}\vs{19} When Isaac's servants dug in the valley and discovered a well with fresh flowing water there,
\markboth{Genesis 26:20}{Genesis 26:20}\vs{20} the herdsmen of Gerar quarreled with Isaac's herdsmen, saying, ``The water belongs to us!'' So Isaac named the well Esek because they argued with him about it.
\markboth{Genesis 26:21}{Genesis 26:21}\vs{21} His servants dug another well, but they quarreled over it too, so Isaac named it Sitnah.
\markboth{Genesis 26:22}{Genesis 26:22}\vs{22} Then he moved away from there and dug another well. They did not quarrel over it, so Isaac named it Rehoboth, saying, ``For now the Lord has made room for us, and we will prosper in the land.''
\everypar{}

\parshape=0
\markboth{Genesis 26:23}{Genesis 26:23}\vs{23} From there Isaac went up to Beer Sheba.
\markboth{Genesis 26:24}{Genesis 26:24}\vs{24} The Lord appeared to him that night and said, ``I am the God of your father Abraham. Do not be afraid, for I am with you. I will bless you and multiply your descendants for the sake of my servant Abraham.''\gva{i}\marginnote{\gva{i}\,God assures Isaac against all fear by rehearsing the promise made to Abraham.}
\markboth{Genesis 26:25}{Genesis 26:25}\vs{25} Then Isaac built an altar there and worshiped the Lord. He pitched his tent there, and his servants dug a well.\gva{x}\marginnote{\gva{x}\,To signify that he would serve no other God, but the God of his father Abraham.}
\everypar{}

\parshape=0
\markboth{Genesis 26:26}{Genesis 26:26}\vs{26} Now Abimelech had come to him from Gerar along with Ahuzzah his friend and Phicol the commander of his army.
\markboth{Genesis 26:27}{Genesis 26:27}\vs{27} Isaac asked them, ``Why have you come to me? You hate me and sent me away from you.''
\markboth{Genesis 26:28}{Genesis 26:28}\vs{28} They replied, ``We could plainly see that the Lord is with you. So we decided there should be a pact between us---between us and you. Allow us to make a treaty with you
\markboth{Genesis 26:29}{Genesis 26:29}\vs{29} so that you will not do us any harm, just as we have not harmed you, but have always treated you well before sending you away in peace. Now you are blessed by the Lord.''\gva{l}\marginnote{\gva{l}\,The Hebrews in swearing begin commonly with "If" and understand the rest, that is, that God will punish him who breaks the oath: here the wicked show that they are afraid lest that happen to them which they would do to others.}
\everypar{}

\parshape=0
\markboth{Genesis 26:30}{Genesis 26:30}\vs{30} So Isaac held a feast for them and they celebrated.
\markboth{Genesis 26:31}{Genesis 26:31}\vs{31} Early in the morning the men made a treaty with each other. Isaac sent them off; they separated on good terms.
\everypar{}

\parshape=0
\markboth{Genesis 26:32}{Genesis 26:32}\vs{32} That day Isaac's servants came and told him about the well they had dug. ``We've found water,'' they reported.
\markboth{Genesis 26:33}{Genesis 26:33}\vs{33} So he named it Shibah; that is why the name of the city has been Beer Sheba to this day.
\everypar{}

\parshape=0
\markboth{Genesis 26:34}{Genesis 26:34}\vs{34} When Esau was forty years old, he married Judith the daughter of Beeri the Hittite, as well as Basemath the daughter of Elon the Hittite.
\markboth{Genesis 26:35}{Genesis 26:35}\vs{35} They caused Isaac and Rebekah great anxiety.

\markboth{Genesis 27:1}{Genesis 27:1}
\Needspace*{8\baselineskip}
\ch{27} \allowchapbreak\hypertarget{ch-genesis-27}{}\lettrine[lines=5]{\color{pentateuch}W}{\textsc{hen}} Isaac was old and his eyes were so weak that he was almost blind, he called his older son Esau and said to him, ``My son!'' ``Here I am!'' Esau replied.\everypar{}
\markboth{Genesis 27:2}{Genesis 27:2}\vs{2} Isaac said, ``Since I am so old, I could die at any time.
\markboth{Genesis 27:3}{Genesis 27:3}\vs{3} Therefore, take your weapons---your quiver and your bow---and go out into the open fields and hunt down some wild game for me.
\markboth{Genesis 27:4}{Genesis 27:4}\vs{4} Then prepare for me some tasty food, the kind I love, and bring it to me. Then I will eat it so that I may bless you before I die.''\gva{a}\marginnote{\gva{a}\,The carnal affection he had for his son made him forget what God spoke to his wife. Genesis 25:23 .}
\everypar{}

\parshape=0
\markboth{Genesis 27:5}{Genesis 27:5}\vs{5} Now Rebekah had been listening while Isaac spoke to his son Esau. When Esau went out to the open fields to hunt down some wild game and bring it back,
\markboth{Genesis 27:6}{Genesis 27:6}\vs{6} Rebekah said to her son Jacob, ``Look, I overheard your father tell your brother Esau,
\markboth{Genesis 27:7}{Genesis 27:7}\vs{7} `Bring me some wild game and prepare for me some tasty food. Then I will eat it and bless you in the presence of the Lord before I die.'
\markboth{Genesis 27:8}{Genesis 27:8}\vs{8} Now then, my son, do exactly what I tell you!
\markboth{Genesis 27:9}{Genesis 27:9}\vs{9} Go to the flock and get me two of the best young goats. I'll prepare them in a tasty way for your father, just the way he loves them.\gva{b}\marginnote{\gva{b}\,This trickery is worthy of blame because she should have waited for God to perform his promise.}
\markboth{Genesis 27:10}{Genesis 27:10}\vs{10} Then you will take it to your father. Thus he will eat it and bless you before he dies.''
\everypar{}

\parshape=0
\markboth{Genesis 27:11}{Genesis 27:11}\vs{11} ``But Esau my brother is a hairy man,'' Jacob protested to his mother Rebekah, ``and I have smooth skin!
\markboth{Genesis 27:12}{Genesis 27:12}\vs{12} My father may touch me! Then he'll think I'm mocking him and I'll bring a curse on myself instead of a blessing.''
\markboth{Genesis 27:13}{Genesis 27:13}\vs{13} So his mother told him, ``Any curse against you will fall on me, my son! Just obey me! Go and get them for me!''\gva{c}\marginnote{\gva{c}\,The assurance of God's decree made her bold.}
\everypar{}

\parshape=0
\markboth{Genesis 27:14}{Genesis 27:14}\vs{14} So he went and got the goats and brought them to his mother. She prepared some tasty food, just the way his father loved it.
\markboth{Genesis 27:15}{Genesis 27:15}\vs{15} Then Rebekah took her older son Esau's best clothes, which she had with her in the house, and put them on her younger son Jacob.
\markboth{Genesis 27:16}{Genesis 27:16}\vs{16} She put the skins of the young goats on his hands and the smooth part of his neck.
\markboth{Genesis 27:17}{Genesis 27:17}\vs{17} Then she handed the tasty food and the bread she had made to her son Jacob.
\everypar{}

\parshape=0
\markboth{Genesis 27:18}{Genesis 27:18}\vs{18} He went to his father and said, ``My father!'' Isaac replied, ``Here I am. Which are you, my son?''
\markboth{Genesis 27:19}{Genesis 27:19}\vs{19} Jacob said to his father, ``I am Esau, your firstborn. I've done as you told me. Now sit up and eat some of my wild game so that you can bless me.''\gva{d}\marginnote{\gva{d}\,Although Jacob was assured of this blessing by faith: yet he did evil to seek it by lies, even more because he abuses God's name through it.}
\markboth{Genesis 27:20}{Genesis 27:20}\vs{20} But Isaac asked his son, ``How in the world did you find it so quickly, my son?'' ``Because the Lord your God brought it to me,'' he replied.
\markboth{Genesis 27:21}{Genesis 27:21}\vs{21} Then Isaac said to Jacob, ``Come closer so I can touch you, my son, and know for certain if you really are my son Esau.''
\markboth{Genesis 27:22}{Genesis 27:22}\vs{22} So Jacob went over to his father Isaac, who felt him and said, ``The voice is Jacob's, but the hands are Esau's.''\gva{e}\marginnote{\gva{e}\,This declares that he suspected something, yet God would not have his decree altered.}
\markboth{Genesis 27:23}{Genesis 27:23}\vs{23} He did not recognize him because his hands were hairy, like his brother Esau's hands. So Isaac blessed Jacob.
\markboth{Genesis 27:24}{Genesis 27:24}\vs{24} Then he asked, ``Are you really my son Esau?'' ``I am,'' Jacob replied.
\markboth{Genesis 27:25}{Genesis 27:25}\vs{25} Isaac said, ``Bring some of the wild game for me to eat, my son. Then I will bless you.'' So Jacob brought it to him, and he ate it. He also brought him wine, and Isaac drank.
\markboth{Genesis 27:26}{Genesis 27:26}\vs{26} Then his father Isaac said to him, ``Come here and kiss me, my son.''
\everypar{}

\parshape=0
\markboth{Genesis 27:27}{Genesis 27:27}\vs{27} So Jacob went over and kissed him. When Isaac caught the scent of his clothing, he blessed him, saying, ``Yes, my son smells like the scent of an open field which the Lord has blessed.
\everypar{}

\parshape=0
\markboth{Genesis 27:28}{Genesis 27:28}\vs{28} May God give you the dew of the sky and the richness of the earth, and plenty of grain and new wine.
\everypar{}

\parshape=0
\markboth{Genesis 27:29}{Genesis 27:29}\vs{29} May peoples serve you and nations bow down to you. You will be lord over your brothers, and the sons of your mother will bow down to you. May those who curse you be cursed, and those who bless you be blessed.''
\everypar{}

\parshape=0
\markboth{Genesis 27:30}{Genesis 27:30}\vs{30} Isaac had just finished blessing Jacob, and Jacob had scarcely left his father's presence, when his brother Esau returned from the hunt.
\markboth{Genesis 27:31}{Genesis 27:31}\vs{31} He also prepared some tasty food and brought it to his father. Esau said to him, ``My father, get up and eat some of your son's wild game. Then you can bless me.''
\markboth{Genesis 27:32}{Genesis 27:32}\vs{32} His father Isaac asked, ``Who are you?'' ``I am your firstborn son,'' he replied, ``Esau!''
\markboth{Genesis 27:33}{Genesis 27:33}\vs{33} Isaac began to shake violently and asked, ``Then who else hunted game and brought it to me? I ate all of it just before you arrived, and I blessed him. He will indeed be blessed!''\gva{f}\marginnote{\gva{f}\,In perceiving his error, by appointing his heir against God's sentence pronounced before.}
\everypar{}

\parshape=0
\markboth{Genesis 27:34}{Genesis 27:34}\vs{34} When Esau heard his father's words, he wailed loudly and bitterly. He said to his father, ``Bless me too, my father!''
\markboth{Genesis 27:35}{Genesis 27:35}\vs{35} But Isaac replied, ``Your brother came in here deceitfully and took away your blessing.''
\markboth{Genesis 27:36}{Genesis 27:36}\vs{36} Esau exclaimed, ``Jacob is the right name for him! He has tripped me up two times! He took away my birthright, and now, look, he has taken away my blessing!'' Then he asked, ``Have you not kept back a blessing for me?''\gva{g}\marginnote{\gva{g}\,In Genesis 25:26 he was so called because he held his brother by the heel, as though he would overthrow him: and therefore he is here called an overthrower, or deceiver.}
\everypar{}

\parshape=0
\markboth{Genesis 27:37}{Genesis 27:37}\vs{37} Isaac replied to Esau, ``Look! I have made him lord over you. I have made all his relatives his servants and provided him with grain and new wine. What is left that I can do for you, my son?''\gva{h}\marginnote{\gva{h}\,thy lord, and all his brethren have I given to him for servants; and with corn and wine have I sustained him: and what shall I do now unto thee, my son?}\gva{h}\marginnote{\gva{h}\,For Isaac did this as he was the minister and prophet of God.}
\markboth{Genesis 27:38}{Genesis 27:38}\vs{38} Esau said to his father, ``Do you have only that one blessing, my father? Bless me too!'' Then Esau wept loudly.
\everypar{}

\parshape=0
\markboth{Genesis 27:39}{Genesis 27:39}\vs{39} So his father Isaac said to him, ``See here, your home will be by the richness of the earth, and by the dew of the sky above.
\everypar{}

\parshape=0
\markboth{Genesis 27:40}{Genesis 27:40}\vs{40} You will live by your sword but you will serve your brother. When you grow restless, you will tear off his yoke from your neck.''\gva{i}\marginnote{\gva{i}\,Because your enemies will be all around you.}\gva{k}\marginnote{\gva{k}\,Which was fulfilled in his posterity the Idumeans: who were tributaries for a time to Israel, and later came to freedom.}
\everypar{}

\parshape=0
\markboth{Genesis 27:41}{Genesis 27:41}\vs{41} So Esau hated Jacob because of the blessing his father had given to his brother. Esau said privately, ``The time of mourning for my father is near; then I will kill my brother Jacob!''\gva{l}\marginnote{\gva{l}\,Hypocrites only abstain from doing evil for fear of men.}
\everypar{}

\parshape=0
\markboth{Genesis 27:42}{Genesis 27:42}\vs{42} When Rebekah heard what her older son Esau had said, she quickly summoned her younger son Jacob and told him, ``Look, your brother Esau is planning to get revenge by killing you.\gva{m}\marginnote{\gva{m}\,He hopes to recover his birthright by killing you.}
\markboth{Genesis 27:43}{Genesis 27:43}\vs{43} Now then, my son, do what I say. Run away immediately to my brother Laban in Haran.
\markboth{Genesis 27:44}{Genesis 27:44}\vs{44} Live with him for a little while until your brother's rage subsides.
\markboth{Genesis 27:45}{Genesis 27:45}\vs{45} Stay there until your brother's anger against you subsides and he forgets what you did to him. Then I'll send someone to bring you back from there. Why should I lose both of you in one day?''\gva{n}\marginnote{\gva{n}\,For the wicked son will kill the godly: and the plague of God will later come on the wicked son.}
\everypar{}

\parshape=0
\markboth{Genesis 27:46}{Genesis 27:46}\vs{46} Then Rebekah said to Isaac, ``I am deeply depressed because of the daughters of Heth. If Jacob were to marry one of these daughters of Heth who live in this land, I would want to die!''\gva{o}\marginnote{\gva{o}\,Who were Esau's wives.}\gva{p}\marginnote{\gva{p}\,By this she persuaded Isaac to agree to Jacob's leaving.}

\markboth{Genesis 28:1}{Genesis 28:1}
\Needspace*{8\baselineskip}
\ch{28} \allowchapbreak\hypertarget{ch-genesis-28}{}\lettrine[lines=5]{\color{pentateuch}S}{\textsc{o}} Isaac called for Jacob and blessed him. Then he commanded him, ``You must not marry a Canaanite woman!\gva{a}\marginnote{\gva{a}\,This second blessing was to confirm Jacob's faith, lest he should think that his father had given it without God's leading.}\everypar{}
\markboth{Genesis 28:2}{Genesis 28:2}\vs{2} Leave immediately for Paddan Aram! Go to the house of Bethuel, your mother's father, and find yourself a wife there, among the daughters of Laban, your mother's brother.
\markboth{Genesis 28:3}{Genesis 28:3}\vs{3} May the Sovereign God bless you! May he make you fruitful and give you a multitude of descendants! Then you will become a large nation.
\markboth{Genesis 28:4}{Genesis 28:4}\vs{4} May he give you and your descendants the blessing he gave to Abraham so that you may possess the land God gave to Abraham, the land where you have been living as a temporary resident.''\gva{b}\marginnote{\gva{b}\,The godly fathers were continually reminded that they were but strangers in this world: so that they would lift up their eyes to the heavens where they have a certain dwelling.}
\markboth{Genesis 28:5}{Genesis 28:5}\vs{5} So Isaac sent Jacob on his way, and he went to Paddan Aram, to Laban son of Bethuel the Aramean and brother of Rebekah, the mother of Jacob and Esau.
\everypar{}

\parshape=0
\markboth{Genesis 28:6}{Genesis 28:6}\vs{6} Esau saw that Isaac had blessed Jacob and sent him off to Paddan Aram to find a wife there. As he blessed him, Isaac commanded him, ``You must not marry a Canaanite woman.''
\markboth{Genesis 28:7}{Genesis 28:7}\vs{7} Jacob obeyed his father and mother and left for Paddan Aram.
\markboth{Genesis 28:8}{Genesis 28:8}\vs{8} Then Esau realized that the Canaanite women were displeasing to his father Isaac.
\markboth{Genesis 28:9}{Genesis 28:9}\vs{9} So Esau went to Ishmael and married Mahalath, the sister of Nebaioth and daughter of Abraham's son Ishmael, along with the wives he already had.\gva{c}\marginnote{\gva{c}\,Thinking by this to have reconciled himself to his father, but all in vain: for he does not take away the cause of the evil.}
\everypar{}

\parshape=0
\markboth{Genesis 28:10}{Genesis 28:10}\vs{10} Meanwhile Jacob left Beer Sheba and set out for Haran.
\markboth{Genesis 28:11}{Genesis 28:11}\vs{11} He reached a certain place where he decided to camp because the sun had gone down. He took one of the stones and placed it near his head. Then he fell asleep in that place
\markboth{Genesis 28:12}{Genesis 28:12}\vs{12} and had a dream. He saw a stairway erected on the earth with its top reaching to the heavens. The angels of God were going up and coming down it\gva{d}\marginnote{\gva{d}\,Christ is the ladder by which God and man are joined together, and by whom the angels minister to us: all graces are given to us by him, and we ascend to heaven by him.}
\markboth{Genesis 28:13}{Genesis 28:13}\vs{13} and the Lord stood at its top. He said, ``I am the Lord, the God of your grandfather Abraham and the God of your father Isaac. I will give you and your descendants the ground you are lying on.\gva{e}\marginnote{\gva{e}\,He felt the force of this promise only by faith: for all his life he was a stranger in this land.}
\markboth{Genesis 28:14}{Genesis 28:14}\vs{14} Your descendants will be like the dust of the earth, and you will spread out to the west, east, north, and south. And so all the families of the earth may receive blessings through you and through your descendants.
\markboth{Genesis 28:15}{Genesis 28:15}\vs{15} I am with you! I will protect you wherever you go and will bring you back to this land. I will not leave you until I have done what I promised you!''
\everypar{}

\parshape=0
\markboth{Genesis 28:16}{Genesis 28:16}\vs{16} Then Jacob woke up and thought, ``Surely the Lord is in this place, but I did not realize it!''
\markboth{Genesis 28:17}{Genesis 28:17}\vs{17} He was afraid and said, ``What an awesome place this is! This is nothing else than the house of God! This is the gate of heaven!''\gva{f}\marginnote{\gva{f}\,He was touched with a godly fear and reverence.}
\everypar{}

\parshape=0
\markboth{Genesis 28:18}{Genesis 28:18}\vs{18} Early in the morning Jacob took the stone he had placed near his head and set it up as a sacred stone. Then he poured oil on top of it.\gva{g}\marginnote{\gva{g}\,To be a reminder of the vision shown to him.}
\markboth{Genesis 28:19}{Genesis 28:19}\vs{19} He called that place Bethel, although the former name of the town was Luz.
\markboth{Genesis 28:20}{Genesis 28:20}\vs{20} Then Jacob made a vow, saying, ``If God is with me and protects me on this journey I am taking and gives me food to eat and clothing to wear,\gva{h}\marginnote{\gva{h}\,He does not bind God under this condition, but acknowledges his infirmity, and promises to be thankful.}
\markboth{Genesis 28:21}{Genesis 28:21}\vs{21} and I return safely to my father's home, then the Lord will become my God.
\markboth{Genesis 28:22}{Genesis 28:22}\vs{22} Then this stone that I have set up as a sacred stone will be the house of God, and I will surely give you back a tenth of everything you give me.''

\markboth{Genesis 29:1}{Genesis 29:1}
\Needspace*{8\baselineskip}
\ch{29} \allowchapbreak\hypertarget{ch-genesis-29}{}\lettrine[lines=5]{\color{pentateuch}S}{\textsc{o}} Jacob moved on and came to the land of the eastern people.\gva{a}\marginnote{\gva{a}\,Or, "lifted up his feet".}\everypar{}
\markboth{Genesis 29:2}{Genesis 29:2}\vs{2} He saw in the field a well with three flocks of sheep lying beside it, because the flocks were watered from that well. Now a large stone covered the mouth of the well.\gva{b}\marginnote{\gva{b}\,Thus he was directed by the providence of God, who brought him to Laban's house.}
\markboth{Genesis 29:3}{Genesis 29:3}\vs{3} When all the flocks were gathered there, the shepherds would roll the stone off the mouth of the well and water the sheep. Then they would put the stone back in its place over the well's mouth.
\everypar{}

\parshape=0
\markboth{Genesis 29:4}{Genesis 29:4}\vs{4} Jacob asked them, ``My brothers, where are you from?'' They replied, ``We're from Haran.''\gva{c}\marginnote{\gva{c}\,It seems that in those days the custom was to call even strangers, brethren.}
\markboth{Genesis 29:5}{Genesis 29:5}\vs{5} So he said to them, ``Do you know Laban, the grandson of Nahor?'' ``We know him,'' they said.
\markboth{Genesis 29:6}{Genesis 29:6}\vs{6} ``Is he well?'' Jacob asked. They replied, ``He is well. Now look, here comes his daughter Rachel with the sheep.''\gva{d}\marginnote{\gva{d}\,Or, "he is in peace?" by which the Hebrews mean prosperity.}
\markboth{Genesis 29:7}{Genesis 29:7}\vs{7} Then Jacob said, ``Since it is still the middle of the day, it is not time for the flocks to be gathered. You should water the sheep and then go and let them graze some more.''
\markboth{Genesis 29:8}{Genesis 29:8}\vs{8} ``We can't,'' they said, ``until all the flocks are gathered and the stone is rolled off the mouth of the well. Then we water the sheep.''
\everypar{}

\parshape=0
\markboth{Genesis 29:9}{Genesis 29:9}\vs{9} While he was still speaking with them, Rachel arrived with her father's sheep, for she was tending them.
\markboth{Genesis 29:10}{Genesis 29:10}\vs{10} When Jacob saw Rachel, the daughter of his uncle Laban, and the sheep of his uncle Laban, he went over and rolled the stone off the mouth of the well and watered the sheep of his uncle Laban.
\markboth{Genesis 29:11}{Genesis 29:11}\vs{11} Then Jacob kissed Rachel and began to weep loudly.
\markboth{Genesis 29:12}{Genesis 29:12}\vs{12} When Jacob explained to Rachel that he was a relative of her father and the son of Rebekah, she ran and told her father.
\markboth{Genesis 29:13}{Genesis 29:13}\vs{13} When Laban heard this news about Jacob, his sister's son, he rushed out to meet him. He embraced him and kissed him and brought him to his house. Jacob told Laban how he was related to him.\gva{e}\marginnote{\gva{e}\,That is, the reason why he departed from his father's house, and what he saw during his journey.}
\markboth{Genesis 29:14}{Genesis 29:14}\vs{14} Then Laban said to him, ``You are indeed my own flesh and blood.'' So Jacob stayed with him for a month.\gva{f}\marginnote{\gva{f}\,That is, of my blood and kindred.}
\everypar{}

\parshape=0
\markboth{Genesis 29:15}{Genesis 29:15}\vs{15} Then Laban said to Jacob, ``Should you work for me for nothing because you are my relative? Tell me what your wages should be.''
\markboth{Genesis 29:16}{Genesis 29:16}\vs{16} (Now Laban had two daughters; the older one was named Leah, and the younger one Rachel.
\markboth{Genesis 29:17}{Genesis 29:17}\vs{17} Leah's eyes were tender, but Rachel had a lovely figure and beautiful appearance.)
\markboth{Genesis 29:18}{Genesis 29:18}\vs{18} Since Jacob had fallen in love with Rachel, he said, ``I'll serve you seven years in exchange for your younger daughter Rachel.''
\markboth{Genesis 29:19}{Genesis 29:19}\vs{19} Laban replied, ``I'd rather give her to you than to another man. Stay with me.''
\markboth{Genesis 29:20}{Genesis 29:20}\vs{20} So Jacob worked for seven years to acquire Rachel. But they seemed like only a few days to him because his love for her was so great.\gva{g}\marginnote{\gva{g}\,Meaning after the years were accomplished.}
\everypar{}

\parshape=0
\markboth{Genesis 29:21}{Genesis 29:21}\vs{21} Finally Jacob said to Laban, ``Give me my wife, for my time of service is up. And I want to sleep with her.''
\markboth{Genesis 29:22}{Genesis 29:22}\vs{22} So Laban invited all the people of that place and prepared a feast.
\markboth{Genesis 29:23}{Genesis 29:23}\vs{23} In the evening he brought his daughter Leah to Jacob, and he slept with her.\gva{h}\marginnote{\gva{h}\,The reason Jacob was deceived was that in ancient times the wife was covered with a veil, when she was brought to her husband as a sign of purity and humbleness.}
\markboth{Genesis 29:24}{Genesis 29:24}\vs{24} (Laban gave his female servant Zilpah to his daughter Leah to be her servant.)
\everypar{}

\parshape=0
\markboth{Genesis 29:25}{Genesis 29:25}\vs{25} In the morning Jacob discovered it was Leah! So Jacob said to Laban, ``What in the world have you done to me? Didn't I work for you in exchange for Rachel? Why have you tricked me?''
\markboth{Genesis 29:26}{Genesis 29:26}\vs{26} ``It is not our custom here,'' Laban replied, ``to give the younger daughter in marriage before the firstborn.\gva{i}\marginnote{\gva{i}\,He valued the profit he had from Jacob's service more than either his promise or the customs of the country, though he used custom for his excuse.}
\markboth{Genesis 29:27}{Genesis 29:27}\vs{27} Complete my older daughter's bridal week. Then we will give you the younger one too, in exchange for seven more years of work.''
\everypar{}

\parshape=0
\markboth{Genesis 29:28}{Genesis 29:28}\vs{28} Jacob did as Laban said. When Jacob completed Leah's bridal week, Laban gave him his daughter Rachel to be his wife.
\markboth{Genesis 29:29}{Genesis 29:29}\vs{29} (Laban gave his female servant Bilhah to his daughter Rachel to be her servant.)
\markboth{Genesis 29:30}{Genesis 29:30}\vs{30} Jacob slept with Rachel as well. He also loved Rachel more than Leah. Then he worked for Laban for seven more years.
\everypar{}

\parshape=0
\markboth{Genesis 29:31}{Genesis 29:31}\vs{31} When the Lord saw that Leah was unloved, he enabled her to become pregnant while Rachel remained childless.\gva{k}\marginnote{\gva{k}\,This declares that often they who are despised by men are favoured by God.}
\markboth{Genesis 29:32}{Genesis 29:32}\vs{32} So Leah became pregnant and gave birth to a son. She named him Reuben, for she said, ``The Lord has looked with pity on my oppressed condition. Surely my husband will love me now.''\gva{l}\marginnote{\gva{l}\,By this it appears that she had sought help from God in her affliction.}\gva{m}\marginnote{\gva{m}\,For children are a great cause of mutual love between man and wife.}
\everypar{}

\parshape=0
\markboth{Genesis 29:33}{Genesis 29:33}\vs{33} She became pregnant again and had another son. She said, ``Because the Lord heard that I was unloved, he gave me this one too.'' So she named him Simeon.
\everypar{}

\parshape=0
\markboth{Genesis 29:34}{Genesis 29:34}\vs{34} She became pregnant again and had another son. She said, ``Now this time my husband will show me affection, because I have given birth to three sons for him.'' That is why he was named Levi.
\everypar{}

\parshape=0
\markboth{Genesis 29:35}{Genesis 29:35}\vs{35} She became pregnant again and had another son. She said, ``This time I will praise the Lord.'' That is why she named him Judah. Then she stopped having children.

\markboth{Genesis 30:1}{Genesis 30:1}
\Needspace*{8\baselineskip}
\ch{30} \allowchapbreak\hypertarget{ch-genesis-30}{}\lettrine[lines=5]{\color{pentateuch}W}{\textsc{hen}} Rachel saw that she could not give Jacob children, she became jealous of her sister. She said to Jacob, ``Give me children or I'll die!''\everypar{}
\markboth{Genesis 30:2}{Genesis 30:2}\vs{2} Jacob became furious with Rachel and exclaimed, ``Am I in the place of God, who has kept you from having children?''\gva{a}\marginnote{\gva{a}\,It is only God who makes one barren or fruitful, and therefore I am not at fault.}
\markboth{Genesis 30:3}{Genesis 30:3}\vs{3} She replied, ``Here is my servant Bilhah! Sleep with her so that she can bear children for me and I can have a family through her.''\gva{b}\marginnote{\gva{b}\,I will receive her children on my lap, as though they were my own.}
\everypar{}

\parshape=0
\markboth{Genesis 30:4}{Genesis 30:4}\vs{4} So Rachel gave him her servant Bilhah as a wife, and Jacob slept with her.
\markboth{Genesis 30:5}{Genesis 30:5}\vs{5} Bilhah became pregnant and gave Jacob a son.
\markboth{Genesis 30:6}{Genesis 30:6}\vs{6} Then Rachel said, ``God has vindicated me. He has responded to my prayer and given me a son.'' That is why she named him Dan.
\everypar{}

\parshape=0
\markboth{Genesis 30:7}{Genesis 30:7}\vs{7} Bilhah, Rachel's servant, became pregnant again and gave Jacob another son.
\markboth{Genesis 30:8}{Genesis 30:8}\vs{8} Then Rachel said, ``I have fought a desperate struggle with my sister, but I have won.'' So she named him Naphtali.\gva{c}\marginnote{\gva{c}\,The arrogancy of man's nature appears in that she condemns her sister, after she has received this benefit from God to bear children.}
\everypar{}

\parshape=0
\markboth{Genesis 30:9}{Genesis 30:9}\vs{9} When Leah saw that she had stopped having children, she gave her servant Zilpah to Jacob as a wife.
\markboth{Genesis 30:10}{Genesis 30:10}\vs{10} Soon Leah's servant Zilpah gave Jacob a son.
\markboth{Genesis 30:11}{Genesis 30:11}\vs{11} Leah said, ``How fortunate!'' So she named him Gad.\gva{d}\marginnote{\gva{d}\,That is, God increases me with a multitude of children for so Jacob explains this name Gad Genesis 49:19 .}
\everypar{}

\parshape=0
\markboth{Genesis 30:12}{Genesis 30:12}\vs{12} Then Leah's servant Zilpah gave Jacob another son.
\markboth{Genesis 30:13}{Genesis 30:13}\vs{13} Leah said, ``How happy I am, for women will call me happy!'' So she named him Asher.
\everypar{}

\parshape=0
\markboth{Genesis 30:14}{Genesis 30:14}\vs{14} At the time of the wheat harvest Reuben went out and found some mandrake plants in a field and brought them to his mother Leah. Rachel said to Leah, ``Give me some of your son's mandrakes.''\gva{e}\marginnote{\gva{e}\,Which is a kind of herb whose root has a likeness to the figure of a man.}
\markboth{Genesis 30:15}{Genesis 30:15}\vs{15} But Leah replied, ``Wasn't it enough that you've taken away my husband? Would you take away my son's mandrakes too?'' ``All right,'' Rachel said, ``he may go to bed with you tonight in exchange for your son's mandrakes.''
\markboth{Genesis 30:16}{Genesis 30:16}\vs{16} When Jacob came in from the fields that evening, Leah went out to meet him and said, ``You must sleep with me because I have paid for your services with my son's mandrakes.'' So he went to bed with her that night.
\markboth{Genesis 30:17}{Genesis 30:17}\vs{17} God paid attention to Leah; she became pregnant and gave Jacob a son for the fifth time.
\markboth{Genesis 30:18}{Genesis 30:18}\vs{18} Then Leah said, ``God has granted me a reward because I gave my servant to my husband as a wife.'' So she named him Issachar.\gva{f}\marginnote{\gva{f}\,Instead of acknowledging her fault she boasts as if God had rewarded her for it.}
\everypar{}

\parshape=0
\markboth{Genesis 30:19}{Genesis 30:19}\vs{19} Leah became pregnant again and gave Jacob a son for the sixth time.
\markboth{Genesis 30:20}{Genesis 30:20}\vs{20} Then Leah said, ``God has given me a good gift. Now my husband will honor me because I have given him six sons.'' So she named him Zebulun.
\everypar{}

\parshape=0
\markboth{Genesis 30:21}{Genesis 30:21}\vs{21} After that she gave birth to a daughter and named her Dinah.
\everypar{}

\parshape=0
\markboth{Genesis 30:22}{Genesis 30:22}\vs{22} Then God took note of Rachel. He paid attention to her and enabled her to become pregnant.
\markboth{Genesis 30:23}{Genesis 30:23}\vs{23} She became pregnant and gave birth to a son. Then she said, ``God has taken away my shame.''\gva{g}\marginnote{\gva{g}\,Because fruitfulness came as God's blessing, who said "Increase and multiply", barrenness was counted as a curse.}
\markboth{Genesis 30:24}{Genesis 30:24}\vs{24} She named him Joseph, saying, ``May the Lord give me yet another son.''
\everypar{}

\parshape=0
\markboth{Genesis 30:25}{Genesis 30:25}\vs{25} After Rachel had given birth to Joseph, Jacob said to Laban, ``Send me on my way so that I can go home to my own country.
\markboth{Genesis 30:26}{Genesis 30:26}\vs{26} Let me take my wives and my children whom I have acquired by working for you. Then I'll depart, because you know how hard I've worked for you.''
\everypar{}

\parshape=0
\markboth{Genesis 30:27}{Genesis 30:27}\vs{27} But Laban said to him, ``If I have found favor in your sight, please stay here, for I have learned by divination that the Lord has blessed me on account of you.''
\markboth{Genesis 30:28}{Genesis 30:28}\vs{28} He added, ``Just name your wages---I'll pay whatever you want.''
\everypar{}

\parshape=0
\markboth{Genesis 30:29}{Genesis 30:29}\vs{29} ``You know how I have worked for you,'' Jacob replied, ``and how well your livestock have fared under my care.
\markboth{Genesis 30:30}{Genesis 30:30}\vs{30} Indeed, you had little before I arrived, but now your possessions have increased many times over. The Lord has blessed you wherever I worked. But now, how long must it be before I do something for my own family too?''\gva{h}\marginnote{\gva{h}\,The order of nature requires that every one provide for his own family.}
\everypar{}

\parshape=0
\markboth{Genesis 30:31}{Genesis 30:31}\vs{31} So Laban asked, ``What should I give you?'' ``You don't need to give me a thing,'' Jacob replied, ``but if you agree to this one condition, I will continue to care for your flocks and protect them:
\markboth{Genesis 30:32}{Genesis 30:32}\vs{32} Let me walk among all your flocks today and remove from them every speckled or spotted sheep, every dark-colored lamb, and the spotted or speckled goats. These animals will be my wages.\gva{i}\marginnote{\gva{i}\,That which is spotted, from now on.}
\markboth{Genesis 30:33}{Genesis 30:33}\vs{33} My integrity will testify for me later on. When you come to verify that I've taken only the wages we agreed on, if I have in my possession any goat that is not speckled or spotted or any sheep that is not dark-colored, it will be considered stolen.''\gva{k}\marginnote{\gva{k}\,God shall attest to my righteous dealing by rewarding my labours.}
\markboth{Genesis 30:34}{Genesis 30:34}\vs{34} ``Agreed!'' said Laban, ``It will be as you say.''
\everypar{}

\parshape=0
\markboth{Genesis 30:35}{Genesis 30:35}\vs{35} So that day Laban removed the male goats that were streaked or spotted, all the female goats that were speckled or spotted (all that had any white on them), and all the dark-colored lambs, and put them in the care of his sons.
\markboth{Genesis 30:36}{Genesis 30:36}\vs{36} Then he separated them from Jacob by a three-day journey, while Jacob was taking care of the rest of Laban's flocks.
\everypar{}

\parshape=0
\markboth{Genesis 30:37}{Genesis 30:37}\vs{37} But Jacob took fresh-cut branches from poplar, almond, and plane trees. He made white streaks by peeling them, making the white inner wood in the branches visible.\gva{l}\marginnote{\gva{l}\,Jacob used no deceit in this for it was God's commandment as he declares in Genesis 31:9 ; Genesis 31:11 .}
\markboth{Genesis 30:38}{Genesis 30:38}\vs{38} Then he set up the peeled branches in all the watering troughs where the flocks came to drink. He set up the branches in front of the flocks when they were in heat and came to drink.
\markboth{Genesis 30:39}{Genesis 30:39}\vs{39} When the sheep mated in front of the branches, they gave birth to young that were streaked or speckled or spotted.
\markboth{Genesis 30:40}{Genesis 30:40}\vs{40} Jacob removed these lambs, but he made the rest of the flock face the streaked and completely dark-colored animals in Laban's flock. So he made separate flocks for himself and did not mix them with Laban's flocks.
\markboth{Genesis 30:41}{Genesis 30:41}\vs{41} When the stronger females were in heat, Jacob would set up the branches in the troughs in front of the flock, so they would mate near the branches.\gva{m}\marginnote{\gva{m}\,As they which took the ram about September and brought forth about March: so the feebler in March and lamb in September.}
\markboth{Genesis 30:42}{Genesis 30:42}\vs{42} But if the animals were weaker, he did not set the branches there. So the weaker animals ended up belonging to Laban and the stronger animals to Jacob.
\markboth{Genesis 30:43}{Genesis 30:43}\vs{43} In this way Jacob became extremely prosperous. He owned large flocks, male and female servants, camels, and donkeys.

\markboth{Genesis 31:1}{Genesis 31:1}
\Needspace*{8\baselineskip}
\ch{31} \allowchapbreak\hypertarget{ch-genesis-31}{}\lettrine[lines=5]{\color{pentateuch}J}{\textsc{acob}} heard that Laban's sons were complaining, ``Jacob has taken everything that belonged to our father! He has gotten rich at our father's expense!''\gva{a}\marginnote{\gva{a}\,The children put in words what the father disguised in his heart for the covetous think that whatever they cannot take, is taken from them.}\everypar{}
\markboth{Genesis 31:2}{Genesis 31:2}\vs{2} When Jacob saw the look on Laban's face, he could tell his attitude toward him had changed.
\\\indent\markboth{Genesis 31:3}{Genesis 31:3}\vs{3} The Lord said to Jacob, ``Return to the land of your fathers and to your relatives. I will be with you.''
\markboth{Genesis 31:4}{Genesis 31:4}\vs{4} So Jacob sent a message for Rachel and Leah to come to the field where his flocks were.
\markboth{Genesis 31:5}{Genesis 31:5}\vs{5} There he said to them, ``I can tell that your father's attitude toward me has changed, but the God of my father has been with me.\gva{b}\marginnote{\gva{b}\,The God whom my fathers worshipped.}
\markboth{Genesis 31:6}{Genesis 31:6}\vs{6} You know that I've worked for your father as hard as I could,
\markboth{Genesis 31:7}{Genesis 31:7}\vs{7} but your father has humiliated me and changed my wages ten times. But God has not permitted him to do me any harm.
\markboth{Genesis 31:8}{Genesis 31:8}\vs{8} If he said, `The speckled animals will be your wage,' then the entire flock gave birth to speckled offspring. But if he said, `The streaked animals will be your wage,' then the entire flock gave birth to streaked offspring.
\markboth{Genesis 31:9}{Genesis 31:9}\vs{9} In this way God has snatched away your father's livestock and given them to me.\gva{c}\marginnote{\gva{c}\,This declares that the thing Jacob did before, was by God's commandment, and not through deceit.}
\everypar{}

\parshape=0
\markboth{Genesis 31:10}{Genesis 31:10}\vs{10} ``Once during breeding season I saw in a dream that the male goats mating with the flock were streaked, speckled, and spotted.
\markboth{Genesis 31:11}{Genesis 31:11}\vs{11} In the dream the angel of God said to me, `Jacob!' `Here I am!' I replied.
\markboth{Genesis 31:12}{Genesis 31:12}\vs{12} Then he said, `Observe that all the male goats mating with the flock are streaked, speckled, or spotted, for I have observed all that Laban has done to you.
\markboth{Genesis 31:13}{Genesis 31:13}\vs{13} I am the God of Bethel, where you anointed the sacred stone and made a vow to me. Now leave this land immediately and return to your native land.'''\gva{d}\marginnote{\gva{d}\,This angel was Christ who appeared to Jacob in Bethel: and by this it appears that he had taught his wives the fear of God: for he talks as though they knew this thing.}
\everypar{}

\parshape=0
\markboth{Genesis 31:14}{Genesis 31:14}\vs{14} Then Rachel and Leah replied to him, ``Do we still have any portion or inheritance in our father's house?
\markboth{Genesis 31:15}{Genesis 31:15}\vs{15} Hasn't he treated us like foreigners? He not only sold us, but completely wasted the money paid for us!\gva{e}\marginnote{\gva{e}\,For they were given to Jacob as payment for his service, which was a kind of sale.}
\markboth{Genesis 31:16}{Genesis 31:16}\vs{16} Surely all the wealth that God snatched away from our father belongs to us and to our children. So now do everything God has told you.''
\everypar{}

\parshape=0
\markboth{Genesis 31:17}{Genesis 31:17}\vs{17} So Jacob immediately put his children and his wives on the camels.
\markboth{Genesis 31:18}{Genesis 31:18}\vs{18} He took away all the livestock he had acquired in Paddan Aram and all his moveable property that he had accumulated. Then he set out toward the land of Canaan to return to his father Isaac.
\everypar{}

\parshape=0
\markboth{Genesis 31:19}{Genesis 31:19}\vs{19} While Laban had gone to shear his sheep, Rachel stole the household idols that belonged to her father.\gva{f}\marginnote{\gva{f}\,For so the word here signifies, because Laban calls them gods, Genesis 31:30 .}
\markboth{Genesis 31:20}{Genesis 31:20}\vs{20} Jacob also deceived Laban the Aramean by not telling him that he was leaving.
\markboth{Genesis 31:21}{Genesis 31:21}\vs{21} He left with all he owned. He quickly crossed the Euphrates River and headed for the hill country of Gilead.
\everypar{}

\parshape=0
\markboth{Genesis 31:22}{Genesis 31:22}\vs{22} Three days later Laban discovered Jacob had left.
\markboth{Genesis 31:23}{Genesis 31:23}\vs{23} So he took his relatives with him and pursued Jacob for seven days. He caught up with him in the hill country of Gilead.
\markboth{Genesis 31:24}{Genesis 31:24}\vs{24} But God came to Laban the Aramean in a dream at night and warned him, ``Be careful that you neither bless nor curse Jacob.''
\everypar{}

\parshape=0
\markboth{Genesis 31:25}{Genesis 31:25}\vs{25} Laban overtook Jacob, and when Jacob pitched his tent in the hill country of Gilead, Laban and his relatives set up camp there too.
\markboth{Genesis 31:26}{Genesis 31:26}\vs{26} ``What have you done?'' Laban demanded of Jacob. ``You've deceived me and carried away my daughters as if they were captives of war!
\markboth{Genesis 31:27}{Genesis 31:27}\vs{27} Why did you run away secretly and deceive me? Why didn't you tell me so I could send you off with a celebration complete with singing, tambourines, and harps?
\markboth{Genesis 31:28}{Genesis 31:28}\vs{28} You didn't even allow me to kiss my daughters and my grandchildren goodbye. You have acted foolishly!
\markboth{Genesis 31:29}{Genesis 31:29}\vs{29} I have the power to do you harm, but the God of your father told me last night, `Be careful that you neither bless nor curse Jacob.'\gva{g}\marginnote{\gva{g}\,He was an idolater and therefore would not acknowledge the God of Jacob for his God.}
\markboth{Genesis 31:30}{Genesis 31:30}\vs{30} Now I understand that you have gone away because you longed desperately for your father's house. Yet why did you steal my gods?''
\everypar{}

\parshape=0
\markboth{Genesis 31:31}{Genesis 31:31}\vs{31} ``I left secretly because I was afraid!'' Jacob replied to Laban. ``I thought you might take your daughters away from me by force.
\markboth{Genesis 31:32}{Genesis 31:32}\vs{32} Whoever has taken your gods will be put to death! In the presence of our relatives identify whatever is yours and take it.'' (Now Jacob did not know that Rachel had stolen them.)
\everypar{}

\parshape=0
\markboth{Genesis 31:33}{Genesis 31:33}\vs{33} So Laban entered Jacob's tent, and Leah's tent, and the tent of the two female servants, but he did not find the idols. Then he left Leah's tent and entered Rachel's.
\markboth{Genesis 31:34}{Genesis 31:34}\vs{34} (Now Rachel had taken the idols and put them inside her camel's saddle and sat on them.) Laban searched the whole tent, but did not find them.
\markboth{Genesis 31:35}{Genesis 31:35}\vs{35} Rachel said to her father, ``Don't be angry, my lord. I cannot stand up in your presence because I am having my period.'' So he searched thoroughly, but did not find the idols.
\everypar{}

\parshape=0
\markboth{Genesis 31:36}{Genesis 31:36}\vs{36} Jacob became angry and argued with Laban. ``What did I do wrong?'' he demanded of Laban. ``What sin of mine prompted you to chase after me in hot pursuit?
\markboth{Genesis 31:37}{Genesis 31:37}\vs{37} When you searched through all my goods, did you find anything that belonged to you? Set it here before my relatives and yours, and let them settle the dispute between the two of us!
\everypar{}

\parshape=0
\markboth{Genesis 31:38}{Genesis 31:38}\vs{38} ``I have been with you for the past 20 years. Your ewes and female goats have not miscarried, nor have I eaten rams from your flocks.
\markboth{Genesis 31:39}{Genesis 31:39}\vs{39} Animals torn by wild beasts I never brought to you; I always absorbed the loss myself. You always made me pay for every missing animal, whether it was taken by day or at night.
\markboth{Genesis 31:40}{Genesis 31:40}\vs{40} I was consumed by scorching heat during the day and by piercing cold at night, and I went without sleep.
\markboth{Genesis 31:41}{Genesis 31:41}\vs{41} This was my lot for 20 years in your house: I worked like a slave for you---14 years for your two daughters and 6 years for your flocks---but you changed my wages 10 times!
\markboth{Genesis 31:42}{Genesis 31:42}\vs{42} If the God of my father---the God of Abraham, the one whom Isaac fears---had not been with me, you would certainly have sent me away empty-handed! But God saw how I was oppressed and how hard I worked, and he rebuked you last night.''\gva{h}\marginnote{\gva{h}\,That is, the God whom Isaac feared and reverenced.}
\everypar{}

\parshape=0
\markboth{Genesis 31:43}{Genesis 31:43}\vs{43} Laban replied to Jacob, ``These women are my daughters, these children are my grandchildren, and these flocks are my flocks. All that you see belongs to me. But how can I harm these daughters of mine today or the children to whom they have given birth?
\markboth{Genesis 31:44}{Genesis 31:44}\vs{44} So now, come, let's make a formal agreement, you and I, and it will be proof that we have made peace.''\gva{i}\marginnote{\gva{i}\,His conscience reproved him for his misbehaviour toward Jacob, and therefore moved him to seek peace.}
\everypar{}

\parshape=0
\markboth{Genesis 31:45}{Genesis 31:45}\vs{45} So Jacob took a stone and set it up as a memorial pillar.
\markboth{Genesis 31:46}{Genesis 31:46}\vs{46} Then he said to his relatives, ``Gather stones.'' So they brought stones and put them in a pile. They ate there by the pile of stones.
\markboth{Genesis 31:47}{Genesis 31:47}\vs{47} Laban called it Jegar Sahadutha, but Jacob called it Galeed.\gva{k}\marginnote{\gva{k}\,The one named the place in the Syrian tongue, and the other in the Hebrew tongue.}
\everypar{}

\parshape=0
\markboth{Genesis 31:48}{Genesis 31:48}\vs{48} Laban said, ``This pile of stones is a witness of our agreement today.'' That is why it was called Galeed.
\markboth{Genesis 31:49}{Genesis 31:49}\vs{49} It was also called Mizpah because he said, ``May the Lord watch between us when we are out of sight of one another.\gva{l}\marginnote{\gva{l}\,To punish the trespasser.}
\markboth{Genesis 31:50}{Genesis 31:50}\vs{50} If you mistreat my daughters or if you take wives besides my daughters, although no one else is with us, realize that God is witness to your actions.''\gva{m}\marginnote{\gva{m}\,Nature compels him to condemn that vice, to which through covetousness he forced Jacob.}
\everypar{}

\parshape=0
\markboth{Genesis 31:51}{Genesis 31:51}\vs{51} ``Here is this pile of stones and this pillar I have set up between me and you,'' Laban said to Jacob.
\markboth{Genesis 31:52}{Genesis 31:52}\vs{52} ``This pile of stones and the pillar are reminders that I will not pass beyond this pile to come to harm you and that you will not pass beyond this pile and this pillar to come to harm me.
\markboth{Genesis 31:53}{Genesis 31:53}\vs{53} May the God of Abraham and the god of Nahor, the gods of their father, judge between us.'' Jacob took an oath by the God whom his father Isaac feared.\gva{n}\marginnote{\gva{n}\,Behold, how the idolaters mingle the true God with their false gods.}\gva{o}\marginnote{\gva{o}\,Meaning, by the true God whom Isaac worshipped.}
\markboth{Genesis 31:54}{Genesis 31:54}\vs{54} Then Jacob offered a sacrifice on the mountain and invited his relatives to eat the meal. They ate the meal and spent the night on the mountain.
\everypar{}

\parshape=0
\markboth{Genesis 31:55}{Genesis 31:55}\vs{55} (32:1) Early in the morning Laban kissed his grandchildren and his daughters goodbye and blessed them. Then Laban left and returned home.\gva{p}\marginnote{\gva{p}\,We see that there is always some seed of the knowledge of God in the hearts of the wicked.}

\markboth{Genesis 32:1}{Genesis 32:1}
\Needspace*{8\baselineskip}
\ch{32} \allowchapbreak\hypertarget{ch-genesis-32}{}\lettrine[lines=5]{\color{pentateuch}S}{\textsc{o}} Jacob went on his way and the angels of God met him.\everypar{}
\markboth{Genesis 32:2}{Genesis 32:2}\vs{2} When Jacob saw them, he exclaimed, ``This is the camp of God!'' So he named that place Mahanaim.\gva{a}\marginnote{\gva{a}\,He acknowledges God's benefits: who for the preservation of his, sends hosts of angels.}
\\\indent\markboth{Genesis 32:3}{Genesis 32:3}\vs{3} Jacob sent messengers on ahead to his brother Esau in the land of Seir, the region of Edom.
\markboth{Genesis 32:4}{Genesis 32:4}\vs{4} He commanded them, ``This is what you must say to my lord Esau: `This is what your servant Jacob says: I have been staying with Laban until now.\gva{b}\marginnote{\gva{b}\,He reverenced his brother in worldly things, because he mainly looked to be preferred to the spiritual promise.}
\markboth{Genesis 32:5}{Genesis 32:5}\vs{5} I have oxen, donkeys, sheep, and male and female servants. I have sent this message to inform my lord, so that I may find favor in your sight.'''
\everypar{}

\parshape=0
\markboth{Genesis 32:6}{Genesis 32:6}\vs{6} The messengers returned to Jacob and said, ``We went to your brother Esau. He is coming to meet you and has 400 men with him.''
\markboth{Genesis 32:7}{Genesis 32:7}\vs{7} Jacob was very afraid and upset. So he divided the people who were with him into two camps, as well as the flocks, herds, and camels.\gva{c}\marginnote{\gva{c}\,Though he was comforted by the angels, yet the infirmity of the flesh appears.}
\markboth{Genesis 32:8}{Genesis 32:8}\vs{8} ``If Esau attacks one camp,'' he thought, ``then the other camp will be able to escape.''
\everypar{}

\parshape=0
\markboth{Genesis 32:9}{Genesis 32:9}\vs{9} Then Jacob prayed, ``O God of my father Abraham, God of my father Isaac, O Lord, you said to me, `Return to your land and to your relatives and I will make you prosper.'
\markboth{Genesis 32:10}{Genesis 32:10}\vs{10} I am not worthy of all the faithful love you have shown your servant. With only my walking stick I crossed the Jordan, but now I have become two camps.\gva{d}\marginnote{\gva{d}\,that is, poor and without all provision.}
\markboth{Genesis 32:11}{Genesis 32:11}\vs{11} Rescue me, I pray, from the hand of my brother Esau, for I am afraid he will come and attack me, as well as the mothers with their children.\gva{e}\marginnote{\gva{e}\,Meaning, he will put all to death. This proverb comes from those who kill the bird together with the young ones.}
\markboth{Genesis 32:12}{Genesis 32:12}\vs{12} But you said, `I will certainly make you prosper and will make your descendants like the sand on the seashore, too numerous to count.'''
\everypar{}

\parshape=0
\markboth{Genesis 32:13}{Genesis 32:13}\vs{13} Jacob stayed there that night. Then he sent as a gift to his brother Esau\gva{f}\marginnote{\gva{f}\,Not distrusting God's assistance, but using such means as God had given him.}
\markboth{Genesis 32:14}{Genesis 32:14}\vs{14} 200 female goats and 20 male goats, 200 ewes and 20 rams,
\markboth{Genesis 32:15}{Genesis 32:15}\vs{15} 30 female camels with their young, 40 cows and 10 bulls, and 20 female donkeys and 10 male donkeys.
\markboth{Genesis 32:16}{Genesis 32:16}\vs{16} He entrusted them to his servants, who divided them into herds. He told his servants, ``Pass over before me, and keep some distance between one herd and the next.''
\markboth{Genesis 32:17}{Genesis 32:17}\vs{17} He instructed the servant leading the first herd, ``When my brother Esau meets you and asks, `To whom do you belong? Where are you going? Whose herds are you driving?'
\markboth{Genesis 32:18}{Genesis 32:18}\vs{18} then you must say, `They belong to your servant Jacob. They have been sent as a gift to my lord Esau. In fact Jacob himself is behind us.'''
\everypar{}

\parshape=0
\markboth{Genesis 32:19}{Genesis 32:19}\vs{19} He also gave these instructions to the second and third servants, as well as all those who were following the herds, saying, ``You must say the same thing to Esau when you meet him.
\markboth{Genesis 32:20}{Genesis 32:20}\vs{20} You must also say, `In fact your servant Jacob is behind us.''' Jacob thought, ``I will first appease him by sending a gift ahead of me. After that I will meet him. Perhaps he will accept me.''\gva{g}\marginnote{\gva{g}\,He thought it no less to depart with these goods with the intent that he might follow the vocation to which God called him.}
\markboth{Genesis 32:21}{Genesis 32:21}\vs{21} So the gifts were sent on ahead of him while he spent that night in the camp.
\everypar{}

\parshape=0
\markboth{Genesis 32:22}{Genesis 32:22}\vs{22} During the night Jacob quickly took his two wives, his two female servants, and his eleven sons and crossed the ford of the Jabbok.
\markboth{Genesis 32:23}{Genesis 32:23}\vs{23} He took them and sent them across the stream along with all his possessions.
\markboth{Genesis 32:24}{Genesis 32:24}\vs{24} So Jacob was left alone. Then a man wrestled with him until daybreak.\gva{h}\marginnote{\gva{h}\,That is, God in the form of a man.}
\markboth{Genesis 32:25}{Genesis 32:25}\vs{25} When the man saw that he could not defeat Jacob, he struck the socket of his hip so the socket of Jacob's hip was dislocated while he wrestled with him.\gva{i}\marginnote{\gva{i}\,For God assails his with the one hand, and upholds them with the other.}
\everypar{}

\parshape=0
\markboth{Genesis 32:26}{Genesis 32:26}\vs{26} Then the man said, ``Let me go, for the dawn is breaking.'' ``I will not let you go,'' Jacob replied, ``unless you bless me.''
\markboth{Genesis 32:27}{Genesis 32:27}\vs{27} The man asked him, ``What is your name?'' He answered, ``Jacob.''
\markboth{Genesis 32:28}{Genesis 32:28}\vs{28} ``No longer will your name be Jacob,'' the man told him, ``but Israel, because you have fought with God and with men and have prevailed.''\gva{k}\marginnote{\gva{k}\,God gave Jacob both power to overcome, and also the praise of the victory.}
\everypar{}

\parshape=0
\markboth{Genesis 32:29}{Genesis 32:29}\vs{29} Then Jacob asked, ``Please tell me your name.'' ``Why do you ask my name?'' the man replied. Then he blessed Jacob there.
\markboth{Genesis 32:30}{Genesis 32:30}\vs{30} So Jacob named the place Peniel, explaining, ``Certainly I have seen God face to face and have survived.''
\everypar{}

\parshape=0
\markboth{Genesis 32:31}{Genesis 32:31}\vs{31} The sun rose over him as he crossed over Penuel, but he was limping because of his hip.\gva{l}\marginnote{\gva{l}\,The faithful to overcome their temptations, so that they feel the pain of it, so they would not boast, except in their humility.}
\markboth{Genesis 32:32}{Genesis 32:32}\vs{32} That is why to this day the Israelites do not eat the sinew which is attached to the socket of the hip, because he struck the socket of Jacob's hip near the attached sinew.

\markboth{Genesis 33:1}{Genesis 33:1}
\Needspace*{8\baselineskip}
\ch{33} \allowchapbreak\hypertarget{ch-genesis-33}{}\lettrine[lines=5]{\color{pentateuch}J}{\textsc{acob}} looked up and saw that Esau was coming along with 400 men. So he divided the children among Leah, Rachel, and the two female servants.\gva{a}\marginnote{\gva{a}\,That if the one part were assailed, the other might escape.}\everypar{}
\markboth{Genesis 33:2}{Genesis 33:2}\vs{2} He put the servants and their children in front, with Leah and her children behind them, and Rachel and Joseph behind them.
\markboth{Genesis 33:3}{Genesis 33:3}\vs{3} But Jacob himself went on ahead of them, and he bowed toward the ground seven times as he approached his brother.\gva{b}\marginnote{\gva{b}\,By this gesture he partly revered his brother and partly prayed to God to appease Esau's wrath.}
\markboth{Genesis 33:4}{Genesis 33:4}\vs{4} But Esau ran to meet him, embraced him, hugged his neck, and kissed him. Then they both wept.
\markboth{Genesis 33:5}{Genesis 33:5}\vs{5} When Esau looked up and saw the women and the children, he asked, ``Who are these people with you?'' Jacob replied, ``The children whom God has graciously given your servant.''
\markboth{Genesis 33:6}{Genesis 33:6}\vs{6} The female servants came forward with their children and bowed down.\gva{c}\marginnote{\gva{c}\,Jacob and his family are the image of the Church under the yoke of tyrants who out of fear are brought to subjection.}
\markboth{Genesis 33:7}{Genesis 33:7}\vs{7} Then Leah came forward with her children and they bowed down. Finally Joseph and Rachel came forward and bowed down.
\everypar{}

\parshape=0
\markboth{Genesis 33:8}{Genesis 33:8}\vs{8} Esau then asked, ``What did you intend by sending all these herds to meet me?'' Jacob replied, ``To find favor in your sight, my lord.''
\markboth{Genesis 33:9}{Genesis 33:9}\vs{9} But Esau said, ``I have plenty, my brother. Keep what belongs to you.''
\markboth{Genesis 33:10}{Genesis 33:10}\vs{10} ``No, please take them,'' Jacob said. ``If I have found favor in your sight, accept my gift from my hand. Now that I have seen your face and you have accepted me, it is as if I have seen the face of God.\gva{d}\marginnote{\gva{d}\,In that his brother embraced him so lovingly, contrary to his expectation, he accepted it as a clear sign of God's presence.}
\markboth{Genesis 33:11}{Genesis 33:11}\vs{11} Please take my present that was brought to you, for God has been generous to me and I have all I need.'' When Jacob urged him, he took it.
\everypar{}

\parshape=0
\markboth{Genesis 33:12}{Genesis 33:12}\vs{12} Then Esau said, ``Let's be on our way! I will go in front of you.''
\markboth{Genesis 33:13}{Genesis 33:13}\vs{13} But Jacob said to him, ``My lord knows that the children are young, and that I have to look after the sheep and cattle that are nursing their young. If they are driven too hard for even a single day, all the animals will die.
\markboth{Genesis 33:14}{Genesis 33:14}\vs{14} Let my lord go on ahead of his servant. I will travel more slowly, at the pace of the herds and the children, until I come to my lord at Seir.''\gva{f}\marginnote{\gva{f}\,He promised that which}
\everypar{}

\parshape=0
\markboth{Genesis 33:15}{Genesis 33:15}\vs{15} So Esau said, ``Let me leave some of my men with you.'' ``Why do that?'' Jacob replied. ``My lord has already been kind enough to me.''
\everypar{}

\parshape=0
\markboth{Genesis 33:16}{Genesis 33:16}\vs{16} So that same day Esau made his way back to Seir.
\markboth{Genesis 33:17}{Genesis 33:17}\vs{17} But Jacob traveled to Sukkoth where he built himself a house and made shelters for his livestock. That is why the place was called Sukkoth.
\everypar{}

\parshape=0
\markboth{Genesis 33:18}{Genesis 33:18}\vs{18} After he left Paddan Aram, Jacob came safely to the city of Shechem in the land of Canaan, and he camped near the city.
\markboth{Genesis 33:19}{Genesis 33:19}\vs{19} Then he purchased the portion of the field where he had pitched his tent; he bought it from the sons of Hamor, Shechem's father, for 100 pieces of money.
\markboth{Genesis 33:20}{Genesis 33:20}\vs{20} There he set up an altar and called it ``The God of Israel is God.''\gva{g}\marginnote{\gva{g}\,He calls the sign, the thing which it signifies, in token that God had mightily delivered him.}

\markboth{Genesis 34:1}{Genesis 34:1}
\Needspace*{8\baselineskip}
\ch{34} \allowchapbreak\hypertarget{ch-genesis-34}{}\lettrine[lines=5]{\color{pentateuch}N}{\textsc{ow}} Dinah, Leah's daughter whom she bore to Jacob, went to meet the young women of the land.\gva{a}\marginnote{\gva{a}\,This example teaches us that too much liberty is not to be given to youth.}\everypar{}
\markboth{Genesis 34:2}{Genesis 34:2}\vs{2} When Shechem son of Hamor the Hivite, who ruled that area, saw her, he grabbed her, forced himself on her, and sexually assaulted her.
\markboth{Genesis 34:3}{Genesis 34:3}\vs{3} Then he became very attached to Dinah, Jacob's daughter. He fell in love with the young woman and spoke romantically to her.
\markboth{Genesis 34:4}{Genesis 34:4}\vs{4} Shechem said to his father Hamor, ``Acquire this young girl as my wife.''\gva{b}\marginnote{\gva{b}\,This proves that the consent of parents is required in marriage, seeing that even the infidels observed it as a necessary thing.}
\markboth{Genesis 34:5}{Genesis 34:5}\vs{5} When Jacob heard that Shechem had violated his daughter Dinah, his sons were with the livestock in the field. So Jacob remained silent until they came in.
\everypar{}

\parshape=0
\markboth{Genesis 34:6}{Genesis 34:6}\vs{6} Then Shechem's father Hamor went to speak with Jacob about Dinah.
\markboth{Genesis 34:7}{Genesis 34:7}\vs{7} Now Jacob's sons had come in from the field when they heard the news. They were offended and very angry because Shechem had disgraced Israel by sexually assaulting Jacob's daughter, a crime that should not be committed.
\everypar{}

\parshape=0
\markboth{Genesis 34:8}{Genesis 34:8}\vs{8} But Hamor made this appeal to them: ``My son Shechem is in love with your daughter. Please give her to him as his wife.
\markboth{Genesis 34:9}{Genesis 34:9}\vs{9} Intermarry with us. Let us marry your daughters, and take our daughters as wives for yourselves.
\markboth{Genesis 34:10}{Genesis 34:10}\vs{10} You may live among us, and the land will be open to you. Live in it, travel freely in it, and acquire property in it.''
\everypar{}

\parshape=0
\markboth{Genesis 34:11}{Genesis 34:11}\vs{11} Then Shechem said to Dinah's father and brothers, ``Let me find favor in your sight, and whatever you require of me I'll give.
\markboth{Genesis 34:12}{Genesis 34:12}\vs{12} You can make the bride price and the gift I must bring very expensive, and I'll give whatever you ask of me. Just give me the young woman as my wife!''
\everypar{}

\parshape=0
\markboth{Genesis 34:13}{Genesis 34:13}\vs{13} Jacob's sons answered Shechem and his father Hamor deceitfully when they spoke because Shechem had violated their sister Dinah.
\markboth{Genesis 34:14}{Genesis 34:14}\vs{14} They said to them, ``We cannot give our sister to a man who is not circumcised, for it would be a disgrace to us.\gva{c}\marginnote{\gva{c}\,They used the holy ordinance of God a means to accomplish their wicked purpose.}\gva{d}\marginnote{\gva{d}\,As it is abomination for those who are baptized to be joined to infidels.}
\markboth{Genesis 34:15}{Genesis 34:15}\vs{15} We will give you our consent on this one condition: You must become like us by circumcising all your males.\gva{e}\marginnote{\gva{e}\,Their fault is even greater since they made religion a disguise for their deceit.}
\markboth{Genesis 34:16}{Genesis 34:16}\vs{16} Then we will give you our daughters to marry, and we will take your daughters as wives for ourselves, and we will live among you and become one people.
\markboth{Genesis 34:17}{Genesis 34:17}\vs{17} But if you do not agree to our terms by being circumcised, then we will take our sister and depart.''
\everypar{}

\parshape=0
\markboth{Genesis 34:18}{Genesis 34:18}\vs{18} Their offer pleased Hamor and his son Shechem.
\markboth{Genesis 34:19}{Genesis 34:19}\vs{19} The young man did not delay in doing what they asked because he wanted Jacob's daughter Dinah badly. (Now he was more important than anyone in his father's household.)
\markboth{Genesis 34:20}{Genesis 34:20}\vs{20} So Hamor and his son Shechem went to the gate of their city and spoke to the men of their city,\gva{f}\marginnote{\gva{f}\,For the people used to assemble there, and justice was administered.}
\markboth{Genesis 34:21}{Genesis 34:21}\vs{21} ``These men are at peace with us. So let them live in the land and travel freely in it, for the land is wide enough for them. We will take their daughters for wives, and we will give them our daughters to marry.\gva{g}\marginnote{\gva{g}\,Thus many pretend to speak for a public profit, when in reality they are only speaking for their own private gain and convenience.}
\markboth{Genesis 34:22}{Genesis 34:22}\vs{22} Only on this one condition will these men consent to live with us and become one people: They demand that every male among us be circumcised just as they are circumcised.
\markboth{Genesis 34:23}{Genesis 34:23}\vs{23} If we do so, won't their livestock, their property, and all their animals become ours? So let's consent to their demand, so they will live among us.''\gva{h}\marginnote{\gva{h}\,Thus they do not lack any form of perversion, who prefer their own convenience before the common good.}
\everypar{}

\parshape=0
\markboth{Genesis 34:24}{Genesis 34:24}\vs{24} All the men who assembled at the city gate agreed with Hamor and his son Shechem. Every male who assembled at the city gate was circumcised.
\markboth{Genesis 34:25}{Genesis 34:25}\vs{25} In three days, when they were still in pain, two of Jacob's sons, Simeon and Levi, Dinah's brothers, each took his sword and went to the unsuspecting city and slaughtered every male.\gva{i}\marginnote{\gva{i}\,For they were the leaders of the company.}\gva{k}\marginnote{\gva{k}\,The people are punished because of their wicked princes.}
\markboth{Genesis 34:26}{Genesis 34:26}\vs{26} They killed Hamor and his son Shechem with the sword, took Dinah from Shechem's house, and left.
\markboth{Genesis 34:27}{Genesis 34:27}\vs{27} Jacob's sons killed them and looted the city because their sister had been violated.
\markboth{Genesis 34:28}{Genesis 34:28}\vs{28} They took their flocks, herds, and donkeys, as well as everything in the city and in the surrounding fields.
\markboth{Genesis 34:29}{Genesis 34:29}\vs{29} They captured as plunder all their wealth, all their little ones, and their wives, including everything in the houses.
\everypar{}

\parshape=0
\markboth{Genesis 34:30}{Genesis 34:30}\vs{30} Then Jacob said to Simeon and Levi, ``You have brought ruin on me by making me a foul odor among the inhabitants of the land---among the Canaanites and the Perizzites. I am few in number; they will join forces against me and attack me, and both I and my family will be destroyed!''
\markboth{Genesis 34:31}{Genesis 34:31}\vs{31} But Simeon and Levi replied, ``Should he treat our sister like a common prostitute?''

\markboth{Genesis 35:1}{Genesis 35:1}
\Needspace*{8\baselineskip}
\ch{35} \allowchapbreak\hypertarget{ch-genesis-35}{}\lettrine[lines=5]{\color{pentateuch}T}{\textsc{hen}} God said to Jacob, ``Go up at once to Bethel and live there. Make an altar there to God, who appeared to you when you fled from your brother Esau.''\gva{a}\marginnote{\gva{a}\,God is ever at hand to comfort his people in their troubles.}\everypar{}
\markboth{Genesis 35:2}{Genesis 35:2}\vs{2} So Jacob told his household and all who were with him, ``Get rid of the foreign gods you have among you. Purify yourselves and change your clothes.\gva{b}\marginnote{\gva{b}\,That by this outward act they should show their inward repentance.}
\markboth{Genesis 35:3}{Genesis 35:3}\vs{3} Let us go up at once to Bethel. Then I will make an altar there to God, who responded to me in my time of distress and has been with me wherever I went.''
\everypar{}

\parshape=0
\markboth{Genesis 35:4}{Genesis 35:4}\vs{4} So they gave Jacob all the foreign gods that were in their possession and the rings that were in their ears. Jacob buried them under the oak near Shechem\gva{e}\marginnote{\gva{e}\,For in this was some sign of superstition, as in tablets and Agnus deis}
\markboth{Genesis 35:5}{Genesis 35:5}\vs{5} and they started on their journey. The surrounding cities were afraid of God, and they did not pursue the sons of Jacob.\gva{d}\marginnote{\gva{d}\,terror of God was upon the cities that [were] round about them, and they did not pursue after the sons of Jacob.}\gva{d}\marginnote{\gva{d}\,Thus, despite the inconvenience that came before, God delivered Jacob.}
\everypar{}

\parshape=0
\markboth{Genesis 35:6}{Genesis 35:6}\vs{6} Jacob and all those who were with him arrived at Luz (that is, Bethel) in the land of Canaan.
\markboth{Genesis 35:7}{Genesis 35:7}\vs{7} He built an altar there and named the place El Bethel because there God had revealed himself to him when he was fleeing from his brother.
\markboth{Genesis 35:8}{Genesis 35:8}\vs{8} (Deborah, Rebekah's nurse, died and was buried under the oak below Bethel; thus it was named Oak of Weeping.)
\everypar{}

\parshape=0
\markboth{Genesis 35:9}{Genesis 35:9}\vs{9} God appeared to Jacob again after he returned from Paddan Aram and blessed him.
\markboth{Genesis 35:10}{Genesis 35:10}\vs{10} God said to him, ``Your name is Jacob, but your name will no longer be called Jacob; Israel will be your name.'' So God named him Israel.
\markboth{Genesis 35:11}{Genesis 35:11}\vs{11} Then God said to him, ``I am the Sovereign God. Be fruitful and multiply! A nation---even a company of nations---will descend from you; kings will be among your descendants!
\markboth{Genesis 35:12}{Genesis 35:12}\vs{12} The land I gave to Abraham and Isaac I will give to you. To your descendants I will also give this land.''
\markboth{Genesis 35:13}{Genesis 35:13}\vs{13} Then God went up from the place where he spoke with him.\gva{e}\marginnote{\gva{e}\,As God is said to descend, when he shows some sign of his presence: so he is said to ascend when a vision is ended.}
\markboth{Genesis 35:14}{Genesis 35:14}\vs{14} So Jacob set up a sacred stone pillar in the place where God spoke with him. He poured out a drink offering on it, and then he poured oil on it.
\markboth{Genesis 35:15}{Genesis 35:15}\vs{15} Jacob named the place where God spoke with him Bethel.
\everypar{}

\parshape=0
\markboth{Genesis 35:16}{Genesis 35:16}\vs{16} They traveled on from Bethel, and when Ephrath was still some distance away, Rachel went into labor---and her labor was hard.\gva{f}\marginnote{\gva{f}\,The Hebrew word signifies as much ground as one can cover from resting point to resting point, which is taken for half a days journey.}
\markboth{Genesis 35:17}{Genesis 35:17}\vs{17} When her labor was at its hardest, the midwife said to her, ``Don't be afraid, for you are having another son.''
\markboth{Genesis 35:18}{Genesis 35:18}\vs{18} With her dying breath, she named him Ben Oni. But his father called him Benjamin instead.
\markboth{Genesis 35:19}{Genesis 35:19}\vs{19} So Rachel died and was buried on the way to Ephrath (that is, Bethlehem).
\markboth{Genesis 35:20}{Genesis 35:20}\vs{20} Jacob set up a marker over her grave; it is the Marker of Rachel's Grave to this day.\gva{g}\marginnote{\gva{g}\,The ancient fathers used this ceremony to testify their hope of the resurrection to come, which was not generally revealed.}
\everypar{}

\parshape=0
\markboth{Genesis 35:21}{Genesis 35:21}\vs{21} Then Israel traveled on and pitched his tent beyond Migdal Eder.
\everypar{}

\parshape=0
\markboth{Genesis 35:22}{Genesis 35:22}\vs{22} While Israel was living in that land, Reuben went to bed with Bilhah, his father's concubine, and Israel heard about it. Jacob had twelve sons:\gva{h}\marginnote{\gva{h}\,This teaches that the fathers were not chosen for their merits, but only by God's mercies, whose election was not changed by their faults.}
\everypar{}

\parshape=0
\markboth{Genesis 35:23}{Genesis 35:23}\vs{23} The sons of Leah were Reuben, Jacob's firstborn, as well as Simeon, Levi, Judah, Issachar, and Zebulun.
\everypar{}

\parshape=0
\markboth{Genesis 35:24}{Genesis 35:24}\vs{24} The sons of Rachel were Joseph and Benjamin.
\everypar{}

\parshape=0
\markboth{Genesis 35:25}{Genesis 35:25}\vs{25} The sons of Bilhah, Rachel's servant, were Dan and Naphtali.
\everypar{}

\parshape=0
\markboth{Genesis 35:26}{Genesis 35:26}\vs{26} The sons of Zilpah, Leah's servant, were Gad and Asher. These were the sons of Jacob who were born to him in Paddan Aram.
\everypar{}

\parshape=0
\markboth{Genesis 35:27}{Genesis 35:27}\vs{27} So Jacob came back to his father Isaac in Mamre, to Kiriath Arba (that is, Hebron), where Abraham and Isaac had stayed.
\markboth{Genesis 35:28}{Genesis 35:28}\vs{28} Isaac lived to be 180 years old.
\markboth{Genesis 35:29}{Genesis 35:29}\vs{29} Then Isaac breathed his last and joined his ancestors. He died an old man who had lived a full life. His sons Esau and Jacob buried him.

\markboth{Genesis 36:1}{Genesis 36:1}
\Needspace*{8\baselineskip}
\ch{36} \allowchapbreak\hypertarget{ch-genesis-36}{}\lettrine[lines=5]{\color{pentateuch}W}{\textsc{hat}} follows is the account of Esau (also known as Edom).\gva{a}\marginnote{\gva{a}\,This genealogy declares that Esau was blessed physically and that his father's blessing took place in worldly things.}\everypar{}
\\\indent\markboth{Genesis 36:2}{Genesis 36:2}\vs{2} Esau took his wives from the Canaanites: Adah the daughter of Elon the Hittite, and Oholibamah the daughter of Anah and granddaughter of Zibeon the Hivite,\gva{b}\marginnote{\gva{b}\,Besides those wives spoken of in Genesis 26:34 .}
\markboth{Genesis 36:3}{Genesis 36:3}\vs{3} in addition to Basemath the daughter of Ishmael and sister of Nebaioth.
\\\indent\markboth{Genesis 36:4}{Genesis 36:4}\vs{4} Adah bore Eliphaz to Esau, Basemath bore Reuel,
\markboth{Genesis 36:5}{Genesis 36:5}\vs{5} and Oholibamah bore Jeush, Jalam, and Korah. These were the sons of Esau who were born to him in the land of Canaan.
\everypar{}

\parshape=0
\markboth{Genesis 36:6}{Genesis 36:6}\vs{6} Esau took his wives, his sons, his daughters, all the people in his household, his livestock, his animals, and all his possessions that he had acquired in the land of Canaan, and he went to a land some distance away from Jacob his brother\gva{c}\marginnote{\gva{c}\,In this, God's providence appears, which causes the wicked to give place to the godly, that Jacob might enjoy Canaan according to God's promise.}
\markboth{Genesis 36:7}{Genesis 36:7}\vs{7} because they had too many possessions to be able to stay together, and the land where they had settled was not able to support them because of their livestock.
\markboth{Genesis 36:8}{Genesis 36:8}\vs{8} So Esau (also known as Edom) lived in the hill country of Seir.
\everypar{}

\parshape=0
\markboth{Genesis 36:9}{Genesis 36:9}\vs{9} This is the account of Esau, the father of the Edomites, in the hill country of Seir.
\everypar{}

\parshape=0
\markboth{Genesis 36:10}{Genesis 36:10}\vs{10} These were the names of Esau's sons: Eliphaz, the son of Esau's wife Adah, and Reuel, the son of Esau's wife Basemath.
\everypar{}

\parshape=0
\markboth{Genesis 36:11}{Genesis 36:11}\vs{11} These were the sons of Eliphaz: Teman, Omar, Zepho, Gatam, and Kenaz.
\everypar{}

\parshape=0
\markboth{Genesis 36:12}{Genesis 36:12}\vs{12} Timna, a concubine of Esau's son Eliphaz, bore Amalek to Eliphaz. These were the sons of Esau's wife Adah.
\everypar{}

\parshape=0
\markboth{Genesis 36:13}{Genesis 36:13}\vs{13} These were the sons of Reuel: Nahath, Zerah, Shammah, and Mizzah. These were the sons of Esau's wife Basemath.
\everypar{}

\parshape=0
\markboth{Genesis 36:14}{Genesis 36:14}\vs{14} These were the sons of Esau's wife Oholibamah the daughter of Anah and granddaughter of Zibeon: She bore Jeush, Jalam, and Korah to Esau.
\everypar{}

\parshape=0
\markboth{Genesis 36:15}{Genesis 36:15}\vs{15} These were the chiefs among the descendants of Esau, the sons of Eliphaz, Esau's firstborn: chief Teman, chief Omar, chief Zepho, chief Kenaz,\gva{d}\marginnote{\gva{d}\,If God's promises are so sure towards those who are not of his household, how much more will he perform the same for us?}
\markboth{Genesis 36:16}{Genesis 36:16}\vs{16} chief Korah, chief Gatam, chief Amalek. These were the chiefs descended from Eliphaz in the land of Edom; these were the sons of Adah.
\everypar{}

\parshape=0
\markboth{Genesis 36:17}{Genesis 36:17}\vs{17} These were the sons of Esau's son Reuel: chief Nahath, chief Zerah, chief Shammah, chief Mizzah. These were the chiefs descended from Reuel in the land of Edom; these were the sons of Esau's wife Basemath.
\everypar{}

\parshape=0
\markboth{Genesis 36:18}{Genesis 36:18}\vs{18} These were the sons of Esau's wife Oholibamah: chief Jeush, chief Jalam, chief Korah. These were the chiefs descended from Esau's wife Oholibamah, the daughter of Anah.
\everypar{}

\parshape=0
\markboth{Genesis 36:19}{Genesis 36:19}\vs{19} These were the sons of Esau (also known as Edom), and these were their chiefs.
\everypar{}

\parshape=0
\markboth{Genesis 36:20}{Genesis 36:20}\vs{20} These were the sons of Seir the Horite, who were living in the land: Lotan, Shobal, Zibeon, Anah,\gva{e}\marginnote{\gva{e}\,Esau lived there before that.}
\markboth{Genesis 36:21}{Genesis 36:21}\vs{21} Dishon, Ezer, and Dishan. These were the chiefs of the Horites, the descendants of Seir in the land of Edom.
\everypar{}

\parshape=0
\markboth{Genesis 36:22}{Genesis 36:22}\vs{22} The sons of Lotan were Hori and Homam; Lotan's sister was Timna.
\everypar{}

\parshape=0
\markboth{Genesis 36:23}{Genesis 36:23}\vs{23} These were the sons of Shobal: Alvan, Manahath, Ebal, Shepho, and Onam.
\everypar{}

\parshape=0
\markboth{Genesis 36:24}{Genesis 36:24}\vs{24} These were the sons of Zibeon: Aiah and Anah (who discovered the hot springs in the wilderness as he pastured the donkeys of his father Zibeon).\gva{f}\marginnote{\gva{f}\,Who not contented with those kinds of beasts, which God had created, discovered the monstrous generation of mules between the ass and the mare.}
\everypar{}

\parshape=0
\markboth{Genesis 36:25}{Genesis 36:25}\vs{25} These were the children of Anah: Dishon and Oholibamah, the daughter of Anah.
\everypar{}

\parshape=0
\markboth{Genesis 36:26}{Genesis 36:26}\vs{26} These were the sons of Dishon: Hemdan, Eshban, Ithran, and Keran.
\everypar{}

\parshape=0
\markboth{Genesis 36:27}{Genesis 36:27}\vs{27} These were the sons of Ezer: Bilhan, Zaavan, and Akan.
\everypar{}

\parshape=0
\markboth{Genesis 36:28}{Genesis 36:28}\vs{28} These were the sons of Dishan: Uz and Aran.
\everypar{}

\parshape=0
\markboth{Genesis 36:29}{Genesis 36:29}\vs{29} These were the chiefs of the Horites: chief Lotan, chief Shobal, chief Zibeon, chief Anah,
\markboth{Genesis 36:30}{Genesis 36:30}\vs{30} chief Dishon, chief Ezer, chief Dishan. These were the chiefs of the Horites, according to their chief lists in the land of Seir.
\everypar{}

\parshape=0
\markboth{Genesis 36:31}{Genesis 36:31}\vs{31} These were the kings who reigned in the land of Edom before any king ruled over the Israelites:\gva{g}\marginnote{\gva{g}\,The wicked rise up suddenly to honour and perish as quickly: but the inheritance of the children of God continues forever, Psalms 102:28 .}
\everypar{}

\parshape=0
\markboth{Genesis 36:32}{Genesis 36:32}\vs{32} Bela the son of Beor reigned in Edom; the name of his city was Dinhabah.
\everypar{}

\parshape=0
\markboth{Genesis 36:33}{Genesis 36:33}\vs{33} When Bela died, Jobab the son of Zerah from Bozrah reigned in his place.
\everypar{}

\parshape=0
\markboth{Genesis 36:34}{Genesis 36:34}\vs{34} When Jobab died, Husham from the land of the Temanites reigned in his place.
\everypar{}

\parshape=0
\markboth{Genesis 36:35}{Genesis 36:35}\vs{35} When Husham died, Hadad the son of Bedad, who defeated the Midianites in the land of Moab, reigned in his place; the name of his city was Avith.
\everypar{}

\parshape=0
\markboth{Genesis 36:36}{Genesis 36:36}\vs{36} When Hadad died, Samlah from Masrekah reigned in his place.
\everypar{}

\parshape=0
\markboth{Genesis 36:37}{Genesis 36:37}\vs{37} When Samlah died, Shaul from Rehoboth on the River reigned in his place.\gva{h}\marginnote{\gva{h}\,Which is by the river Euphrates.}
\everypar{}

\parshape=0
\markboth{Genesis 36:38}{Genesis 36:38}\vs{38} When Shaul died, Baal Hanan the son of Achbor reigned in his place.
\everypar{}

\parshape=0
\markboth{Genesis 36:39}{Genesis 36:39}\vs{39} When Baal Hanan the son of Achbor died, Hadad reigned in his place; the name of his city was Pau. His wife's name was Mehetabel, the daughter of Matred, the daughter of Me-Zahab.
\everypar{}

\parshape=0
\markboth{Genesis 36:40}{Genesis 36:40}\vs{40} These were the names of the chiefs of Esau, according to their families, according to their places, by their names: chief Timna, chief Alvah, chief Jetheth,
\markboth{Genesis 36:41}{Genesis 36:41}\vs{41} chief Oholibamah, chief Elah, chief Pinon,
\markboth{Genesis 36:42}{Genesis 36:42}\vs{42} chief Kenaz, chief Teman, chief Mibzar,
\markboth{Genesis 36:43}{Genesis 36:43}\vs{43} chief Magdiel, chief Iram. These were the chiefs of Edom, according to their settlements in the land they possessed. This was Esau, the father of the Edomites.\gva{i}\marginnote{\gva{i}\,Of Edom came the Idumeans.}

\markboth{Genesis 37:1}{Genesis 37:1}
\Needspace*{8\baselineskip}
\ch{37} \allowchapbreak\hypertarget{ch-genesis-37}{}\lettrine[lines=5]{\color{pentateuch}B}{\textsc{ut}} Jacob lived in the land where his father had stayed, in the land of Canaan.\gva{a}\marginnote{\gva{a}\,That is, the story of such things as came to him and his family as in Genesis 5:1}\everypar{}
\\\indent\markboth{Genesis 37:2}{Genesis 37:2}\vs{2} This is the account of Jacob. Joseph, his seventeen-year-old son, was taking care of the flocks with his brothers. Now he was a youngster working with the sons of Bilhah and Zilpah, his father's wives. Joseph brought back a bad report about them to their father.\gva{b}\marginnote{\gva{b}\,He complained of the evil words and injuries which they spoke and did to him.}
\everypar{}

\parshape=0
\markboth{Genesis 37:3}{Genesis 37:3}\vs{3} Now Israel loved Joseph more than all his sons because he was a son born to him late in life, and he made a special tunic for him.
\markboth{Genesis 37:4}{Genesis 37:4}\vs{4} When Joseph's brothers saw that their father loved him more than any of them, they hated Joseph and were not able to speak to him kindly.
\everypar{}

\parshape=0
\markboth{Genesis 37:5}{Genesis 37:5}\vs{5} Joseph had a dream, and when he told his brothers about it they hated him even more.\gva{c}\marginnote{\gva{c}\,God revealed to him by a dream what should come to pass.}
\markboth{Genesis 37:6}{Genesis 37:6}\vs{6} He said to them, ``Listen to this dream I had:
\markboth{Genesis 37:7}{Genesis 37:7}\vs{7} There we were, binding sheaves of grain in the middle of the field. Suddenly my sheaf rose up and stood upright and your sheaves surrounded my sheaf and bowed down to it!''
\markboth{Genesis 37:8}{Genesis 37:8}\vs{8} Then his brothers asked him, ``Do you really think you will rule over us or have dominion over us?'' They hated him even more because of his dream and because of what he said.\gva{d}\marginnote{\gva{d}\,The more God shows himself favourable to his own, the more the malice of the wicked rages against them.}
\everypar{}

\parshape=0
\markboth{Genesis 37:9}{Genesis 37:9}\vs{9} Then he had another dream, and told it to his brothers. ``Look,'' he said. ``I had another dream. The sun, the moon, and eleven stars were bowing down to me.''
\markboth{Genesis 37:10}{Genesis 37:10}\vs{10} When he told his father and his brothers, his father rebuked him, saying, ``What is this dream that you had? Will I, your mother, and your brothers really come and bow down to you?''\gva{e}\marginnote{\gva{e}\,Not despising the vision, but seeking to appease his brethren.}
\markboth{Genesis 37:11}{Genesis 37:11}\vs{11} His brothers were jealous of him, but his father kept in mind what Joseph said.\gva{f}\marginnote{\gva{f}\,He knew that God was the author of the dream, but he did not understand the meaning.}
\everypar{}

\parshape=0
\markboth{Genesis 37:12}{Genesis 37:12}\vs{12} When his brothers had gone to graze their father's flocks near Shechem,
\markboth{Genesis 37:13}{Genesis 37:13}\vs{13} Israel said to Joseph, ``Your brothers are grazing the flocks near Shechem. Come, I will send you to them.'' ``I'm ready,'' Joseph replied.
\markboth{Genesis 37:14}{Genesis 37:14}\vs{14} So Jacob said to him, ``Go now and check on the welfare of your brothers and of the flocks, and bring me word.'' So Jacob sent him from the valley of Hebron.
\everypar{}

\parshape=0
\markboth{Genesis 37:15}{Genesis 37:15}\vs{15} When Joseph reached Shechem, a man found him wandering in the field, so the man asked him, ``What are you looking for?''
\markboth{Genesis 37:16}{Genesis 37:16}\vs{16} He replied, ``I'm looking for my brothers. Please tell me where they are grazing their flocks.''
\markboth{Genesis 37:17}{Genesis 37:17}\vs{17} The man said, ``They left this area, for I heard them say, `Let's go to Dothan.''' So Joseph went after his brothers and found them at Dothan.
\everypar{}

\parshape=0
\markboth{Genesis 37:18}{Genesis 37:18}\vs{18} Now Joseph's brothers saw him from a distance, and before he reached them, they plotted to kill him.\gva{g}\marginnote{\gva{g}\,conspired against him to slay him.}\gva{g}\marginnote{\gva{g}\,The Holy Spirit does not cover the faults of men, as vain writers do, who make virtues out of vices.}
\markboth{Genesis 37:19}{Genesis 37:19}\vs{19} They said to one another, ``Here comes this master of dreams!
\markboth{Genesis 37:20}{Genesis 37:20}\vs{20} Come now, let's kill him, throw him into one of the cisterns, and then say that a wild animal ate him. Then we'll see how his dreams turn out!''
\everypar{}

\parshape=0
\markboth{Genesis 37:21}{Genesis 37:21}\vs{21} When Reuben heard this, he rescued Joseph from their hands, saying, ``Let's not take his life!''
\markboth{Genesis 37:22}{Genesis 37:22}\vs{22} Reuben continued, ``Don't shed blood! Throw him into this cistern that is here in the wilderness, but don't lay a hand on him.'' (Reuben said this so he could rescue Joseph from them and take him back to his father.)
\everypar{}

\parshape=0
\markboth{Genesis 37:23}{Genesis 37:23}\vs{23} When Joseph reached his brothers, they stripped him of his tunic, the special tunic that he wore.
\markboth{Genesis 37:24}{Genesis 37:24}\vs{24} Then they took him and threw him into the cistern. (Now the cistern was empty; there was no water in it.)\gva{h}\marginnote{\gva{h}\,Their hypocrisy appears in this that they feared man more than God: and thought it was not murder, if they did not shed his blood or had excuses to cover their fault.}
\everypar{}

\parshape=0
\markboth{Genesis 37:25}{Genesis 37:25}\vs{25} When they sat down to eat their food, they looked up and saw a caravan of Ishmaelites coming from Gilead. Their camels were carrying spices, balm, and myrrh down to Egypt.
\markboth{Genesis 37:26}{Genesis 37:26}\vs{26} Then Judah said to his brothers, ``What profit is there if we kill our brother and cover up his blood?
\markboth{Genesis 37:27}{Genesis 37:27}\vs{27} Come, let's sell him to the Ishmaelites, but let's not lay a hand on him, for after all, he is our brother, our own flesh.'' His brothers agreed.
\markboth{Genesis 37:28}{Genesis 37:28}\vs{28} So when the Midianite merchants passed by, Joseph's brothers pulled him out of the cistern and sold him to the Ishmaelites for 20 pieces of silver. The Ishmaelites then took Joseph to Egypt.\gva{i}\marginnote{\gva{i}\,Moses writes according to the opinion of those who took the Midianites and Ishmaelites to be one, and here mixes their names: as also appears in Genesis 37:36 ; Genesis 39:1 or else he was first offered to the Midianites, but sold to the Ishmaelites.}
\everypar{}

\parshape=0
\markboth{Genesis 37:29}{Genesis 37:29}\vs{29} Later Reuben returned to the cistern to find that Joseph was not in it! He tore his clothes,
\markboth{Genesis 37:30}{Genesis 37:30}\vs{30} returned to his brothers, and said, ``The boy isn't there! And I, where can I go?''
\markboth{Genesis 37:31}{Genesis 37:31}\vs{31} So they took Joseph's tunic, killed a young goat, and dipped the tunic in the blood.
\markboth{Genesis 37:32}{Genesis 37:32}\vs{32} Then they brought the special tunic to their father and said, ``We found this. Determine now whether it is your son's tunic or not.''\gva{k}\marginnote{\gva{k}\,That is, the messengers who were sent.}
\everypar{}

\parshape=0
\markboth{Genesis 37:33}{Genesis 37:33}\vs{33} He recognized it and exclaimed, ``It is my son's tunic! A wild animal has eaten him! Joseph has surely been torn to pieces!''
\markboth{Genesis 37:34}{Genesis 37:34}\vs{34} Then Jacob tore his clothes, put on sackcloth, and mourned for his son many days.
\markboth{Genesis 37:35}{Genesis 37:35}\vs{35} All his sons and daughters stood by him to console him, but he refused to be consoled. ``No,'' he said, ``I will go to the grave mourning my son.'' So Joseph's father wept for him.
\everypar{}

\parshape=0
\markboth{Genesis 37:36}{Genesis 37:36}\vs{36} Now in Egypt the Midianites sold Joseph to Potiphar, one of Pharaoh's officials, the captain of the guard.\gva{l}\marginnote{\gva{l}\,Or "eunuch", which does not always signify a man that is gelded, but also someone that is in some high position.}

\markboth{Genesis 38:1}{Genesis 38:1}
\Needspace*{8\baselineskip}
\ch{38} \allowchapbreak\hypertarget{ch-genesis-38}{}\lettrine[lines=5]{\color{pentateuch}A}{\textsc{t}} that time Judah left his brothers and stayed with an Adullamite man named Hirah.\gva{a}\marginnote{\gva{a}\,Moses describes the genealogy of Judah, because the Messiah should come from him.}\everypar{}
\markboth{Genesis 38:2}{Genesis 38:2}\vs{2} There Judah saw the daughter of a Canaanite man named Shua. Judah acquired her as a wife and slept with her.\gva{b}\marginnote{\gva{b}\,A relationship which nonetheless was condemned by God.}
\markboth{Genesis 38:3}{Genesis 38:3}\vs{3} She became pregnant and had a son. Judah named him Er.
\markboth{Genesis 38:4}{Genesis 38:4}\vs{4} She became pregnant again and had another son, whom she named Onan.
\markboth{Genesis 38:5}{Genesis 38:5}\vs{5} Then she had yet another son, whom she named Shelah. She gave birth to him in Kezib.
\everypar{}

\parshape=0
\markboth{Genesis 38:6}{Genesis 38:6}\vs{6} Judah acquired a wife for Er his firstborn; her name was Tamar.
\markboth{Genesis 38:7}{Genesis 38:7}\vs{7} But Er, Judah's firstborn, was evil in the Lord's sight, so the Lord killed him.
\everypar{}

\parshape=0
\markboth{Genesis 38:8}{Genesis 38:8}\vs{8} Then Judah said to Onan, ``Sleep with your brother's wife and fulfill the duty of a brother-in-law to her so that you may raise up a descendant for your brother.''\gva{c}\marginnote{\gva{c}\,This order was for the preservation of the stock, since the child begotten by the second brother would have the name and inheritance of the first: a practice which is abolished in the New Testament.}
\markboth{Genesis 38:9}{Genesis 38:9}\vs{9} But Onan knew that the child would not be considered his. So whenever he slept with his brother's wife, he wasted his emission on the ground so as not to give his brother a descendant.
\markboth{Genesis 38:10}{Genesis 38:10}\vs{10} What he did was evil in the Lord's sight, so the Lord killed him too.
\everypar{}

\parshape=0
\markboth{Genesis 38:11}{Genesis 38:11}\vs{11} Then Judah said to his daughter-in-law Tamar, ``Live as a widow in your father's house until Shelah my son grows up.'' For he thought, ``I don't want him to die like his brothers.'' So Tamar went and lived in her father's house.\gva{d}\marginnote{\gva{d}\,For she could not marry in any other family so long as Judah would retain her in his.}
\everypar{}

\parshape=0
\markboth{Genesis 38:12}{Genesis 38:12}\vs{12} After some time Judah's wife, the daughter of Shua, died. After Judah was consoled, he left for Timnah to visit his sheepshearers, along with his friend Hirah the Adullamite.
\markboth{Genesis 38:13}{Genesis 38:13}\vs{13} Tamar was told, ``Look, your father-in-law is going up to Timnah to shear his sheep.''
\markboth{Genesis 38:14}{Genesis 38:14}\vs{14} So she removed her widow's clothes and covered herself with a veil. She wrapped herself and sat at the entrance to Enaim which is on the way to Timnah. (She did this because she saw that she had not been given to Shelah as a wife, even though he had now grown up.)
\everypar{}

\parshape=0
\markboth{Genesis 38:15}{Genesis 38:15}\vs{15} When Judah saw her, he thought she was a prostitute because she had covered her face.
\markboth{Genesis 38:16}{Genesis 38:16}\vs{16} He turned aside to her along the road and said, ``Come, please, I want to sleep with you.'' (He did not realize it was his daughter-in-law.) She asked, ``What will you give me so that you may sleep with me?''\gva{e}\marginnote{\gva{e}\,God miraculously blinded him so that he could not know her by her voice.}
\markboth{Genesis 38:17}{Genesis 38:17}\vs{17} He replied, ``I'll send you a young goat from the flock.'' She asked, ``Will you give me a pledge until you send it?''
\markboth{Genesis 38:18}{Genesis 38:18}\vs{18} He said, ``What pledge should I give you?'' She replied, ``Your seal, your cord, and the staff that's in your hand.'' So he gave them to her, then slept with her, and she became pregnant by him.
\markboth{Genesis 38:19}{Genesis 38:19}\vs{19} She left immediately, removed her veil, and put on her widow's clothes.
\everypar{}

\parshape=0
\markboth{Genesis 38:20}{Genesis 38:20}\vs{20} Then Judah had his friend Hirah the Adullamite take a young goat to get back from the woman the items he had given in pledge, but Hirah could not find her.\gva{f}\marginnote{\gva{f}\,That his wickedness might not be known to others.}
\markboth{Genesis 38:21}{Genesis 38:21}\vs{21} He asked the men who were there, ``Where is the cult prostitute who was at Enaim by the road?'' But they replied, ``There has been no cult prostitute here.''
\markboth{Genesis 38:22}{Genesis 38:22}\vs{22} So he returned to Judah and said, ``I couldn't find her. Moreover, the men of the place said, `There has been no cult prostitute here.'''
\markboth{Genesis 38:23}{Genesis 38:23}\vs{23} Judah said, ``Let her keep the things for herself. Otherwise we will appear to be dishonest. I did indeed send this young goat, but you couldn't find her.''\gva{g}\marginnote{\gva{g}\,He fears man more than God.}
\everypar{}

\parshape=0
\markboth{Genesis 38:24}{Genesis 38:24}\vs{24} After three months Judah was told, ``Your daughter-in-law Tamar has turned to prostitution, and as a result she has become pregnant.'' Judah said, ``Bring her out and let her be burned!''\gva{h}\marginnote{\gva{h}\,We see that the Law, which was written in man's heart, taught them that adultery should be punished with death, even though no law had been given yet.}
\markboth{Genesis 38:25}{Genesis 38:25}\vs{25} While they were bringing her out, she sent word to her father-in-law: ``I am pregnant by the man to whom these belong.'' Then she said, ``Identify the one to whom the seal, cord, and staff belong.''
\markboth{Genesis 38:26}{Genesis 38:26}\vs{26} Judah recognized them and said, ``She is more upright than I am, because I wouldn't give her to Shelah my son.'' He was not physically intimate with her again.\gva{i}\marginnote{\gva{i}\,That is, she ought rather to accuse me than I her.}\gva{k}\marginnote{\gva{k}\,for the horror of the sin condemned him.}
\everypar{}

\parshape=0
\markboth{Genesis 38:27}{Genesis 38:27}\vs{27} When it was time for her to give birth, there were twins in her womb.
\markboth{Genesis 38:28}{Genesis 38:28}\vs{28} While she was giving birth, one child put out his hand, and the midwife took a scarlet thread and tied it on his hand, saying, ``This one came out first.''
\markboth{Genesis 38:29}{Genesis 38:29}\vs{29} But then he drew back his hand, and his brother came out before him. She said, ``How you have broken out of the womb!'' So he was named Perez.\gva{l}\marginnote{\gva{l}\,Their heinous sin was signified by this monstrous birth.}\gva{m}\marginnote{\gva{m}\,Or the separation between you and your brother.}
\markboth{Genesis 38:30}{Genesis 38:30}\vs{30} Afterward his brother came out---the one who had the scarlet thread on his hand---and he was named Zerah.

\markboth{Genesis 39:1}{Genesis 39:1}
\Needspace*{8\baselineskip}
\ch{39} \allowchapbreak\hypertarget{ch-genesis-39}{}\lettrine[lines=5]{\color{pentateuch}N}{\textsc{ow}} Joseph had been brought down to Egypt. An Egyptian named Potiphar, an official of Pharaoh and the captain of the guard, purchased him from the Ishmaelites who had brought him there.\gva{a}\marginnote{\gva{a}\,See Genesis 37:36 .}\everypar{}
\markboth{Genesis 39:2}{Genesis 39:2}\vs{2} The Lord was with Joseph. He was successful and lived in the household of his Egyptian master.\gva{b}\marginnote{\gva{b}\,The favour of God is the fountain of all prosperity.}
\markboth{Genesis 39:3}{Genesis 39:3}\vs{3} His master observed that the Lord was with him and that the Lord made everything he was doing successful.
\markboth{Genesis 39:4}{Genesis 39:4}\vs{4} So Joseph found favor in his sight and became his personal attendant. Potiphar appointed Joseph overseer of his household and put him in charge of everything he owned.\gva{c}\marginnote{\gva{c}\,Because God prospered him: and so he made religion serve his profit.}
\markboth{Genesis 39:5}{Genesis 39:5}\vs{5} From the time Potiphar appointed him over his household and over all that he owned, the Lord blessed the Egyptian's household for Joseph's sake. The blessing of the Lord was on everything that he had, both in his house and in his fields.\gva{d}\marginnote{\gva{d}\,The wicked are blessed by the company of the godly.}
\everypar{}

\parshape=0
\markboth{Genesis 39:6}{Genesis 39:6}\vs{6} So Potiphar left everything he had in Joseph's care; he gave no thought to anything except the food he ate. Now Joseph was well built and good-looking.\gva{e}\marginnote{\gva{e}\,For he was sure that everything would prosper: therefore he ate and drank and did not worry.}
\markboth{Genesis 39:7}{Genesis 39:7}\vs{7} Soon after these things, his master's wife took notice of Joseph and said, ``Come to bed with me.''\gva{f}\marginnote{\gva{f}\,In this word he declares the purpose she was working towards.}
\markboth{Genesis 39:8}{Genesis 39:8}\vs{8} But he refused, saying to his master's wife, ``Look, my master does not give any thought to his household with me here, and everything that he owns he has put into my care.
\markboth{Genesis 39:9}{Genesis 39:9}\vs{9} There is no one greater in this household than I am. He has withheld nothing from me except you because you are his wife. So how could I do such a great evil and sin against God?''\gva{g}\marginnote{\gva{g}\,The fear of God preserved him against her continual temptations.}
\markboth{Genesis 39:10}{Genesis 39:10}\vs{10} Even though she continued to speak to Joseph day after day, he did not respond to her invitation to go to bed with her.
\everypar{}

\parshape=0
\markboth{Genesis 39:11}{Genesis 39:11}\vs{11} One day he went into the house to do his work when none of the household servants were there in the house.
\markboth{Genesis 39:12}{Genesis 39:12}\vs{12} She grabbed him by his outer garment, saying, ``Come to bed with me!'' But he left his outer garment in her hand and ran outside.
\markboth{Genesis 39:13}{Genesis 39:13}\vs{13} When she saw that he had left his outer garment in her hand and had run outside,
\markboth{Genesis 39:14}{Genesis 39:14}\vs{14} she called for her household servants and said to them, ``See, my husband brought in a Hebrew man to us to humiliate us. He tried to go to bed with me, but I screamed loudly.\gva{h}\marginnote{\gva{h}\,This declares that in which lack of restraint exists and to this is joined extreme impudency and deceit.}
\markboth{Genesis 39:15}{Genesis 39:15}\vs{15} When he heard me raise my voice and scream, he left his outer garment beside me and ran outside.''
\everypar{}

\parshape=0
\markboth{Genesis 39:16}{Genesis 39:16}\vs{16} So she laid his outer garment beside her until his master came home.
\markboth{Genesis 39:17}{Genesis 39:17}\vs{17} This is what she said to him: ``That Hebrew slave you brought to us tried to humiliate me,
\markboth{Genesis 39:18}{Genesis 39:18}\vs{18} but when I raised my voice and screamed, he left his outer garment and ran outside.''
\everypar{}

\parshape=0
\markboth{Genesis 39:19}{Genesis 39:19}\vs{19} When his master heard his wife say, ``This is the way your slave treated me,'' he became furious.
\markboth{Genesis 39:20}{Genesis 39:20}\vs{20} Joseph's master took him and threw him into the prison, the place where the king's prisoners were confined. So he was there in the prison.\gva{i}\marginnote{\gva{i}\,His bad treatment in the prison may be gathered from Psalms 105:18 .}
\everypar{}

\parshape=0
\markboth{Genesis 39:21}{Genesis 39:21}\vs{21} But the Lord was with Joseph and showed him kindness. He granted him favor in the sight of the prison warden.
\markboth{Genesis 39:22}{Genesis 39:22}\vs{22} The warden put all the prisoners under Joseph's care. He was in charge of whatever they were doing.\gva{k}\marginnote{\gva{k}\,That is, nothing was done without his commandment.}
\markboth{Genesis 39:23}{Genesis 39:23}\vs{23} The warden did not concern himself with anything that was in Joseph's care because the Lord was with him and whatever he was doing the Lord was making successful.

\markboth{Genesis 40:1}{Genesis 40:1}
\Needspace*{8\baselineskip}
\ch{40} \allowchapbreak\hypertarget{ch-genesis-40}{}\lettrine[lines=5]{\color{pentateuch}A}{\textsc{fter}} these things happened, the cupbearer to the king of Egypt and the royal baker offended their master, the king of Egypt.\everypar{}
\markboth{Genesis 40:2}{Genesis 40:2}\vs{2} Pharaoh was enraged with his two officials, the cupbearer and the baker,
\markboth{Genesis 40:3}{Genesis 40:3}\vs{3} so he imprisoned them in the house of the captain of the guard in the same facility where Joseph was confined.\gva{a}\marginnote{\gva{a}\,God works in many wonderful ways to deliver his own.}
\everypar{}

\parshape=0
\markboth{Genesis 40:4}{Genesis 40:4}\vs{4} The captain of the guard appointed Joseph to be their attendant, and he served them. They spent some time in custody.
\markboth{Genesis 40:5}{Genesis 40:5}\vs{5} Both of them, the cupbearer and the baker of the king of Egypt, who were confined in the prison, had a dream the same night. Each man's dream had its own meaning.\gva{b}\marginnote{\gva{b}\,That is, every dream had his interpretation, as the thing afterward declared.}
\markboth{Genesis 40:6}{Genesis 40:6}\vs{6} When Joseph came to them in the morning, he saw that they were looking depressed.
\markboth{Genesis 40:7}{Genesis 40:7}\vs{7} So he asked Pharaoh's officials, who were with him in custody in his master's house, ``Why do you look so sad today?''
\markboth{Genesis 40:8}{Genesis 40:8}\vs{8} They told him, ``We both had dreams, but there is no one to interpret them.'' Joseph responded, ``Don't interpretations belong to God? Tell them to me.''\gva{c}\marginnote{\gva{c}\,Cannot God raise up such as shall interpret such things.}
\everypar{}

\parshape=0
\markboth{Genesis 40:9}{Genesis 40:9}\vs{9} So the chief cupbearer told his dream to Joseph: ``In my dream, there was a vine in front of me.
\markboth{Genesis 40:10}{Genesis 40:10}\vs{10} On the vine there were three branches. As it budded, its blossoms opened and its clusters ripened into grapes.
\markboth{Genesis 40:11}{Genesis 40:11}\vs{11} Now Pharaoh's cup was in my hand, so I took the grapes, squeezed them into his cup, and put the cup in Pharaoh's hand.''
\everypar{}

\parshape=0
\markboth{Genesis 40:12}{Genesis 40:12}\vs{12} ``This is its meaning,'' Joseph said to him. ``The three branches represent three days.\gva{d}\marginnote{\gva{d}\,He was reassured by the spirit of God, that his interpretation was true.}
\markboth{Genesis 40:13}{Genesis 40:13}\vs{13} In three more days Pharaoh will reinstate you and restore you to your office. You will put Pharaoh's cup in his hand, just as you did before when you were cupbearer.
\markboth{Genesis 40:14}{Genesis 40:14}\vs{14} But remember me when it goes well for you, and show me kindness. Make mention of me to Pharaoh and bring me out of this prison,\gva{e}\marginnote{\gva{e}\,He does not refuse the method of deliverance which he thought God had appointed.}
\markboth{Genesis 40:15}{Genesis 40:15}\vs{15} for I really was kidnapped from the land of the Hebrews and I have done nothing wrong here for which they should put me in a dungeon.''
\everypar{}

\parshape=0
\markboth{Genesis 40:16}{Genesis 40:16}\vs{16} When the chief baker saw that the interpretation of the first dream was favorable, he said to Joseph, ``I also appeared in my dream and there were three baskets of white bread on my head.\gva{f}\marginnote{\gva{f}\,That is made of white twigs, or as some read, baskets full of holes.}
\markboth{Genesis 40:17}{Genesis 40:17}\vs{17} In the top basket there were baked goods of every kind for Pharaoh, but the birds were eating them from the basket that was on my head.''
\everypar{}

\parshape=0
\markboth{Genesis 40:18}{Genesis 40:18}\vs{18} Joseph replied, ``This is its meaning: The three baskets represent three days.\gva{g}\marginnote{\gva{g}\,He shows that the ministers of God should not conceal that, which God reveals to them.}
\markboth{Genesis 40:19}{Genesis 40:19}\vs{19} In three more days Pharaoh will decapitate you and impale you on a pole. Then the birds will eat your flesh from you.''
\everypar{}

\parshape=0
\markboth{Genesis 40:20}{Genesis 40:20}\vs{20} On the third day it was Pharaoh's birthday, so he gave a feast for all his servants. He ``lifted up'' the head of the chief cupbearer and the head of the chief baker in the midst of his servants.\gva{h}\marginnote{\gva{h}\,Which was an occasion to appoint his officers, and to examine those who were in prison.}
\markboth{Genesis 40:21}{Genesis 40:21}\vs{21} He restored the chief cupbearer to his former position so that he placed the cup in Pharaoh's hand,
\markboth{Genesis 40:22}{Genesis 40:22}\vs{22} but the chief baker he impaled, just as Joseph had predicted.
\markboth{Genesis 40:23}{Genesis 40:23}\vs{23} But the chief cupbearer did not remember Joseph---he forgot him.

\markboth{Genesis 41:1}{Genesis 41:1}
\Needspace*{8\baselineskip}
\ch{41} \allowchapbreak\hypertarget{ch-genesis-41}{}\lettrine[lines=5]{\color{pentateuch}A}{\textsc{t}} the end of two full years Pharaoh had a dream. As he was standing by the Nile,\gva{a}\marginnote{\gva{a}\,This dream was not so much for Pharaoh, as is was a means to deliver Joseph and to provide for God's Church.}\everypar{}
\markboth{Genesis 41:2}{Genesis 41:2}\vs{2} seven fine-looking, fat cows were coming up out of the Nile, and they grazed in the reeds.
\markboth{Genesis 41:3}{Genesis 41:3}\vs{3} Then seven bad-looking, thin cows were coming up after them from the Nile, and they stood beside the other cows at the edge of the river.
\markboth{Genesis 41:4}{Genesis 41:4}\vs{4} The bad-looking, thin cows ate the seven fine-looking, fat cows. Then Pharaoh woke up.
\everypar{}

\parshape=0
\markboth{Genesis 41:5}{Genesis 41:5}\vs{5} Then he fell asleep again and had a second dream: There were seven heads of grain growing on one stalk, healthy and good.\gva{b}\marginnote{\gva{b}\,All these means God used to deliver his servant, and to bring him into favour and authority.}
\markboth{Genesis 41:6}{Genesis 41:6}\vs{6} Then seven heads of grain, thin and burned by the east wind, were sprouting up after them.
\markboth{Genesis 41:7}{Genesis 41:7}\vs{7} The thin heads swallowed up the seven healthy and full heads. Then Pharaoh woke up and realized it was a dream.
\everypar{}

\parshape=0
\markboth{Genesis 41:8}{Genesis 41:8}\vs{8} In the morning he was troubled, so he called for all the diviner-priests of Egypt and all its wise men. Pharaoh told them his dreams, but no one could interpret them for him.\gva{c}\marginnote{\gva{c}\,This fear was enough to teach him that this vision was sent by God.}
\markboth{Genesis 41:9}{Genesis 41:9}\vs{9} Then the chief cupbearer said to Pharaoh, ``Today I recall my failures.\gva{e}\marginnote{\gva{e}\,He confesses his fault against the king before he speaks of Joseph.}
\markboth{Genesis 41:10}{Genesis 41:10}\vs{10} Pharaoh was enraged with his servants, and he put me in prison in the house of the captain of the guards---me and the chief baker.
\markboth{Genesis 41:11}{Genesis 41:11}\vs{11} We each had a dream one night; each of us had a dream with its own meaning.
\markboth{Genesis 41:12}{Genesis 41:12}\vs{12} Now a young man, a Hebrew, a servant of the captain of the guards, was with us there. We told him our dreams, and he interpreted the meaning of each of our respective dreams for us.
\markboth{Genesis 41:13}{Genesis 41:13}\vs{13} It happened just as he had said to us---Pharaoh restored me to my office, but he impaled the baker.''
\everypar{}

\parshape=0
\markboth{Genesis 41:14}{Genesis 41:14}\vs{14} Then Pharaoh summoned Joseph. So they brought him quickly out of the dungeon; he shaved himself, changed his clothes, and came before Pharaoh.\gva{f}\marginnote{\gva{f}\,The wicked seek the prophets of God in their time of need, while in their prosperity they abhor them.}
\markboth{Genesis 41:15}{Genesis 41:15}\vs{15} Pharaoh said to Joseph, ``I had a dream, and there is no one who can interpret it. But I have heard about you, that you can interpret dreams.''
\markboth{Genesis 41:16}{Genesis 41:16}\vs{16} Joseph replied to Pharaoh, ``It is not within my power, but God will speak concerning the welfare of Pharaoh.''\gva{g}\marginnote{\gva{g}\,As though he would say if I interpret your dream it comes from God, and not from me.}
\everypar{}

\parshape=0
\markboth{Genesis 41:17}{Genesis 41:17}\vs{17} Then Pharaoh said to Joseph, ``In my dream I was standing by the edge of the Nile.
\markboth{Genesis 41:18}{Genesis 41:18}\vs{18} Then seven fat and fine-looking cows were coming up out of the Nile, and they grazed in the reeds.
\markboth{Genesis 41:19}{Genesis 41:19}\vs{19} Then seven other cows came up after them; they were scrawny, very bad looking, and lean. I had never seen such bad-looking cows as these in all the land of Egypt!
\markboth{Genesis 41:20}{Genesis 41:20}\vs{20} The lean, bad-looking cows ate up the seven fat cows.
\markboth{Genesis 41:21}{Genesis 41:21}\vs{21} When they had eaten them, no one would have known that they had done so, for they were just as bad looking as before. Then I woke up.
\markboth{Genesis 41:22}{Genesis 41:22}\vs{22} I also saw in my dream seven heads of grain growing on one stalk, full and good.
\markboth{Genesis 41:23}{Genesis 41:23}\vs{23} Then seven heads of grain, withered and thin and burned with the east wind, were sprouting up after them.
\markboth{Genesis 41:24}{Genesis 41:24}\vs{24} The thin heads of grain swallowed up the seven good heads of grain. So I told all this to the diviner-priests, but no one could tell me its meaning.''
\everypar{}

\parshape=0
\markboth{Genesis 41:25}{Genesis 41:25}\vs{25} Then Joseph said to Pharaoh, ``Both dreams of Pharaoh have the same meaning. God has revealed to Pharaoh what he is about to do.\gva{h}\marginnote{\gva{h}\,Both his dreams have the same message.}
\markboth{Genesis 41:26}{Genesis 41:26}\vs{26} The seven good cows represent seven years, and the seven good heads of grain represent seven years. Both dreams have the same meaning.
\markboth{Genesis 41:27}{Genesis 41:27}\vs{27} The seven lean, bad-looking cows that came up after them represent seven years, as do the seven empty heads of grain burned with the east wind. They represent seven years of famine.
\markboth{Genesis 41:28}{Genesis 41:28}\vs{28} This is just what I told Pharaoh: God has shown Pharaoh what he is about to do.
\markboth{Genesis 41:29}{Genesis 41:29}\vs{29} Seven years of great abundance are coming throughout the whole land of Egypt.
\markboth{Genesis 41:30}{Genesis 41:30}\vs{30} But seven years of famine will occur after them, and all the abundance will be forgotten in the land of Egypt. The famine will devastate the land.
\markboth{Genesis 41:31}{Genesis 41:31}\vs{31} The previous abundance of the land will not be remembered because of the famine that follows, for the famine will be very severe.
\markboth{Genesis 41:32}{Genesis 41:32}\vs{32} The dream was repeated to Pharaoh because the matter has been decreed by God, and God will make it happen soon.
\everypar{}

\parshape=0
\markboth{Genesis 41:33}{Genesis 41:33}\vs{33} ``So now Pharaoh should look for a wise and discerning man and give him authority over all the land of Egypt.\gva{i}\marginnote{\gva{i}\,The office of a true prophet is not only to show the evils to come, but also the remedies for the same.}
\markboth{Genesis 41:34}{Genesis 41:34}\vs{34} Pharaoh should do this---he should appoint officials throughout the land to collect one-fifth of the produce of the land of Egypt during the seven years of abundance.
\markboth{Genesis 41:35}{Genesis 41:35}\vs{35} They should gather all the excess food during these good years that are coming. By Pharaoh's authority they should store up grain so the cities will have food, and they should preserve it.
\markboth{Genesis 41:36}{Genesis 41:36}\vs{36} This food should be held in storage for the land in preparation for the seven years of famine that will occur throughout the land of Egypt. In this way the land will survive the famine.''
\everypar{}

\parshape=0
\markboth{Genesis 41:37}{Genesis 41:37}\vs{37} This advice made sense to Pharaoh and all his officials.
\markboth{Genesis 41:38}{Genesis 41:38}\vs{38} So Pharaoh asked his officials, ``Can we find a man like Joseph, one in whom the Spirit of God is present?''\gva{k}\marginnote{\gva{k}\,No one should be honoured who does not have gifts from God fitting for the same.}
\markboth{Genesis 41:39}{Genesis 41:39}\vs{39} So Pharaoh said to Joseph, ``Because God has enabled you to know all this, there is no one as wise and discerning as you are!
\markboth{Genesis 41:40}{Genesis 41:40}\vs{40} You will oversee my household, and all my people will submit to your commands. Only I, the king, will be greater than you.\gva{l}\marginnote{\gva{l}\,Some read, "the people will kill your mouth", that is obey you in all things.}
\everypar{}

\parshape=0
\markboth{Genesis 41:41}{Genesis 41:41}\vs{41} ``See here,'' Pharaoh said to Joseph, ``I place you in authority over all the land of Egypt.''
\markboth{Genesis 41:42}{Genesis 41:42}\vs{42} Then Pharaoh took his signet ring from his own hand and put it on Joseph's. He clothed him with fine linen clothes and put a gold chain around his neck.
\markboth{Genesis 41:43}{Genesis 41:43}\vs{43} Pharaoh had him ride in the chariot used by his second-in-command, and they cried out before him, ``Kneel down!'' So he placed him over all the land of Egypt.\gva{b}\marginnote{\gva{b}\,Or "Abrech": a sign of honour; a word some translate, tender father or father of the king, or kneel down.}
\markboth{Genesis 41:44}{Genesis 41:44}\vs{44} Pharaoh also said to Joseph, ``I am Pharaoh, but without your permission no one will move his hand or his foot in all the land of Egypt.''
\markboth{Genesis 41:45}{Genesis 41:45}\vs{45} Pharaoh gave Joseph the name Zaphenath-Paneah. He also gave him Asenath daughter of Potiphera, priest of On, to be his wife. So Joseph took charge of all the land of Egypt.
\everypar{}

\parshape=0
\markboth{Genesis 41:46}{Genesis 41:46}\vs{46} Now Joseph was 30 years old when he began serving Pharaoh king of Egypt. Joseph was commissioned by Pharaoh and was in charge of all the land of Egypt.\gva{n}\marginnote{\gva{n}\,His age is mentioned both to show that his authority came from God, and also that he endured imprisonment and exile for twelve years or more.}
\markboth{Genesis 41:47}{Genesis 41:47}\vs{47} During the seven years of abundance the land produced large, bountiful harvests.
\markboth{Genesis 41:48}{Genesis 41:48}\vs{48} Joseph collected all the excess food in the land of Egypt during the seven years and stored it in the cities. In every city he put the food gathered from the fields around it.
\markboth{Genesis 41:49}{Genesis 41:49}\vs{49} Joseph stored up a vast amount of grain, like the sand of the sea, until he stopped measuring it because it was impossible to measure.
\everypar{}

\parshape=0
\markboth{Genesis 41:50}{Genesis 41:50}\vs{50} Two sons were born to Joseph before the famine came. Asenath daughter of Potiphera, priest of On, was their mother.
\markboth{Genesis 41:51}{Genesis 41:51}\vs{51} Joseph named the firstborn Manasseh, saying, ``Certainly God has made me forget all my trouble and all my father's house.''\gva{o}\marginnote{\gva{o}\,Nonetheless, his father's house was the true Church of God: yet the company of the wicked and prosperity caused him to forget it.}
\markboth{Genesis 41:52}{Genesis 41:52}\vs{52} He named the second child Ephraim, saying, ``Certainly God has made me fruitful in the land of my suffering.''
\everypar{}

\parshape=0
\markboth{Genesis 41:53}{Genesis 41:53}\vs{53} The seven years of abundance in the land of Egypt came to an end.
\markboth{Genesis 41:54}{Genesis 41:54}\vs{54} Then the seven years of famine began, just as Joseph had predicted. There was famine in all the other lands, but throughout the land of Egypt there was food.
\markboth{Genesis 41:55}{Genesis 41:55}\vs{55} When all the land of Egypt experienced the famine, the people cried out to Pharaoh for food. Pharaoh said to all the people of Egypt, ``Go to Joseph and do whatever he tells you.''
\everypar{}

\parshape=0
\markboth{Genesis 41:56}{Genesis 41:56}\vs{56} While the famine was over all the earth, Joseph opened the storehouses and sold grain to the Egyptians. The famine was severe throughout the land of Egypt.
\markboth{Genesis 41:57}{Genesis 41:57}\vs{57} People from every country came to Joseph in Egypt to buy grain because the famine was severe throughout the earth.

\markboth{Genesis 42:1}{Genesis 42:1}
\Needspace*{8\baselineskip}
\ch{42} \allowchapbreak\hypertarget{ch-genesis-42}{}\lettrine[lines=5]{\color{pentateuch}W}{\textsc{hen}} Jacob heard there was grain in Egypt, he said to his sons, ``Why are you looking at each other?''\gva{a}\marginnote{\gva{a}\,This story shows plainly that all things are governed by God's providence for the profit of his Church.}\gva{b}\marginnote{\gva{b}\,As men destitute of counsel.}\everypar{}
\markboth{Genesis 42:2}{Genesis 42:2}\vs{2} He then said, ``Look, I hear that there is grain in Egypt. Go down there and buy grain for us so that we may live and not die.''
\\\indent\markboth{Genesis 42:3}{Genesis 42:3}\vs{3} So ten of Joseph's brothers went down to buy grain from Egypt.
\markboth{Genesis 42:4}{Genesis 42:4}\vs{4} But Jacob did not send Joseph's brother Benjamin with his brothers, for he said, ``What if some accident happens to him?''
\markboth{Genesis 42:5}{Genesis 42:5}\vs{5} So Israel's sons came to buy grain among the other travelers, for the famine was severe in the land of Canaan.
\everypar{}

\parshape=0
\markboth{Genesis 42:6}{Genesis 42:6}\vs{6} Now Joseph was the ruler of the country, the one who sold grain to all the people of the country. Joseph's brothers came and bowed down before him with their faces to the ground.
\markboth{Genesis 42:7}{Genesis 42:7}\vs{7} When Joseph saw his brothers, he recognized them, but he pretended to be a stranger to them and spoke to them harshly. He asked, ``Where do you come from?'' They answered, ``From the land of Canaan, to buy grain for food.''\gva{c}\marginnote{\gva{c}\,This concealing is not to be followed, nor any actions of the father's not approved by God's word.}
\everypar{}

\parshape=0
\markboth{Genesis 42:8}{Genesis 42:8}\vs{8} Joseph recognized his brothers, but they did not recognize him.
\markboth{Genesis 42:9}{Genesis 42:9}\vs{9} Then Joseph remembered the dreams he had dreamed about them, and he said to them, ``You are spies; you have come to see if our land is vulnerable!''
\everypar{}

\parshape=0
\markboth{Genesis 42:10}{Genesis 42:10}\vs{10} But they exclaimed, ``No, my lord! Your servants have come to buy grain for food!
\markboth{Genesis 42:11}{Genesis 42:11}\vs{11} We are all the sons of one man; we are honest men! Your servants are not spies.''
\everypar{}

\parshape=0
\markboth{Genesis 42:12}{Genesis 42:12}\vs{12} ``No,'' he insisted, ``but you have come to see if our land is vulnerable.''
\markboth{Genesis 42:13}{Genesis 42:13}\vs{13} They replied, ``Your servants are from a family of twelve brothers. We are the sons of one man in the land of Canaan. The youngest is with our father at this time, and one is no longer alive.''
\everypar{}

\parshape=0
\markboth{Genesis 42:14}{Genesis 42:14}\vs{14} But Joseph told them, ``It is just as I said to you: You are spies!
\markboth{Genesis 42:15}{Genesis 42:15}\vs{15} You will be tested in this way: As surely as Pharaoh lives, you will not depart from this place unless your youngest brother comes here.\gva{d}\marginnote{\gva{d}\,The Egyptians who were idolaters, used to swear by their king's life: but God forbids swearing by anyone but him: yet Joseph dwelling among the wicked was corrupted by them.}
\markboth{Genesis 42:16}{Genesis 42:16}\vs{16} One of you must go and get your brother, while the rest of you remain in prison. In this way your words may be tested to see if you are telling the truth. If not, then, as surely as Pharaoh lives, you are spies!''
\markboth{Genesis 42:17}{Genesis 42:17}\vs{17} He imprisoned them all for three days.
\markboth{Genesis 42:18}{Genesis 42:18}\vs{18} On the third day Joseph said to them, ``Do as I say and you will live, for I fear God.\gva{e}\marginnote{\gva{e}\,And therefore am true and just.}
\markboth{Genesis 42:19}{Genesis 42:19}\vs{19} If you are honest men, leave one of your brothers confined here in prison while the rest of you go and take grain back for your hungry families.
\markboth{Genesis 42:20}{Genesis 42:20}\vs{20} But you must bring your youngest brother to me. Then your words will be verified and you will not die.'' They did as he said.
\everypar{}

\parshape=0
\markboth{Genesis 42:21}{Genesis 42:21}\vs{21} They said to one another, ``Surely we're being punished because of our brother, because we saw how distressed he was when he cried to us for mercy, but we refused to listen. That is why this distress has come on us!''\gva{f}\marginnote{\gva{f}\,Affliction makes men acknowledge their faults, which otherwise they would conceal.}
\markboth{Genesis 42:22}{Genesis 42:22}\vs{22} Reuben said to them, ``Didn't I say to you, `Don't sin against the boy,' but you wouldn't listen? So now we must pay for shedding his blood!''\gva{g}\marginnote{\gva{g}\,God will take vengeance on us, and measure us with our own measure.}
\markboth{Genesis 42:23}{Genesis 42:23}\vs{23} (Now they did not know that Joseph could understand them, for he was speaking through an interpreter.)
\markboth{Genesis 42:24}{Genesis 42:24}\vs{24} He turned away from them and wept. When he turned around and spoke to them again, he had Simeon taken from them and tied up before their eyes.\gva{h}\marginnote{\gva{h}\,Though he acts harshly, yet his brotherly affection remained.}
\everypar{}

\parshape=0
\markboth{Genesis 42:25}{Genesis 42:25}\vs{25} Then Joseph gave orders to fill their bags with grain, to return each man's money to his sack, and to give them provisions for the journey. His orders were carried out.
\markboth{Genesis 42:26}{Genesis 42:26}\vs{26} So they loaded their grain on their donkeys and left.
\everypar{}

\parshape=0
\markboth{Genesis 42:27}{Genesis 42:27}\vs{27} When one of them opened his sack to get feed for his donkey at their resting place, he saw his money in the mouth of his sack.
\markboth{Genesis 42:28}{Genesis 42:28}\vs{28} He said to his brothers, ``My money was returned! Here it is in my sack!'' They were dismayed; they turned trembling to one another and said, ``What in the world has God done to us?''\gva{i}\marginnote{\gva{i}\,Because their conscience accused them of their sin, they thought God had brought them trouble through the money.}
\everypar{}

\parshape=0
\markboth{Genesis 42:29}{Genesis 42:29}\vs{29} They returned to their father Jacob in the land of Canaan and told him all the things that had happened to them, saying,
\markboth{Genesis 42:30}{Genesis 42:30}\vs{30} ``The man, the lord of the land, spoke harshly to us and treated us as if we were spying on the land.
\markboth{Genesis 42:31}{Genesis 42:31}\vs{31} But we said to him, `We are honest men; we are not spies!
\markboth{Genesis 42:32}{Genesis 42:32}\vs{32} We are from a family of twelve brothers; we are the sons of one father. One is no longer alive, and the youngest is with our father at this time in the land of Canaan.'
\everypar{}

\parshape=0
\markboth{Genesis 42:33}{Genesis 42:33}\vs{33} ``Then the man, the lord of the land, said to us, `This is how I will find out if you are honest men. Leave one of your brothers with me, and take grain for your hungry households and go.
\markboth{Genesis 42:34}{Genesis 42:34}\vs{34} But bring your youngest brother back to me so I will know that you are honest men and not spies. Then I will give your brother back to you and you may move about freely in the land.'''
\everypar{}

\parshape=0
\markboth{Genesis 42:35}{Genesis 42:35}\vs{35} When they were emptying their sacks, there was each man's bag of money in his sack! When they and their father saw the bags of money, they were afraid.
\markboth{Genesis 42:36}{Genesis 42:36}\vs{36} Their father Jacob said to them, ``You are making me childless! Joseph is gone. Simeon is gone. And now you want to take Benjamin! Everything is against me.''\gva{k}\marginnote{\gva{k}\,For they did not seem to be concerned or have any love for their brother which increased his sorrow: and partly as it appears he suspected them for Joseph.}
\everypar{}

\parshape=0
\markboth{Genesis 42:37}{Genesis 42:37}\vs{37} Then Reuben said to his father, ``You may put my two sons to death if I do not bring him back to you. Put him in my care and I will bring him back to you.''
\markboth{Genesis 42:38}{Genesis 42:38}\vs{38} But Jacob replied, ``My son will not go down there with you, for his brother is dead and he alone is left. If an accident happens to him on the journey you have to make, then you will bring down my gray hair in sorrow to the grave.''

\markboth{Genesis 43:1}{Genesis 43:1}
\Needspace*{8\baselineskip}
\ch{43} \allowchapbreak\hypertarget{ch-genesis-43}{}\lettrine[lines=5]{\color{pentateuch}N}{\textsc{ow}} the famine was severe in the land.\gva{a}\marginnote{\gva{a}\,This was a great temptation to Jacob to suffer such a great famine in the land where God had promised to bless him.}\everypar{}
\markboth{Genesis 43:2}{Genesis 43:2}\vs{2} When they finished eating the grain they had brought from Egypt, their father said to them, ``Return, buy us a little more food.''
\\\indent\markboth{Genesis 43:3}{Genesis 43:3}\vs{3} But Judah said to him, ``The man solemnly warned us, `You will not see my face unless your brother is with you.'
\markboth{Genesis 43:4}{Genesis 43:4}\vs{4} If you send our brother with us, we'll go down and buy food for you.
\markboth{Genesis 43:5}{Genesis 43:5}\vs{5} But if you will not send him, we won't go down there because the man said to us, `You will not see my face unless your brother is with you.'''
\everypar{}

\parshape=0
\markboth{Genesis 43:6}{Genesis 43:6}\vs{6} Israel said, ``Why did you bring this trouble on me by telling the man you had one more brother?''
\everypar{}

\parshape=0
\markboth{Genesis 43:7}{Genesis 43:7}\vs{7} They replied, ``The man questioned us thoroughly about ourselves and our family, saying, `Is your father still alive? Do you have another brother?' So we answered him in this way. How could we possibly know that he would say, `Bring your brother down'?''
\everypar{}

\parshape=0
\markboth{Genesis 43:8}{Genesis 43:8}\vs{8} Then Judah said to his father Israel, ``Send the boy with me and we will go immediately. Then we will live and not die---we and you and our little ones.
\markboth{Genesis 43:9}{Genesis 43:9}\vs{9} I myself pledge security for him; you may hold me liable. If I do not bring him back to you and place him here before you, I will bear the blame before you all my life.
\markboth{Genesis 43:10}{Genesis 43:10}\vs{10} But if we had not delayed, we could have traveled there and back twice by now!''
\everypar{}

\parshape=0
\markboth{Genesis 43:11}{Genesis 43:11}\vs{11} Then their father Israel said to them, ``If it must be so, then do this: Take some of the best products of the land in your bags, and take a gift down to the man---a little balm and a little honey, spices and myrrh, pistachios and almonds.
\markboth{Genesis 43:12}{Genesis 43:12}\vs{12} Take double the money with you; you must take back the money that was returned in the mouths of your sacks---perhaps it was an oversight.\gva{b}\marginnote{\gva{b}\,When we are in need or danger, God does not forbid us to use honest means to better our estate and condition.}
\markboth{Genesis 43:13}{Genesis 43:13}\vs{13} Take your brother too, and go right away to the man.
\markboth{Genesis 43:14}{Genesis 43:14}\vs{14} May the Sovereign God grant you mercy before the man so that he may release your other brother and Benjamin! As for me, if I lose my children I lose them.''\gva{c}\marginnote{\gva{c}\,Our main trust should be in God, not in worldly means.}\gva{d}\marginnote{\gva{d}\,He speaks these words not so much in despair, but to make his sons more careful to return with their brother.}
\everypar{}

\parshape=0
\markboth{Genesis 43:15}{Genesis 43:15}\vs{15} So the men took these gifts, and they took double the money with them, along with Benjamin. Then they hurried down to Egypt and stood before Joseph.
\markboth{Genesis 43:16}{Genesis 43:16}\vs{16} When Joseph saw Benjamin with them, he said to the servant who was over his household, ``Bring the men to the house. Slaughter an animal and prepare it, for the men will eat with me at noon.''
\markboth{Genesis 43:17}{Genesis 43:17}\vs{17} The man did just as Joseph said; he brought the men into Joseph's house.
\everypar{}

\parshape=0
\markboth{Genesis 43:18}{Genesis 43:18}\vs{18} But the men were afraid when they were brought to Joseph's house. They said, ``We are being brought in because of the money that was returned in our sacks last time. He wants to capture us, make us slaves, and take our donkeys!''\gva{e}\marginnote{\gva{e}\,So the judgment of God weighed on their consciences.}
\markboth{Genesis 43:19}{Genesis 43:19}\vs{19} So they approached the man who was in charge of Joseph's household and spoke to him at the entrance to the house.
\markboth{Genesis 43:20}{Genesis 43:20}\vs{20} They said, ``My lord, we did indeed come down the first time to buy food.
\markboth{Genesis 43:21}{Genesis 43:21}\vs{21} But when we came to the place where we spent the night, we opened our sacks and each of us found his money---the full amount---in the mouth of his sack. So we have returned it.
\markboth{Genesis 43:22}{Genesis 43:22}\vs{22} We have brought additional money with us to buy food. We do not know who put the money in our sacks!''
\everypar{}

\parshape=0
\markboth{Genesis 43:23}{Genesis 43:23}\vs{23} ``Everything is fine,'' the man in charge of Joseph's household told them. ``Don't be afraid. Your God and the God of your father has given you treasure in your sacks. I had your money.'' Then he brought Simeon out to them.\gva{f}\marginnote{\gva{f}\,Despite the corruption of Egypt, Joseph taught his family to fear God.}
\everypar{}

\parshape=0
\markboth{Genesis 43:24}{Genesis 43:24}\vs{24} The servant in charge brought the men into Joseph's house. He gave them water, and they washed their feet. Then he gave food to their donkeys.
\markboth{Genesis 43:25}{Genesis 43:25}\vs{25} They got their gifts ready for Joseph's arrival at noon, for they had heard that they were to have a meal there.
\everypar{}

\parshape=0
\markboth{Genesis 43:26}{Genesis 43:26}\vs{26} When Joseph came home, they presented him with the gifts they had brought inside, and they bowed down to the ground before him.
\markboth{Genesis 43:27}{Genesis 43:27}\vs{27} He asked them how they were doing. Then he said, ``Is your aging father well, the one you spoke about? Is he still alive?''
\markboth{Genesis 43:28}{Genesis 43:28}\vs{28} ``Your servant our father is well,'' they replied. ``He is still alive.'' They bowed down in humility.
\everypar{}

\parshape=0
\markboth{Genesis 43:29}{Genesis 43:29}\vs{29} When Joseph looked up and saw his brother Benjamin, his mother's son, he said, ``Is this your youngest brother, whom you told me about?'' Then he said, ``May God be gracious to you, my son.''\gva{g}\marginnote{\gva{g}\,For only these two were born of Rachel.}
\markboth{Genesis 43:30}{Genesis 43:30}\vs{30} Joseph hurried out, for he was overcome by affection for his brother and was at the point of tears. So he went to his room and wept there.
\everypar{}

\parshape=0
\markboth{Genesis 43:31}{Genesis 43:31}\vs{31} Then he washed his face and came out. With composure he said, ``Set out the food.''
\markboth{Genesis 43:32}{Genesis 43:32}\vs{32} They set a place for him, a separate place for his brothers, and another for the Egyptians who were eating with him. (The Egyptians are not able to eat with Hebrews, for the Egyptians think it is disgusting to do so.)\gva{h}\marginnote{\gva{h}\,To signify his dignity.}\gva{i}\marginnote{\gva{i}\,The nature of the superstitions is to condemn all others in respect to themselves.}
\markboth{Genesis 43:33}{Genesis 43:33}\vs{33} They sat before him, arranged by order of birth, beginning with the firstborn and ending with the youngest. The men looked at each other in astonishment.
\markboth{Genesis 43:34}{Genesis 43:34}\vs{34} He gave them portions of the food set before him, but the portion for Benjamin was five times greater than the portions for any of the others. They drank with Joseph until they all became drunk.\gva{k}\marginnote{\gva{k}\,Sometimes this word means "to be drunken", but here it means that they had enough, and drank of the best wine.}

\markboth{Genesis 44:1}{Genesis 44:1}
\Needspace*{8\baselineskip}
\ch{44} \allowchapbreak\hypertarget{ch-genesis-44}{}\lettrine[lines=5]{\color{pentateuch}H}{\textsc{e}} instructed the servant who was over his household, ``Fill the sacks of the men with as much food as they can carry and put each man's money in the mouth of his sack.\everypar{}
\markboth{Genesis 44:2}{Genesis 44:2}\vs{2} Then put my cup---the silver cup---in the mouth of the youngest one's sack, along with the money for his grain.'' He did as Joseph instructed.\gva{a}\marginnote{\gva{a}\,We may not use this example to justify any unlawful practices, seeing God has commanded us to walk in simplicity.}
\\\indent\markboth{Genesis 44:3}{Genesis 44:3}\vs{3} When morning came, the men and their donkeys were sent off.
\markboth{Genesis 44:4}{Genesis 44:4}\vs{4} They had not gone very far from the city when Joseph said to the servant who was over his household, ``Pursue the men at once! When you overtake them, say to them, `Why have you repaid good with evil?
\markboth{Genesis 44:5}{Genesis 44:5}\vs{5} Doesn't my master drink from this cup and use it for divination? You have done wrong!'''\gva{b}\marginnote{\gva{b}\,Because the people thought he could divine, he attributes to himself that knowledge: or else he pretends that he consults with soothsayers: which deceit is worthy to be reproved.}
\everypar{}

\parshape=0
\markboth{Genesis 44:6}{Genesis 44:6}\vs{6} When the man overtook them, he spoke these words to them.
\markboth{Genesis 44:7}{Genesis 44:7}\vs{7} They answered him, ``Why does my lord say such things? Far be it from your servants to do such a thing!
\markboth{Genesis 44:8}{Genesis 44:8}\vs{8} Look, the money that we found in the mouths of our sacks we brought back to you from the land of Canaan. Why then would we steal silver or gold from your master's house?
\markboth{Genesis 44:9}{Genesis 44:9}\vs{9} If one of us has it, he will die, and the rest of us will become my lord's slaves!''
\everypar{}

\parshape=0
\markboth{Genesis 44:10}{Genesis 44:10}\vs{10} He replied, ``You have suggested your own punishment! The one who has it will become my slave, but the rest of you will go free.''
\markboth{Genesis 44:11}{Genesis 44:11}\vs{11} So each man quickly lowered his sack to the ground and opened it.
\markboth{Genesis 44:12}{Genesis 44:12}\vs{12} Then the man searched. He began with the oldest and finished with the youngest. The cup was found in Benjamin's sack!
\markboth{Genesis 44:13}{Genesis 44:13}\vs{13} They all tore their clothes! Then each man loaded his donkey, and they returned to the city.\gva{c}\marginnote{\gva{c}\,To show how greatly the thing displeased them and how sorry they were for it.}
\everypar{}

\parshape=0
\markboth{Genesis 44:14}{Genesis 44:14}\vs{14} So Judah and his brothers came back to Joseph's house. He was still there, and they threw themselves to the ground before him.
\markboth{Genesis 44:15}{Genesis 44:15}\vs{15} Joseph said to them, ``What did you think you were doing? Don't you know that a man like me can find out things like this by divination?''
\everypar{}

\parshape=0
\markboth{Genesis 44:16}{Genesis 44:16}\vs{16} Judah replied, ``What can we say to my lord? What can we speak? How can we clear ourselves? God has exposed the sin of your servants! We are now my lord's slaves, we and the one in whose possession the cup was found.''\gva{d}\marginnote{\gva{d}\,If we see no obvious cause for our affliction, let us look to the secret counsel of God, who punishes us justly for our sins.}
\everypar{}

\parshape=0
\markboth{Genesis 44:17}{Genesis 44:17}\vs{17} But Joseph said, ``Far be it from me to do this! The man in whose hand the cup was found will become my slave, but the rest of you may go back to your father in peace.''
\everypar{}

\parshape=0
\markboth{Genesis 44:18}{Genesis 44:18}\vs{18} Then Judah approached him and said, ``My lord, please allow your servant to speak a word with you. Please do not get angry with your servant, for you are just like Pharaoh.\gva{e}\marginnote{\gva{e}\,Equal in authority or, next to the king.}
\markboth{Genesis 44:19}{Genesis 44:19}\vs{19} My lord asked his servants, `Do you have a father or a brother?'
\markboth{Genesis 44:20}{Genesis 44:20}\vs{20} We said to my lord, `We have an aged father, and there is a young boy who was born when our father was old. The boy's brother is dead. He is the only one of his mother's sons left, and his father loves him.'
\everypar{}

\parshape=0
\markboth{Genesis 44:21}{Genesis 44:21}\vs{21} ``Then you told your servants, `Bring him down to me so I can see him.'
\markboth{Genesis 44:22}{Genesis 44:22}\vs{22} We said to my lord, `The boy cannot leave his father. If he leaves his father, his father will die.'
\markboth{Genesis 44:23}{Genesis 44:23}\vs{23} But you said to your servants, `If your youngest brother does not come down with you, you will not see my face again.'
\markboth{Genesis 44:24}{Genesis 44:24}\vs{24} When we returned to your servant my father, we told him the words of my lord.
\everypar{}

\parshape=0
\markboth{Genesis 44:25}{Genesis 44:25}\vs{25} ``Then our father said, `Go back and buy us a little food.'
\markboth{Genesis 44:26}{Genesis 44:26}\vs{26} But we replied, `We cannot go down there. If our youngest brother is with us, then we will go, for we won't be permitted to see the man's face if our youngest brother is not with us.'
\everypar{}

\parshape=0
\markboth{Genesis 44:27}{Genesis 44:27}\vs{27} ``Then your servant my father said to us, `You know that my wife gave me two sons.\gva{f}\marginnote{\gva{f}\,Rachel bore to Jacob, Joseph and Benjamin.}
\markboth{Genesis 44:28}{Genesis 44:28}\vs{28} The first disappeared and I said, ``He has surely been torn to pieces.'' I have not seen him since.
\markboth{Genesis 44:29}{Genesis 44:29}\vs{29} If you take this one from me too and an accident happens to him, then you will bring down my gray hair in tragedy to the grave.'
\everypar{}

\parshape=0
\markboth{Genesis 44:30}{Genesis 44:30}\vs{30} ``So now, when I return to your servant my father, and the boy is not with us---his very life is bound up in his son's life.
\markboth{Genesis 44:31}{Genesis 44:31}\vs{31} When he sees the boy is not with us, he will die, and your servants will bring down the gray hair of your servant our father in sorrow to the grave.
\markboth{Genesis 44:32}{Genesis 44:32}\vs{32} Indeed, your servant pledged security for the boy with my father, saying, `If I do not bring him back to you, then I will bear the blame before my father all my life.'
\everypar{}

\parshape=0
\markboth{Genesis 44:33}{Genesis 44:33}\vs{33} ``So now, please let your servant remain as my lord's slave instead of the boy. As for the boy, let him go back with his brothers.
\markboth{Genesis 44:34}{Genesis 44:34}\vs{34} For how can I go back to my father if the boy is not with me? I couldn't bear to see my father's pain.''\gva{h}\marginnote{\gva{h}\,Meaning, he would rather remain as their prisoner, than to return and see his father in sorrow.}

\markboth{Genesis 45:1}{Genesis 45:1}
\Needspace*{8\baselineskip}
\ch{45} \allowchapbreak\hypertarget{ch-genesis-45}{}\lettrine[lines=5]{\color{pentateuch}J}{\textsc{oseph}} was no longer able to control himself before all his attendants, so he cried out, ``Make everyone go out from my presence!'' No one remained with Joseph when he made himself known to his brothers.\gva{a}\marginnote{\gva{a}\,Not because he was ashamed of his kindred, but rather because he wanted to cover his brother's sin.}\everypar{}
\markboth{Genesis 45:2}{Genesis 45:2}\vs{2} He wept loudly; the Egyptians heard it and Pharaoh's household heard about it.
\everypar{}

\parshape=0
\markboth{Genesis 45:3}{Genesis 45:3}\vs{3} Joseph said to his brothers, ``I am Joseph! Is my father still alive?'' His brothers could not answer him because they were dumbfounded before him.
\markboth{Genesis 45:4}{Genesis 45:4}\vs{4} Joseph said to his brothers, ``Come closer to me,'' so they came near. Then he said, ``I am Joseph your brother, whom you sold into Egypt.
\markboth{Genesis 45:5}{Genesis 45:5}\vs{5} Now, do not be upset and do not be angry with yourselves because you sold me here, for God sent me ahead of you to preserve life!\gva{b}\marginnote{\gva{b}\,This example teaches that we must by all means comfort those who are truly ashamed and sorry for their sins.}
\markboth{Genesis 45:6}{Genesis 45:6}\vs{6} For these past two years there has been famine in the land and for five more years there will be neither plowing nor harvesting.
\markboth{Genesis 45:7}{Genesis 45:7}\vs{7} God sent me ahead of you to preserve you on the earth and to save your lives by a great deliverance.
\markboth{Genesis 45:8}{Genesis 45:8}\vs{8} So now, it is not you who sent me here, but God. He has made me an adviser to Pharaoh, lord over all his household, and ruler over all the land of Egypt.\gva{c}\marginnote{\gva{c}\,Though God detests sin, yet he turns man's wickedness into his glory.}
\markboth{Genesis 45:9}{Genesis 45:9}\vs{9} Now go up to my father quickly and tell him, `This is what your son Joseph says: ``God has made me lord of all Egypt. Come down to me; do not delay!
\markboth{Genesis 45:10}{Genesis 45:10}\vs{10} You will live in the land of Goshen, and you will be near me---you, your children, your grandchildren, your flocks, your herds, and everything you have.
\markboth{Genesis 45:11}{Genesis 45:11}\vs{11} I will provide you with food there because there will be five more years of famine. Otherwise you would become poor---you, your household, and everyone who belongs to you.'''
\markboth{Genesis 45:12}{Genesis 45:12}\vs{12} You and my brother Benjamin can certainly see with your own eyes that I really am the one who speaks to you.\gva{d}\marginnote{\gva{d}\,That is, that I speak in your own language and have no interpreter.}
\markboth{Genesis 45:13}{Genesis 45:13}\vs{13} So tell my father about all my honor in Egypt and about everything you have seen. But bring my father down here quickly!''
\everypar{}

\parshape=0
\markboth{Genesis 45:14}{Genesis 45:14}\vs{14} Then he threw himself on the neck of his brother Benjamin and wept, and Benjamin wept on his neck.
\markboth{Genesis 45:15}{Genesis 45:15}\vs{15} He kissed all his brothers and wept over them. After this his brothers talked with him.
\everypar{}

\parshape=0
\markboth{Genesis 45:16}{Genesis 45:16}\vs{16} Now it was reported in the household of Pharaoh, ``Joseph's brothers have arrived.'' It pleased Pharaoh and his servants.
\markboth{Genesis 45:17}{Genesis 45:17}\vs{17} Pharaoh said to Joseph, ``Say to your brothers, `Do this: Load your animals and go to the land of Canaan!
\markboth{Genesis 45:18}{Genesis 45:18}\vs{18} Get your father and your households and come to me! Then I will give you the best land in Egypt and you will eat the best of the land.'\gva{e}\marginnote{\gva{e}\,The most plentiful ground.}\gva{f}\marginnote{\gva{f}\,The main fruits and conveniences.}
\markboth{Genesis 45:19}{Genesis 45:19}\vs{19} You are also commanded to say, `Do this: Take for yourselves wagons from the land of Egypt for your little ones and for your wives. Bring your father and come.
\markboth{Genesis 45:20}{Genesis 45:20}\vs{20} Don't worry about your belongings, for the best of all the land of Egypt will be yours.'''
\everypar{}

\parshape=0
\markboth{Genesis 45:21}{Genesis 45:21}\vs{21} So the sons of Israel did as he said. Joseph gave them wagons as Pharaoh had instructed, and he gave them provisions for the journey.
\markboth{Genesis 45:22}{Genesis 45:22}\vs{22} He gave sets of clothes to each one of them, but to Benjamin he gave 300 pieces of silver and five sets of clothes.
\markboth{Genesis 45:23}{Genesis 45:23}\vs{23} To his father he sent the following: ten donkeys loaded with the best products of Egypt and ten female donkeys loaded with grain, food, and provisions for his father's journey.
\markboth{Genesis 45:24}{Genesis 45:24}\vs{24} Then he sent his brothers on their way and they left. He said to them, ``As you travel don't be overcome with fear.''\gva{g}\marginnote{\gva{g}\,Seeing he had remitted the fault done to him, he did not want them to accuse one another.}
\everypar{}

\parshape=0
\markboth{Genesis 45:25}{Genesis 45:25}\vs{25} So they went up from Egypt and came to their father Jacob in the land of Canaan.
\markboth{Genesis 45:26}{Genesis 45:26}\vs{26} They told him, ``Joseph is still alive and he is ruler over all the land of Egypt!'' Jacob was stunned, for he did not believe them.\gva{h}\marginnote{\gva{h}\,As one between hope and fear.}
\markboth{Genesis 45:27}{Genesis 45:27}\vs{27} But when they related to him everything Joseph had said to them, and when he saw the wagons that Joseph had sent to transport him, their father Jacob's spirit revived.
\markboth{Genesis 45:28}{Genesis 45:28}\vs{28} Then Israel said, ``Enough! My son Joseph is still alive! I will go and see him before I die.''

\markboth{Genesis 46:1}{Genesis 46:1}
\Needspace*{8\baselineskip}
\ch{46} \allowchapbreak\hypertarget{ch-genesis-46}{}\lettrine[lines=5]{\color{pentateuch}S}{\textsc{o}} Israel began his journey, taking with him all that he had. When he came to Beer Sheba he offered sacrifices to the God of his father Isaac.\gva{a}\marginnote{\gva{a}\,By this he signified both that he worshipped the true God, and that he kept in his heart the possession of that land from which need drove him at that time.}\everypar{}
\markboth{Genesis 46:2}{Genesis 46:2}\vs{2} God spoke to Israel in a vision during the night and said, ``Jacob, Jacob!'' He replied, ``Here I am!''
\markboth{Genesis 46:3}{Genesis 46:3}\vs{3} He said, ``I am God, the God of your father. Do not be afraid to go down to Egypt, for I will make you into a great nation there.
\markboth{Genesis 46:4}{Genesis 46:4}\vs{4} I will go down with you to Egypt and I myself will certainly bring you back from there. Joseph will close your eyes.''\gva{b}\marginnote{\gva{b}\,Conducting you by my power.}\gva{c}\marginnote{\gva{c}\,In your posterity.}\gva{d}\marginnote{\gva{d}\,Shall shut your eyes when you die: which belongs to him that was most dear or chief of the kindred.}
\everypar{}

\parshape=0
\markboth{Genesis 46:5}{Genesis 46:5}\vs{5} Then Jacob started out from Beer Sheba, and the sons of Israel carried their father Jacob, their little children, and their wives in the wagons that Pharaoh had sent along to transport him.
\markboth{Genesis 46:6}{Genesis 46:6}\vs{6} Jacob and all his descendants took their livestock and the possessions they had acquired in the land of Canaan, and they went to Egypt.
\markboth{Genesis 46:7}{Genesis 46:7}\vs{7} He brought with him to Egypt his sons and grandsons, his daughters and granddaughters---all his descendants.
\everypar{}

\parshape=0
\markboth{Genesis 46:8}{Genesis 46:8}\vs{8} These are the names of the sons of Israel who went to Egypt---Jacob and his sons: Reuben, the firstborn of Jacob.
\everypar{}

\parshape=0
\markboth{Genesis 46:9}{Genesis 46:9}\vs{9} The sons of Reuben: Hanoch, Pallu, Hezron, and Carmi.
\everypar{}

\parshape=0
\markboth{Genesis 46:10}{Genesis 46:10}\vs{10} The sons of Simeon: Jemuel, Jamin, Ohad, Jakin, Zohar, and Shaul (the son of a Canaanite woman).
\everypar{}

\parshape=0
\markboth{Genesis 46:11}{Genesis 46:11}\vs{11} The sons of Levi: Gershon, Kohath, and Merari.
\everypar{}

\parshape=0
\markboth{Genesis 46:12}{Genesis 46:12}\vs{12} The sons of Judah: Er, Onan, Shelah, Perez, and Zerah (but Er and Onan died in the land of Canaan). The sons of Perez were Hezron and Hamul.
\everypar{}

\parshape=0
\markboth{Genesis 46:13}{Genesis 46:13}\vs{13} The sons of Issachar: Tola, Puah, Jashub, and Shimron.
\everypar{}

\parshape=0
\markboth{Genesis 46:14}{Genesis 46:14}\vs{14} The sons of Zebulun: Sered, Elon, and Jahleel.
\everypar{}

\parshape=0
\markboth{Genesis 46:15}{Genesis 46:15}\vs{15} These were the sons of Leah, whom she bore to Jacob in Paddan Aram, along with Dinah his daughter. His sons and daughters numbered thirty-three in all.
\everypar{}

\parshape=0
\markboth{Genesis 46:16}{Genesis 46:16}\vs{16} The sons of Gad: Zephon, Haggi, Shuni, Ezbon, Eri, Arodi, and Areli.
\everypar{}

\parshape=0
\markboth{Genesis 46:17}{Genesis 46:17}\vs{17} The sons of Asher: Imnah, Ishvah, Ishvi, Beriah, and Serah their sister. The sons of Beriah were Heber and Malkiel.
\everypar{}

\parshape=0
\markboth{Genesis 46:18}{Genesis 46:18}\vs{18} These were the sons of Zilpah, whom Laban gave to Leah his daughter. She bore these to Jacob, sixteen in all.
\everypar{}

\parshape=0
\markboth{Genesis 46:19}{Genesis 46:19}\vs{19} The sons of Rachel the wife of Jacob: Joseph and Benjamin.
\everypar{}

\parshape=0
\markboth{Genesis 46:20}{Genesis 46:20}\vs{20} Manasseh and Ephraim were born to Joseph in the land of Egypt. Asenath daughter of Potiphera, priest of On, bore them to him.
\everypar{}

\parshape=0
\markboth{Genesis 46:21}{Genesis 46:21}\vs{21} The sons of Benjamin: Bela, Beker, Ashbel, Gera, Naaman, Ehi, Rosh, Muppim, Huppim and Ard.
\everypar{}

\parshape=0
\markboth{Genesis 46:22}{Genesis 46:22}\vs{22} These were the sons of Rachel who were born to Jacob, fourteen in all.
\everypar{}

\parshape=0
\markboth{Genesis 46:23}{Genesis 46:23}\vs{23} The son of Dan: Hushim.
\everypar{}

\parshape=0
\markboth{Genesis 46:24}{Genesis 46:24}\vs{24} The sons of Naphtali: Jahziel, Guni, Jezer, and Shillem.
\everypar{}

\parshape=0
\markboth{Genesis 46:25}{Genesis 46:25}\vs{25} These were the sons of Bilhah, whom Laban gave to Rachel his daughter. She bore these to Jacob, seven in all.
\everypar{}

\parshape=0
\markboth{Genesis 46:26}{Genesis 46:26}\vs{26} All the direct descendants of Jacob who went to Egypt with him were sixty-six in number. (This number does not include the wives of Jacob's sons.)
\markboth{Genesis 46:27}{Genesis 46:27}\vs{27} Counting the two sons of Joseph who were born to him in Egypt, all the people of the household of Jacob who were in Egypt numbered seventy.
\everypar{}

\parshape=0
\markboth{Genesis 46:28}{Genesis 46:28}\vs{28} Jacob sent Judah before him to Joseph to accompany him to Goshen. So they came to the land of Goshen.
\markboth{Genesis 46:29}{Genesis 46:29}\vs{29} Joseph harnessed his chariot and went up to meet his father Israel in Goshen. When he met him, he hugged his neck and wept on his neck for quite some time.
\everypar{}

\parshape=0
\markboth{Genesis 46:30}{Genesis 46:30}\vs{30} Israel said to Joseph, ``Now let me die since I have seen your face and know that you are still alive.''
\markboth{Genesis 46:31}{Genesis 46:31}\vs{31} Then Joseph said to his brothers and his father's household, ``I will go up and tell Pharaoh, `My brothers and my father's household who were in the land of Canaan have come to me.
\markboth{Genesis 46:32}{Genesis 46:32}\vs{32} The men are shepherds; they take care of livestock. They have brought their flocks and their herds and all that they have.'\gva{e}\marginnote{\gva{e}\,He was not ashamed of his father and kindred, though they were of base condition.}
\markboth{Genesis 46:33}{Genesis 46:33}\vs{33} Pharaoh will summon you and say, `What is your occupation?'
\markboth{Genesis 46:34}{Genesis 46:34}\vs{34} Tell him, `Your servants have taken care of cattle from our youth until now, both we and our fathers,' so that you may live in the land of Goshen, for everyone who takes care of sheep is disgusting to the Egyptians.''\gva{f}\marginnote{\gva{f}\,God permits the world to hate his own, so they will forsake the filth of the world, and cling to him.}

\markboth{Genesis 47:1}{Genesis 47:1}
\Needspace*{8\baselineskip}
\ch{47} \allowchapbreak\hypertarget{ch-genesis-47}{}\lettrine[lines=5]{\color{pentateuch}J}{\textsc{oseph}} went and told Pharaoh, ``My father, my brothers, their flocks and herds, and all that they own have arrived from the land of Canaan. They are now in the land of Goshen.''\everypar{}
\markboth{Genesis 47:2}{Genesis 47:2}\vs{2} He took five of his brothers and introduced them to Pharaoh.\gva{a}\marginnote{\gva{a}\,That the king might be assured that they had come, and to see what type of people they were.}
\\\indent\markboth{Genesis 47:3}{Genesis 47:3}\vs{3} Pharaoh said to Joseph's brothers, ``What is your occupation?'' They said to Pharaoh, ``Your servants take care of flocks, just as our ancestors did.''
\markboth{Genesis 47:4}{Genesis 47:4}\vs{4} Then they said to Pharaoh, ``We have come to live as temporary residents in the land. There is no pasture for your servants' flocks because the famine is severe in the land of Canaan. So now, please let your servants live in the land of Goshen.''
\everypar{}

\parshape=0
\markboth{Genesis 47:5}{Genesis 47:5}\vs{5} Pharaoh said to Joseph, ``Your father and your brothers have come to you.
\markboth{Genesis 47:6}{Genesis 47:6}\vs{6} The land of Egypt is before you; settle your father and your brothers in the best region of the land. They may live in the land of Goshen. If you know of any highly capable men among them, put them in charge of my livestock.''\gva{b}\marginnote{\gva{b}\,Joseph's great modesty appears in that he would attempt nothing without the king's commandment.}
\everypar{}

\parshape=0
\markboth{Genesis 47:7}{Genesis 47:7}\vs{7} Then Joseph brought in his father Jacob and presented him before Pharaoh. Jacob blessed Pharaoh.
\markboth{Genesis 47:8}{Genesis 47:8}\vs{8} Pharaoh said to Jacob, ``How long have you lived?''
\markboth{Genesis 47:9}{Genesis 47:9}\vs{9} Jacob said to Pharaoh, ``All the years of my travels are 130. All the years of my life have been few and painful; the years of my travels are not as long as those of my ancestors.''
\markboth{Genesis 47:10}{Genesis 47:10}\vs{10} Then Jacob blessed Pharaoh and went out from his presence.
\everypar{}

\parshape=0
\markboth{Genesis 47:11}{Genesis 47:11}\vs{11} So Joseph settled his father and his brothers. He gave them territory in the land of Egypt, in the best region of the land, the land of Rameses, just as Pharaoh had commanded.\gva{c}\marginnote{\gva{c}\,Which was a city in the country of Goshen, Exodus 1:11 .}
\markboth{Genesis 47:12}{Genesis 47:12}\vs{12} Joseph also provided food for his father, his brothers, and all his father's household, according to the number of their little children.\gva{d}\marginnote{\gva{d}\,Some read that he fed them as little babies, because they could not provide for themselves against that famine.}
\everypar{}

\parshape=0
\markboth{Genesis 47:13}{Genesis 47:13}\vs{13} But there was no food in all the land because the famine was very severe; the land of Egypt and the land of Canaan wasted away because of the famine.
\markboth{Genesis 47:14}{Genesis 47:14}\vs{14} Joseph collected all the money that could be found in the land of Egypt and in the land of Canaan as payment for the grain they were buying. Then Joseph brought the money into Pharaoh's palace.\gva{e}\marginnote{\gva{e}\,In which he both declares his faithfulness to the king, and his freedom from covetousness.}
\markboth{Genesis 47:15}{Genesis 47:15}\vs{15} When the money from the lands of Egypt and Canaan was used up, all the Egyptians came to Joseph and said, ``Give us food! Why should we die before your very eyes because our money has run out?''
\everypar{}

\parshape=0
\markboth{Genesis 47:16}{Genesis 47:16}\vs{16} Then Joseph said, ``If your money is gone, bring your livestock, and I will give you food in exchange for your livestock.''
\markboth{Genesis 47:17}{Genesis 47:17}\vs{17} So they brought their livestock to Joseph, and Joseph gave them food in exchange for their horses, the livestock of their flocks and herds, and their donkeys. He got them through that year by giving them food in exchange for all their livestock.
\everypar{}

\parshape=0
\markboth{Genesis 47:18}{Genesis 47:18}\vs{18} When that year was over, they came to him the next year and said to him, ``We cannot hide from our lord that the money is used up and the livestock and the animals belong to our lord. Nothing remains before our lord except our bodies and our land.
\markboth{Genesis 47:19}{Genesis 47:19}\vs{19} Why should we die before your very eyes, both we and our land? Buy us and our land in exchange for food, and we, with our land, will become Pharaoh's slaves. Give us seed that we may live and not die. Then the land will not become desolate.''\gva{f}\marginnote{\gva{f}\,For unless the ground is tilled and sown, it perishes and is as if it was dead.}
\everypar{}

\parshape=0
\markboth{Genesis 47:20}{Genesis 47:20}\vs{20} So Joseph bought all the land of Egypt for Pharaoh. Each of the Egyptians sold his field, for the famine was severe. So the land became Pharaoh's.
\markboth{Genesis 47:21}{Genesis 47:21}\vs{21} Joseph made all the people slaves from one end of Egypt's border to the other end of it.\gva{g}\marginnote{\gva{g}\,By this changing they signified that they had nothing of their own, but received everything from the king's generosity.}
\markboth{Genesis 47:22}{Genesis 47:22}\vs{22} But he did not purchase the land of the priests because the priests had an allotment from Pharaoh and they ate from their allotment that Pharaoh gave them. That is why they did not sell their land.
\everypar{}

\parshape=0
\markboth{Genesis 47:23}{Genesis 47:23}\vs{23} Joseph said to the people, ``Since I have bought you and your land today for Pharaoh, here is seed for you. Cultivate the land.
\markboth{Genesis 47:24}{Genesis 47:24}\vs{24} When the crop comes in, give one-fifth of it to Pharaoh. The remaining four-fifths will be yours for seed for the fields and for you to eat, including those in your households and your little children.''
\markboth{Genesis 47:25}{Genesis 47:25}\vs{25} They replied, ``You have saved our lives! You are showing us favor, and we will be Pharaoh's slaves.''
\everypar{}

\parshape=0
\markboth{Genesis 47:26}{Genesis 47:26}\vs{26} So Joseph made it a statute, which is in effect to this day throughout the land of Egypt: One-fifth belongs to Pharaoh. Only the land of the priests did not become Pharaoh's.\gva{h}\marginnote{\gva{h}\,Pharaoh, in providing for idolatrous priests, will be a condemnation to all those who neglect the true ministers of God's word.}
\everypar{}

\parshape=0
\markboth{Genesis 47:27}{Genesis 47:27}\vs{27} Israel settled in the land of Egypt, in the land of Goshen, and they owned land there. They were fruitful and increased rapidly in number.
\everypar{}

\parshape=0
\markboth{Genesis 47:28}{Genesis 47:28}\vs{28} Jacob lived in the land of Egypt 17 years; the years of Jacob's life were 147 in all.
\markboth{Genesis 47:29}{Genesis 47:29}\vs{29} The time for Israel to die approached, so he called for his son Joseph and said to him, ``If now I have found favor in your sight, put your hand under my thigh and show me kindness and faithfulness. Do not bury me in Egypt,
\markboth{Genesis 47:30}{Genesis 47:30}\vs{30} but when I rest with my fathers, carry me out of Egypt and bury me in their burial place.'' Joseph said, ``I will do as you say.''\gva{i}\marginnote{\gva{i}\,By this he demonstrated that he died in the faith of his fathers, teaching his children to hope for the promised land.}
\everypar{}

\parshape=0
\markboth{Genesis 47:31}{Genesis 47:31}\vs{31} Jacob said, ``Swear to me that you will do so.'' So Joseph gave him his word. Then Israel bowed down at the head of his bed.\gva{k}\marginnote{\gva{k}\,He rejoiced that Joseph had promised him, and setting himself up on his pillows, praised God; 1 Chronicles 29:10 .}

\markboth{Genesis 48:1}{Genesis 48:1}
\Needspace*{8\baselineskip}
\ch{48} \allowchapbreak\hypertarget{ch-genesis-48}{}\lettrine[lines=5]{\color{pentateuch}A}{\textsc{fter}} these things Joseph was told, ``Your father is weakening.'' So he took his two sons Manasseh and Ephraim with him.\gva{a}\marginnote{\gva{a}\,Joseph valued his children being received into Jacob's family, which was the Church of God, more than enjoying all the treasures of Egypt.}\everypar{}
\markboth{Genesis 48:2}{Genesis 48:2}\vs{2} When Jacob was told, ``Your son Joseph has just come to you,'' Israel regained strength and sat up on his bed.
\markboth{Genesis 48:3}{Genesis 48:3}\vs{3} Jacob said to Joseph, ``The Sovereign God appeared to me at Luz in the land of Canaan and blessed me.
\markboth{Genesis 48:4}{Genesis 48:4}\vs{4} He said to me, `I am going to make you fruitful and will multiply you. I will make you into a group of nations, and I will give this land to your descendants as an everlasting possession.'\gva{b}\marginnote{\gva{b}\,Which is true in the carnal Israel until the coming of Christ, and in the spiritual forever.}
\everypar{}

\parshape=0
\markboth{Genesis 48:5}{Genesis 48:5}\vs{5} ``Now, as for your two sons, who were born to you in the land of Egypt before I came to you in Egypt, they will be mine. Ephraim and Manasseh will be mine just as Reuben and Simeon are.
\markboth{Genesis 48:6}{Genesis 48:6}\vs{6} Any children that you father after them will be yours; they will be listed under the names of their brothers in their inheritance.
\markboth{Genesis 48:7}{Genesis 48:7}\vs{7} But as for me, when I was returning from Paddan, Rachel died---to my sorrow---in the land of Canaan. It happened along the way, some distance from Ephrath. So I buried her there on the way to Ephrath'' (that is, Bethlehem).
\everypar{}

\parshape=0
\markboth{Genesis 48:8}{Genesis 48:8}\vs{8} When Israel saw Joseph's sons, he asked, ``Who are these?''
\markboth{Genesis 48:9}{Genesis 48:9}\vs{9} Joseph said to his father, ``They are the sons God has given me in this place.'' His father said, ``Bring them to me so I may bless them.''\gva{c}\marginnote{\gva{c}\,The faithful acknowledge all benefits come from God's free mercy.}
\markboth{Genesis 48:10}{Genesis 48:10}\vs{10} Now Israel's eyes were failing because of his age; he was not able to see well. So Joseph brought his sons near to him, and his father kissed them and embraced them.
\markboth{Genesis 48:11}{Genesis 48:11}\vs{11} Israel said to Joseph, ``I never expected to see you again, but now God has allowed me to see your children too.''
\everypar{}

\parshape=0
\markboth{Genesis 48:12}{Genesis 48:12}\vs{12} So Joseph moved them from Israel's knees and bowed down with his face to the ground.
\markboth{Genesis 48:13}{Genesis 48:13}\vs{13} Joseph positioned them; he put Ephraim on his right hand across from Israel's left hand, and Manasseh on his left hand across from Israel's right hand. Then Joseph brought them closer to his father.
\markboth{Genesis 48:14}{Genesis 48:14}\vs{14} Israel stretched out his right hand and placed it on Ephraim's head, although he was the younger. Crossing his hands, he put his left hand on Manasseh's head, for Manasseh was the firstborn.\gva{d}\marginnote{\gva{d}\,God's judgments are often contrary to man's and he prefers what man despises.}
\everypar{}

\parshape=0
\markboth{Genesis 48:15}{Genesis 48:15}\vs{15} Then he blessed Joseph and said, ``May the God before whom my fathers Abraham and Isaac walked--- the God who has been my shepherd all my life long to this day,
\everypar{}

\parshape=0
\markboth{Genesis 48:16}{Genesis 48:16}\vs{16} the angel who has protected me from all harm--- bless these boys. May my name be named in them, and the name of my fathers Abraham and Isaac. May they grow into a multitude on the earth.''\gva{e}\marginnote{\gva{e}\,This angel must be understood to be Christ, as in Genesis 31:13 ; Genesis 32:1 .}\gva{f}\marginnote{\gva{f}\,Let them be taken as my children.}
\everypar{}

\parshape=0
\markboth{Genesis 48:17}{Genesis 48:17}\vs{17} When Joseph saw that his father placed his right hand on Ephraim's head, it displeased him. So he took his father's hand to move it from Ephraim's head to Manasseh's head.\gva{g}\marginnote{\gva{g}\,Joseph fails by binding God's grace to the order of nature.}
\markboth{Genesis 48:18}{Genesis 48:18}\vs{18} Joseph said to his father, ``Not so, my father, for this is the firstborn. Put your right hand on his head.''
\everypar{}

\parshape=0
\markboth{Genesis 48:19}{Genesis 48:19}\vs{19} But his father refused and said, ``I know, my son, I know. He too will become a nation and he too will become great. In spite of this, his younger brother will be even greater and his descendants will become a multitude of nations.''
\everypar{}

\parshape=0
\markboth{Genesis 48:20}{Genesis 48:20}\vs{20} So he blessed them that day, saying, ``By you will Israel bless, saying, `May God make you like Ephraim and Manasseh.''' Thus he put Ephraim before Manasseh.\gva{h}\marginnote{\gva{h}\,In whom God's graces should manifestly appear.}
\everypar{}

\parshape=0
\markboth{Genesis 48:21}{Genesis 48:21}\vs{21} Then Israel said to Joseph, ``I am about to die, but God will be with you and will bring you back to the land of your fathers.\gva{i}\marginnote{\gva{i}\,Which they had by faith in the promise.}
\markboth{Genesis 48:22}{Genesis 48:22}\vs{22} As one who is above your brothers, I give to you the mountain slope, which I took from the Amorites with my sword and my bow.''\gva{k}\marginnote{\gva{k}\,By my children whom God spared for my sake.}

\markboth{Genesis 49:1}{Genesis 49:1}
\Needspace*{8\baselineskip}
\ch{49} \allowchapbreak\hypertarget{ch-genesis-49}{}\lettrine[lines=5]{\color{pentateuch}J}{\textsc{acob}} called for his sons and said, ``Gather together so I can tell you what will happen to you in future days.\gva{a}\marginnote{\gva{a}\,last days.}\gva{a}\marginnote{\gva{a}\,When God will bring you out of Egypt, and because he speaks of the Messiah, he calls it the last days.}\everypar{}
\\\indent\markboth{Genesis 49:2}{Genesis 49:2}\vs{2} ``Assemble and listen, you sons of Jacob; listen to Israel, your father.
\\\indent\markboth{Genesis 49:3}{Genesis 49:3}\vs{3} Reuben, you are my firstborn, my might and the beginning of my strength, outstanding in dignity, outstanding in power.\gva{b}\marginnote{\gva{b}\,Begotten in my youth.}\gva{c}\marginnote{\gva{c}\,If you have not left your birthright by your offence.}
\everypar{}

\parshape=0
\markboth{Genesis 49:4}{Genesis 49:4}\vs{4} You are destructive like water and will not excel, for you got on your father's bed, then you defiled it---he got on my couch!
\everypar{}

\parshape=0
\markboth{Genesis 49:5}{Genesis 49:5}\vs{5} Simeon and Levi are brothers, weapons of violence are their knives!
\everypar{}

\parshape=0
\markboth{Genesis 49:6}{Genesis 49:6}\vs{6} O my soul, do not come into their council, do not be united to their assembly, my heart, for in their anger they have killed men, and for pleasure they have hamstrung oxen.\gva{d}\marginnote{\gva{d}\,Or, tongue: meaning that he neither consented to them in word or thought.}\gva{e}\marginnote{\gva{e}\,The Shechemites Genesis 34:26 .}
\everypar{}

\parshape=0
\markboth{Genesis 49:7}{Genesis 49:7}\vs{7} Cursed be their anger, for it was fierce, and their fury, for it was cruel. I will divide them in Jacob, and scatter them in Israel!\gva{f}\marginnote{\gva{f}\,For Levi had no part, and Simeon was under Judah, Joshua 19:1 till God gave them the place of the Amalekites, 1 Chronicles 4:43 .}
\everypar{}

\parshape=0
\markboth{Genesis 49:8}{Genesis 49:8}\vs{8} Judah, your brothers will praise you. Your hand will be on the neck of your enemies, your father's sons will bow down before you.\gva{g}\marginnote{\gva{g}\,As was verified in David and Christ.}
\everypar{}

\parshape=0
\markboth{Genesis 49:9}{Genesis 49:9}\vs{9} You are a lion's cub, Judah, from the prey, my son, you have gone up. He crouches and lies down like a lion; like a lioness---who will rouse him?\gva{h}\marginnote{\gva{h}\,His enemies will so fear him.}
\everypar{}

\parshape=0
\markboth{Genesis 49:10}{Genesis 49:10}\vs{10} The scepter will not depart from Judah, nor the ruler's staff from between his feet, until he comes to whom it belongs; the nations will obey him.\gva{i}\marginnote{\gva{i}\,Which is Christ the Messiah, the giver of prosperity who will call the Gentiles to salvation.}
\everypar{}

\parshape=0
\markboth{Genesis 49:11}{Genesis 49:11}\vs{11} Binding his foal to the vine, and his colt to the choicest vine, he will wash his garments in wine, his robes in the blood of grapes.\gva{k}\marginnote{\gva{k}\,A country most abundant with vines and pastures is promised to him.}
\everypar{}

\parshape=0
\markboth{Genesis 49:12}{Genesis 49:12}\vs{12} His eyes will be red from wine, and his teeth white from milk.
\everypar{}

\parshape=0
\markboth{Genesis 49:13}{Genesis 49:13}\vs{13} Zebulun will live by the haven of the sea and become a haven for ships; his border will extend to Sidon.
\everypar{}

\parshape=0
\markboth{Genesis 49:14}{Genesis 49:14}\vs{14} Issachar is a strong-boned donkey lying down between two saddlebags.\gva{l}\marginnote{\gva{l}\,His force will be great, but he will lack courage to resist his enemies.}
\everypar{}

\parshape=0
\markboth{Genesis 49:15}{Genesis 49:15}\vs{15} When he sees a good resting place, and the pleasant land, he will bend his shoulder to the burden and become a slave laborer.
\everypar{}

\parshape=0
\markboth{Genesis 49:16}{Genesis 49:16}\vs{16} Dan will judge his people as one of the tribes of Israel.\gva{m}\marginnote{\gva{m}\,Shall have the honour of a tribe.}
\everypar{}

\parshape=0
\markboth{Genesis 49:17}{Genesis 49:17}\vs{17} May Dan be a snake beside the road, a viper by the path, that bites the heels of the horse so that its rider falls backward.
\everypar{}

\parshape=0
\markboth{Genesis 49:18}{Genesis 49:18}\vs{18} I wait for your deliverance, O Lord.\gva{o}\marginnote{\gva{o}\,Seeing the miseries that his posterity would fall into, he bursts out in prayer to God to remedy it.}
\everypar{}

\parshape=0
\markboth{Genesis 49:19}{Genesis 49:19}\vs{19} Gad will be raided by marauding bands, but he will attack them at their heels.
\everypar{}

\parshape=0
\markboth{Genesis 49:20}{Genesis 49:20}\vs{20} Asher's food will be rich, and he will provide delicacies to royalty.\gva{p}\marginnote{\gva{p}\,He will abound in corn and pleasant fruits.}
\everypar{}

\parshape=0
\markboth{Genesis 49:21}{Genesis 49:21}\vs{21} Naphtali is a free running doe, he speaks delightful words.\gva{q}\marginnote{\gva{q}\,Overcoming more by fair words than by force.}
\everypar{}

\parshape=0
\markboth{Genesis 49:22}{Genesis 49:22}\vs{22} Joseph is a fruitful bough, a fruitful bough near a spring whose branches climb over the wall.
\everypar{}

\parshape=0
\markboth{Genesis 49:23}{Genesis 49:23}\vs{23} The archers will attack him, they will shoot at him and oppose him.\gva{r}\marginnote{\gva{r}\,As his brethren when they were his enemies, Potiphar and others.}
\everypar{}

\parshape=0
\markboth{Genesis 49:24}{Genesis 49:24}\vs{24} But his bow will remain steady, and his hands will be skillful; because of the hands of the Powerful One of Jacob, because of the Shepherd, the Rock of Israel,\gva{s}\marginnote{\gva{s}\,That is God.}
\everypar{}

\parshape=0
\markboth{Genesis 49:25}{Genesis 49:25}\vs{25} because of the God of your father, who will help you, because of the Sovereign God, who will bless you with blessings from the sky above, blessings from the deep that lies below, and blessings of the breasts and womb.
\everypar{}

\parshape=0
\markboth{Genesis 49:26}{Genesis 49:26}\vs{26} The blessings of your father are greater than the blessings of the eternal mountains or the desirable things of the age-old hills. They will be on the head of Joseph and on the brow of the prince of his brothers.\gva{t}\marginnote{\gva{t}\,In as much as he was closer to the accomplishment of the promise and it had been more often confirmed.}\gva{u}\marginnote{\gva{u}\,Either in dignity, or when he was sold from his brethren.}
\everypar{}

\parshape=0
\markboth{Genesis 49:27}{Genesis 49:27}\vs{27} Benjamin is a ravenous wolf; in the morning devouring the prey, and in the evening dividing the plunder.''
\everypar{}

\parshape=0
\markboth{Genesis 49:28}{Genesis 49:28}\vs{28} These are the twelve tribes of Israel. This is what their father said to them when he blessed them. He gave each of them an appropriate blessing.
\everypar{}

\parshape=0
\markboth{Genesis 49:29}{Genesis 49:29}\vs{29} Then he instructed them, ``I am about to go to my people. Bury me with my fathers in the cave in the field of Ephron the Hittite.
\markboth{Genesis 49:30}{Genesis 49:30}\vs{30} It is the cave in the field of Machpelah, near Mamre in the land of Canaan, which Abraham bought for a burial plot from Ephron the Hittite.
\markboth{Genesis 49:31}{Genesis 49:31}\vs{31} There they buried Abraham and his wife Sarah; there they buried Isaac and his wife Rebekah; and there I buried Leah.
\markboth{Genesis 49:32}{Genesis 49:32}\vs{32} The field and the cave in it were acquired from the sons of Heth.''
\everypar{}

\parshape=0
\markboth{Genesis 49:33}{Genesis 49:33}\vs{33} When Jacob finished giving these instructions to his sons, he pulled his feet up onto the bed, breathed his last breath, and went to his people.\gva{x}\marginnote{\gva{x}\,By which is signified how quietly he died.}

\markboth{Genesis 50:1}{Genesis 50:1}
\Needspace*{8\baselineskip}
\ch{50} \allowchapbreak\hypertarget{ch-genesis-50}{}\lettrine[lines=5]{\color{pentateuch}T}{\textsc{hen}} Joseph hugged his father's face. He wept over him and kissed him.\everypar{}
\markboth{Genesis 50:2}{Genesis 50:2}\vs{2} Joseph instructed the physicians in his service to embalm his father, so the physicians embalmed Israel.\gva{a}\marginnote{\gva{a}\,He means those who embalmed the dead and buried them.}
\markboth{Genesis 50:3}{Genesis 50:3}\vs{3} They took 40 days, for that is the full time needed for embalming. The Egyptians mourned for him 70 days.\gva{b}\marginnote{\gva{b}\,They were more excessive in lamenting than the faithful.}
\everypar{}

\parshape=0
\markboth{Genesis 50:4}{Genesis 50:4}\vs{4} When the days of mourning had passed, Joseph said to Pharaoh's royal court, ``If I have found favor in your sight, please say to Pharaoh,
\markboth{Genesis 50:5}{Genesis 50:5}\vs{5} `My father made me swear an oath. He said, ``I am about to die. Bury me in my tomb that I dug for myself there in the land of Canaan.'' Now let me go and bury my father; then I will return.'''
\markboth{Genesis 50:6}{Genesis 50:6}\vs{6} So Pharaoh said, ``Go and bury your father, just as he made you swear to do.''\gva{c}\marginnote{\gva{c}\,Even the infidels would have oaths carried out.}
\everypar{}

\parshape=0
\markboth{Genesis 50:7}{Genesis 50:7}\vs{7} So Joseph went up to bury his father; all Pharaoh's officials went with him---the senior courtiers of his household, all the senior officials of the land of Egypt,
\markboth{Genesis 50:8}{Genesis 50:8}\vs{8} all Joseph's household, his brothers, and his father's household. But they left their little children and their flocks and herds in the land of Goshen.
\markboth{Genesis 50:9}{Genesis 50:9}\vs{9} Chariots and horsemen also went up with him, so it was a very large entourage.
\everypar{}

\parshape=0
\markboth{Genesis 50:10}{Genesis 50:10}\vs{10} When they came to the threshing floor of Atad on the other side of the Jordan, they mourned there with very great and bitter sorrow. There Joseph observed a seven-day period of mourning for his father.
\markboth{Genesis 50:11}{Genesis 50:11}\vs{11} When the Canaanites who lived in the land saw them mourning at the threshing floor of Atad, they said, ``This is a very sad occasion for the Egyptians.'' That is why its name was called Abel Mizraim, which is beyond the Jordan.
\everypar{}

\parshape=0
\markboth{Genesis 50:12}{Genesis 50:12}\vs{12} So the sons of Jacob did for him just as he had instructed them.
\markboth{Genesis 50:13}{Genesis 50:13}\vs{13} His sons carried him to the land of Canaan and buried him in the cave of the field of Machpelah, near Mamre. This is the field Abraham purchased as a burial plot from Ephron the Hittite.
\markboth{Genesis 50:14}{Genesis 50:14}\vs{14} After he buried his father, Joseph returned to Egypt, along with his brothers and all who had accompanied him to bury his father.
\everypar{}

\parshape=0
\markboth{Genesis 50:15}{Genesis 50:15}\vs{15} When Joseph's brothers saw that their father was dead, they said, ``What if Joseph bears a grudge and wants to repay us in full for all the harm we did to him?''\gva{d}\marginnote{\gva{d}\,An evil conscience is never fully at rest.}
\markboth{Genesis 50:16}{Genesis 50:16}\vs{16} So they sent word to Joseph, saying, ``Your father gave these instructions before he died:
\markboth{Genesis 50:17}{Genesis 50:17}\vs{17} `Tell Joseph this: Please forgive the sin of your brothers and the wrong they did when they treated you so badly.' Now please forgive the sin of the servants of the God of your father.'' When this message was reported to him, Joseph wept.\gva{e}\marginnote{\gva{e}\,Meaning, that they who have one God should be joined in most sure love.}
\markboth{Genesis 50:18}{Genesis 50:18}\vs{18} Then his brothers also came and threw themselves down before him; they said, ``Here we are; we are your slaves.''
\markboth{Genesis 50:19}{Genesis 50:19}\vs{19} But Joseph answered them, ``Don't be afraid. Am I in the place of God?\gva{s}\marginnote{\gva{s}\,Who by the good success seems to remit it, and therefore it should not be revenged by me.}
\markboth{Genesis 50:20}{Genesis 50:20}\vs{20} As for you, you meant to harm me, but God intended it for a good purpose, so he could preserve the lives of many people, as you can see this day.
\markboth{Genesis 50:21}{Genesis 50:21}\vs{21} So now, don't be afraid. I will provide for you and your little children.'' Then he consoled them and spoke kindly to them.
\everypar{}

\parshape=0
\markboth{Genesis 50:22}{Genesis 50:22}\vs{22} Joseph lived in Egypt, along with his father's family. Joseph lived 110 years.\gva{g}\marginnote{\gva{g}\,Who, even though he ruled in Egypt about eighty years, yet was joined with the church of God in faith and religion.}
\markboth{Genesis 50:23}{Genesis 50:23}\vs{23} Joseph saw the descendants of Ephraim to the third generation. He also saw the children of Makir the son of Manasseh; they were given special inheritance rights by Joseph.
\everypar{}

\parshape=0
\markboth{Genesis 50:24}{Genesis 50:24}\vs{24} Then Joseph said to his brothers, ``I am about to die. But God will surely come to you and lead you up from this land to the land he swore on oath to give to Abraham, Isaac, and Jacob.''
\markboth{Genesis 50:25}{Genesis 50:25}\vs{25} Joseph made the sons of Israel swear an oath. He said, ``God will surely come to you. Then you must carry my bones up from this place.''\gva{h}\marginnote{\gva{h}\,He speaks this by the spirit of prophecy, exhorting his brethren to have full trust in God's promise for their deliverance.}
\markboth{Genesis 50:26}{Genesis 50:26}\vs{26} So Joseph died at the age of 110. After they embalmed him, his body was placed in a coffin in Egypt.

\end{scripture}
